% Compile from this directory with latexmk -lualatex main.tex
\documentclass[11pt,oneside]{article}
\usepackage[hmargin=0.98in,vmargin=1.05in,a4paper]{geometry}
\usepackage{fontspec}
\usepackage{amsmath,amsthm,mathtools}
\usepackage{unicode-math}
\setmainfont{Latin Modern Roman}
\setmathfont{Latin Modern Math}
\usepackage{microtype,booktabs,longtable,array,calc,enumitem,graphicx,fvextra}
\usepackage[svgnames]{xcolor}
\usepackage{xurl}
\usepackage[colorlinks=true,linkcolor=MidnightBlue,citecolor=MidnightBlue,urlcolor=MidnightBlue,bookmarksnumbered=true]{hyperref}
\usepackage{bookmark}
\hypersetup{pdftitle={Ten-term arithmetic progressions and single shifts in the ideal Prouhet--Tarry--Escott problem},pdfauthor={}}
\newcommand{\Q}{\mathbb Q} \newcommand{\Z}{\mathbb Z} \newcommand{\F}{\mathbb F}
\DeclareMathOperator{\Cl}{Cl} \DeclareMathOperator{\Sel}{Sel} \DeclareMathOperator{\ord}{ord}
\newfontfamily\cyrfont{DejaVu Serif}
\providecommand{\Sha}{\mathop{\text{\cyrfont Ш}}\nolimits}
\newtheorem{theorem}{Theorem}[section]
\newtheorem{proposition}[theorem]{Proposition}
\newtheorem{lemma}[theorem]{Lemma}
\newtheorem{corollary}[theorem]{Corollary}
\theoremstyle{definition}
\newtheorem{definition}[theorem]{Definition}
\newtheorem{remark}[theorem]{Remark}
\newcommand{\certbase}{https://github.com/dongxuelian2/pte-eleven-ideal-solutions/blob/934486426a928fe66bc7b228ffaadea23a4eefd1/scripts/}
\newcommand{\cert}[1]{\href{\certbase#1}{\nolinkurl{#1}}}
\newcommand{\dossier}{\href{https://github.com/dongxuelian2/pte-eleven-ideal-solutions/blob/934486426a928fe66bc7b228ffaadea23a4eefd1/PROOF.md}{arithmetic dossier}}
\setcounter{tocdepth}{1}
\setlist{nosep}
\setlength{\emergencystretch}{2em}
\allowdisplaybreaks
\providecommand{\tightlist}{\setlength{\itemsep}{0pt}\setlength{\parskip}{0pt}}
\urlstyle{same}
\begin{document}
\begin{center}
{\LARGE\bfseries Ten-term arithmetic progressions and single shifts\\[3pt] in the ideal Prouhet--Tarry--Escott problem\par}
\end{center}
\begin{abstract}
We prove that a ten-term rational arithmetic progression is determined, among rational multisets of size ten, by its first nine power sums. Equivalently, the polynomial \(\prod_{j=0}^{9}(T-j)-c\) splits over \(\Q\) only for \(c=0\). The proof singles out one squared root by descent and forces its value on a rank-one elliptic curve over a quartic field. We also classify the eighteen possible arrangements of consecutive chains in a single-shift construction from ten distinct entries per side and exclude each arrangement for a symmetric ideal target of size eleven. Finally, for any two disjoint integer sets of size \(n\) with equal power sums through degree \(n-1\), we show that equality in the ordered distance bound \(2(n-1)\) forces reflection symmetry. At size eleven the distance is at least \(22\), increasing to \(24\) for symmetric pairs. The unrestricted size-eleven existence question remains open.
\par\smallskip\noindent\textit{Keywords.} Equal sums of like powers, rational splitting, elliptic Chabauty, descent, ordered distance.
\end{abstract}
\tableofcontents
\section{Equal moments and three rigidity statements}
Write \(A=_k B\) when the sums of \(j\)th powers on the two sides agree for \(0\le j\le k\), with multiplicities. A nontrivial pair of size \(n\) with \(A=_{n-1}B\) is called \emph{ideal}. Newton's identities say that the monic root polynomials of an ideal pair differ by a nonzero constant. In particular its two sides are disjoint. We study ideal pairs with distinct entries as well as one question about multisets. Coppersmith, Mossinghoff, Scheinerman and VanderKam describe known integer solutions through size ten and at size twelve, recall the family of symmetric ten-term pairs arising from Letac's construction, and report a search for symmetric size-eleven pairs \cite[\S\S1,4.1]{cmsv}.

\begin{theorem}[A ten-term progression]\label{thm:progression}
Let \(A=\{a+jd:0\le j<10\}\), where \(a,d\in\Q\) and \(d\ne0\). If a rational multiset \(B\) of cardinality ten satisfies \(A=_9B\), then \(B=A\), including multiplicities.
\end{theorem}
The assertion depends on both the ground field and the length. The sets \(\{\pm1,\pm3\}\) and \(\{\pm9/5,\pm13/5\}\) have equal moments through degree three, since \(1+9=(9/5)^2+(13/5)^2\). Over the reals the ten simple roots of \(\prod_{j=0}^9(T-j)\) persist under small changes of its constant term. Theorem~\ref{thm:progression} concerns complete rational splitting in this pencil. Banks and van der Poorten treat a different hypothesis in which the \emph{union} of the two sides is an arithmetic progression \cite[\S3]{banks}. Here only one side is prescribed.

For the second statement, start with disjoint sets \(X,Y\) and a nonzero rational \(h\). Take multiset unions \(X\cup(Y+h)\) and \(Y\cup(X+h)\) and cancel common entries with multiplicity. The binomial theorem raises the degree of moment agreement by one. Wright's example after equation~(11) illustrates how cancellation can preserve ideality \cite{wright}. A shift by \(5\) takes \(\{0,4,5\}=_2\{1,2,6\}\) to \(\{0,4,7,11\}=_3\{1,2,9,10\}\) after cancellation of \(5,6\). Choudhry uses cancellation along arithmetic progressions to construct parametric ideal solutions in degrees four through seven \cite[\S\S2--6]{choudhry}. These constructions supply precedents for the operation and for its interaction with chains. The specific question here asks whether one shift suffices when its precursor is ideal of size ten and its output symmetric of size eleven.

Call an ideal pair \emph{symmetric} if reflection about a rational center exchanges its two sides for odd size, or preserves each side for even size.
\begin{theorem}[Single-shift obstruction]\label{thm:shift}
If the reduced output of a single rational shift of disjoint finite rational sets is a symmetric distinct-entry ideal pair of size eleven, the precursor has at least twelve entries per side. In particular no distinct-entry ideal pair of size ten produces such an output.
\end{theorem}
The precursor symmetry follows from finite-support inversion. For a size-ten precursor, normalizing the shift to one leaves eighteen partitions of consecutive chains. The progression in Theorem~\ref{thm:progression} accounts for one partition. The other seventeen involve different arithmetic curves.

For integer sets \(A=\{a_1<\cdots<a_n\}\), \(B=\{b_1<\cdots<b_n\}\), let \(L(A,B)=\sum_i|a_i-b_i|\).
\begin{theorem}[Distance]\label{thm:distance}
For every distinct-entry ideal integer pair of size \(n\ge2\), \(L(A,B)\ge2(n-1)\), with equality only for a symmetric pair. For size eleven, \(L(A,B)\ge22\) in general and \(L(A,B)\ge24\) in the symmetric case.
\end{theorem}
These last inequalities constrain a putative solution irrespective of how it is obtained. The unrestricted size-eleven existence question allows other constructions and both symmetry classes.

Sections~\ref{sec:vertices} and~\ref{sec:marked} prove Theorem~\ref{thm:progression}. Section~\ref{sec:shift} reduces Theorem~\ref{thm:shift} to the cases in the appendices. Section~\ref{sec:distance} proves Theorem~\ref{thm:distance}. The first theorem uses only the marked vertex. The size-eleven distance bounds require only the fifth partition.

\section{From a progression to five squareclasses}\label{sec:vertices}

By an affine change of coordinates the ten terms become \(\{\pm1,\pm3,\pm5,\pm7,\pm9\}\). Their root polynomial is even. Newton’s identities allow only a change of the constant coefficient, so the competing polynomial is also even. Its rational roots come in opposite pairs, with any zero occurring with even multiplicity. We may thus write them as \(\pm y_1,\ldots,\pm y_5\). The odd moments agree automatically, leaving
\begin{equation}\label{eq:four-even-moments}
 \sum_{i=1}^{5}y_i^{2j}=\sum_{a\in\{1,3,5,7,9\}}a^{2j}
 \qquad(1\leq j\leq4).
\end{equation}
No distinctness or nonvanishing assumption is imposed. Put
\[
 \mathcal A=\{1,9,25,49,81\},\qquad S_i=y_i^2,\qquad
 U=\prod_i S_i,\qquad U_0=945^2,
\]
and define
\begin{equation}\label{eq:H-polynomial}
 H(X)=X^4-165X^3+8778X^2-172810X+1057221.
\end{equation}
Newton's identities, applied through degree four, give
\begin{equation}\label{eq:split-quintic}
 \prod_{i=1}^{5}(X-S_i)=XH(X)-U,\qquad
 XH(X)-U_0=\prod_{a\in\mathcal A}(X-a).
\end{equation}
Thus the unfixed coefficient is itself a square, since
$U=(\prod_i y_i)^2$. We will show that the vertex belonging to one
specified squareclass must equal $25$; the second identity will then
recover all five vertices.

\begin{lemma}[Denominators]\label{lem:progression-denominators}
Write $y_i=n_i/D$, where $D>0$ and
$\gcd(n_1,\ldots,n_5,D)=1$. Every prime divisor of $D$ is congruent to
$1$ modulo $10$ and divides none of the $n_i$. All five numerators are
odd. In particular, the five fractions have the same reduced denominator,
and $S_i\equiv1\pmod8$ in $\Z_2$. If one $y_i$ is integral, then
$\{S_i\}_{i=1}^5=\mathcal A$ as multisets.
\end{lemma}
\begin{proof}
Set $M_j=\sum_{a\in\{1,3,5,7,9\}}a^{2j}$. Clearing denominators gives
\begin{equation}\label{eq:cleared-moments}
 \sum_i n_i^{2j}=M_jD^{2j}\qquad(1\leq j\leq4).
\end{equation}
If $2\mid D$, fourth powers modulo $16$ show that the number of odd
$n_i$ is divisible by $16$. That number lies between zero and five,
so every $n_i$ is even, contrary to primitivity. If $3\mid D$, use
sixth powers modulo $9$: a unit has sixth power one and a nonunit
has sixth power zero. The number of units is divisible by nine and
at most five, again contradicting primitivity.

The prime $5$ requires one further modulus. Suppose $5\mid D$.
Fourth powers modulo $5$ imply that the number of unit numerators
is zero or five. The first possibility contradicts primitivity,
so all five are units. Write $n_i^4=1+5a_i\pmod{25}$.
The fourth-power equation in \eqref{eq:cleared-moments} gives
$5+5\sum_i a_i\equiv0\pmod{25}$, or $\sum_i a_i\equiv-1\pmod5$.
Squaring these congruences now gives
\[
 \sum_i n_i^8\equiv5+10\sum_i a_i\equiv20\pmod{25},
\]
whereas the eighth-power equation requires zero. Thus $5\nmid D$.

Let $p>5$ divide $D$ and put $x_i=n_i^2\bmod p$. The four power
sums vanish. Since $1,2,3,4$ are invertible modulo $p$, Newton's
identities make the first four elementary symmetric functions vanish.
The root polynomial is therefore $X^5-e_5$. If $e_5=0$, every $n_i$
vanishes modulo $p$, which is impossible. Otherwise this polynomial
has five distinct nonzero roots in $\F_p$. Their ratios exhaust the
fifth roots of unity, so $5\mid p-1$. As $p$ is odd, $p\equiv1\pmod{10}$.
The roots are nonzero, proving also that $p\nmid n_i$ for every $i$.

We have proved that $D$ is odd. Now $M_2=9669\equiv5\pmod{16}$ and
$D^4\equiv1\pmod{16}$. Fourth powers count the odd numerators, forcing
all five to be odd. Their squares, divided by the odd square $D^2$,
are one modulo eight.

If some $y_i$ is integral, coprimality with $D$ gives $D=1$.
The equation $\sum_i n_i^2=165$ bounds every $|n_i|$ by
$\sqrt{165}<13$. If none belongs to $\{1,3,5,7,9\}$, all five odd
absolute values must equal $11$, contradicting the same equation.
A shared square is a common root of the two quintics in
\eqref{eq:split-quintic}; evaluating their difference gives $U=U_0$.
Their root multisets then agree.
\end{proof}

The restriction on denominator primes is sharp for these four equations
considered locally. If $p\equiv1\pmod{10}$ and $e\ge1$, put $D=p^e$.
Choose five nonzero residues $n_i$ whose squares are the fifth roots of
unity. Such choices exist because squaring is an automorphism of the
subgroup of order five. The four cleared equations hold modulo $p$.
Fix one numerator and lift the other four. The derivative matrix has
entries $2j n_i(n_i^2)^{j-1}$, for $1\le j\le4$. Its determinant is
a product of units times the Vandermonde determinant of four distinct
squares. Hensel's lemma therefore supplies a solution over $\Z_p$,
and all five resulting $y_i$ have valuation $-e$. Thus the global
argument cannot be replaced by a prohibition on negative valuations
at every prime. The dyadic condition, in contrast, places every rational
candidate in the single residue class used in Section~\ref{sec:marked}.

The denominator restriction gives affine reduction at all primes needed
below. It is important here that reduction retains multiplicity: a test
that only looked for five candidate roots would not exclude a repeated
vertex.

\begin{proposition}[Five prescribed classes]\label{prop:five-classes}
For a root $r$ of $H$, the vertices can be labeled by
$a\in\mathcal A$ so that
\begin{equation}\label{eq:five-squareclasses}
 \frac{S_a-r}{a-r}\in\Q(r)^{\times2}.
\end{equation}
Each label occurs exactly once; in particular, all vertices are nonzero
and pairwise distinct.
\end{proposition}
\begin{proof}
The certificate \cert{check_type18_product_classes.m}
gives exact assertions for the finite computations below. They use two-cover descent
for genus-two curves as implemented in Magma; see \cite{BruinStoll,magma}.
Every class-group bound used in this certificate is a Minkowski bound.

Let $d(U)=\operatorname{disc}_X(XH(X)-U)$. Explicitly,
\begin{align*}
 d(U)={}&3125U^4+108276547500U^3-1117199475459133506U^2\\
       &+533350028141359954915500U
         +2470517618419162907684321203125.
\end{align*}
The split quintic has square discriminant, the square of the
Vandermonde product. Consequently there are rational points with
parameters $U$ and $S_i$ on, respectively,
\[
 C_d:\quad V^2=Ud(U),\qquad C_H:\quad V^2=SH(S).
\]
For the second assertion use $S_iH(S_i)=U$.

First enumerate, for
\[
 p\in\{13,17,19,29,37,73,89,139,193\},
\]
all squares $u\in\F_p$ for which $XH(X)-u$ splits into square roots
with total multiplicity five. The only result is $u=U_0\bmod p$;
its roots are the five distinct elements of $\mathcal A\bmod p$.
None of these primes is $1\pmod{10}$, so Lemma~\ref{lem:progression-denominators}
applies. In particular even one such reduction rules out a zero or
repeated global vertex and shows $U\ne0$.

Here is the precise local comparison used to prune the descent classes.
For $f=d$ or $H$, the algebra of the two-cover map for $V^2=Xf(X)$
contains a rational factor at $X=0$. Represent a candidate by a polynomial
$w(X)$ and normalize it to $\widetilde w(X)=w(X)/w(0)$. The class of a
point with parameter $z$ has normalized representative $(z-X)/z$.
At a good test prime, factor $f$ into monic irreducibles $f_\nu$.
A necessary condition for compatibility with that parameter is
\begin{equation}\label{eq:resultant-filter}
 \operatorname{Res}\!\left(f_\nu,
       \widetilde w(X)\frac{z-X}{z}\right)\in\F_p^{\times2}
 \quad\text{for every }\nu.
\end{equation}
Indeed these resultants are norms of the relevant components, and
norms of squares are squares. If a representative has a pole or a
vanishing component modulo $p$, the certificate skips that comparison;
it does not discard the class. Thus this is a necessary local test,
with no assumption that norms detect all squares.

For $C_d$, descent initially gives five classes. Testing the sole
allowed parameter $U_0$ at $13$ and $19$ leaves two and then one,
namely the class represented by $U_0-X$. For $C_H$ there are initially
38 classes. Allowing any of the five seed residues, the successive
tests at $13,17,29,37,73,89,139,193$ leave
\[
 30,\ 22,\ 14,\ 13,\ 10,\ 8,\ 6,\ 5
\]
classes. The final set consists precisely of the classes $a-X$ for
$a\in\mathcal A$. At $89$, each seed residue is compatible with just
its own class. The five global vertices reduce to the five distinct
seed residues, so each global class occurs exactly once. Finally the
rational components $S_a$ and $a$ are squares. Their ratio removes
the possible common scalar in the cover comparison, giving
\eqref{eq:five-squareclasses} in the quartic component.
\end{proof}

The two auxiliary curves separate two roles of the same split quintic.
The discriminant curve records the compatibility of its five roots as
a collection; the vertex curve records an individual square root
parameter. The proof above uses the latter classes individually, with
the discriminant descent providing an additional check on the product
parameter. In particular, the passage to a labeled vertex does not
choose the smallest real root, impose an ordering on the $y_i$, or
assume that their signs agree with the seed. Labels are determined by
the local cover classes. This distinction lets the proof apply directly
to rational multisets arising from an even root polynomial, before
any chain decomposition has been made.

This assignment is stronger than membership in a union of five covers.
It forces a point on each labeled elliptic curve. The label $25$ will
suffice, and its dyadic residue is already fixed by the denominator lemma.

\section{The marked elliptic curve}\label{sec:marked}

Let $S=S_{25}$. We specify all fields and models used in the rank and
local arguments. Set
\begin{align}
 K&=\Q(h),&h^4-h^3-991h^2+13345h-19455&=0,\label{eq:quartic-K}\\
 r&=\frac{-h^3-23h^2+585h+5002}{73},&k&=25-r.\label{eq:r-and-k}
\end{align}
The element $r$ is a root of $H$, and
\[
 q(T)=\frac{H(T)}{T-r}=T^3+c_2T^2+c_1T+c_0,
\]\[
 c_2=r-165,\qquad c_1=r^2-165r+8778,\qquad
 c_0=r^3-165r^2+8778r-172810.
\]
Since $(S-r)/k$ and $H(S)=U/S$ are squares, $kq(S)$ is a square in
$K$. Therefore $S$ supplies a point
\begin{equation}\label{eq:marked-curve}
 P(S)=(kS,Y)\in E(K),\qquad
 E:\ Y^2=X^3+kc_2X^2+k^2c_1X+k^3c_0.
\end{equation}
The marked point is
\begin{equation}\label{eq:marked-point}
 P_0=(25k,189k),
\end{equation}
because $H(25)=189^2$. The following rank computation is unconditional;
its class-group input is established before the descent is used.

\subsection{A certified generating bound for the class group}

\begin{lemma}\label{lem:type18-class-group}
The cubic $q$ is irreducible over $K$. Its root field $M$ has degree
12, signature $(12,0)$, discriminant
\[
 D_M=5^6 11^9 1231^5 2417^5,
\]
and $\Cl(M)\simeq\Z/2\Z$. Prime ideals of norm at most $1000$ generate
$\Cl(M)$.
\end{lemma}
\begin{proof}
An absolute defining polynomial for $M$ is
\begin{align*}
 B(T)={}&T^{12}-3T^{11}-5625T^{10}+44481T^9+9843099T^8
 -99843252T^7\\
 &-5998855201T^6+67870752036T^5+854290494639T^4\\
 &-9750057898095T^3-2947764849192T^2
 +213431600705238T-420745260940011.
\end{align*}
The certificate \cert{type18_descent/rank.m} checks the signature and
maximal-order discriminant, irreducibility of $q$, and an isomorphism
$M\simeq K[T]/(q)$. This is also the field of an ordered pair of
distinct roots of the quartic \eqref{eq:quartic-K}.

We establish first that primes of norm at most $b=85\,000\,000$
generate. Put $X_*=2^{31}$, $\ell=\log X_*$, and
\[
 F(x)=\frac{\max(1-|x|/\ell,0)}{\cosh(x/2)},\qquad
 \Phi(s)=\int_{-\infty}^{\infty}F(x)e^{(s-1/2)x}\,dx.
\]
For a field of degree $n$ with $r_1$ real embeddings the archimedean
part of the explicit formula is
\begin{align}\label{eq:explicit-arch}
 A={}&n(\gamma+\log(8\pi))+\frac{r_1\pi}{2}
 -n\int_0^\infty\frac{1-F(x)}{2\sinh(x/2)}\,dx
 -r_1\int_0^\infty\frac{1-F(x)}{2\cosh(x/2)}\,dx.
\end{align}
We use the general-signature formula in \cite[Section~4]{miller}.
For $F_0(x)=1/\cosh(x/2)$ the integrals are $\log2$ and $\pi/2-1$.
Replacing $F_0$ by $F$ increases them by at most, respectively,
\[
 \frac1\ell\int_0^\infty\frac{x}{\sinh x}\,dx
 =\frac{\pi^2}{4\ell},\qquad
 \frac1\ell\int_0^\infty\frac{x}{2\cosh^2(x/2)}\,dx
 =\frac{2\log2}{\ell}.
\]
Here the same upper bound includes the tail $x>\ell$ since $1\le x/\ell$.
It follows that
\begin{equation}\label{eq:arch-lower}
 A\ge n\left(\gamma+\log(4\pi)-\frac{\pi^2}{4\ell}\right)
       +r_1\left(1-\frac{2\log2}{\ell}\right).
\end{equation}

If the proposed generating subgroup were proper, its quotient would
have a cyclic quotient of prime order $d$. The ordinary class field
correspondence would give an extension $N/M$ of degree $d$, unramified
at finite primes and split at real places, in which all primes of norm
at most $b$ split completely. Subtract the explicit formula for
$\zeta_M$ from that for $\zeta_N$ and divide by $d-1$.
The pole terms cancel and the discriminant term becomes $\log D_M$.
The remaining zeros belong to the nontrivial abelian Hecke
$L$-functions. Their contributions have nonnegative real part without
assuming the generalized Riemann hypothesis. To see this, for
$|a|<1/2$ the Fourier transform of $\cosh(ax)/\cosh(x/2)$ is
\[
 2\pi\Re\!\left(\operatorname{sech}(\pi(t+ia))\right)\ge0.
\]
The triangle function also has nonnegative Fourier transform, and
multiplication convolves these transforms. At $|a|=1/2$ the ratio is
one. The compact piecewise smooth test is admissible by smoothing;
equivalently its transforms decay quadratically and the zero sum
converges absolutely.

For a prime ideal $\mathfrak p$ of norm $q$ and $t=q^m\le X_*$, the
unsigned prime contribution is
\[
 w(t,m)=\frac{4\log t}{m(t+1)}\left(1-\frac{\log t}{\ell}\right).
\]
Its coefficient is the average of the nontrivial characters of the
cyclic quotient at $\operatorname{Frob}_{\mathfrak p}^m$. It is one
when $t\le b$ and is always at least $-1$. We therefore obtain
\begin{equation}\label{eq:class-bound-contradiction}
 \log D_M\ge A+S_{\mathrm{low}},
\end{equation}
where verified terms are counted positively and every other term may
be counted negatively.

The counter \cert{type13_prym_class/prime_counts.cpp}, in its
\texttt{type18} mode, counts all prime-ideal powers through $X_*$.
It decomposes small and exceptional primes in the maximal order.
For $p>\sqrt{X_*}$ outside the quartic discriminant, only degree-one
prime ideals can contribute. If Frobenius fixes $m$ quartic roots,
it fixes $m(m-1)$ ordered pairs, giving exactly that many degree-one
primes of $M$. This avoids factoring the degree-twelve polynomial at
each such prime.

Counts are binned, retaining the exponent $m$, in intervals
\[
 l=2^e(256+j)/256,\qquad
 u=\min\{X_*,2^e(257+j)/256\},\quad0\le j<256.
\]
A bin with $u\le b$ uses the lower bound
$4\log l(1-\log u/\ell)/(m(u+1))$.
Every other bin, including a bin crossing $b$, uses the negative of
$4\log u(1-\log l/\ell)/(m(l+1))$.
The rational interval computation
\cert{type13_prym_class/check_bound.py}, invoked with \texttt{--type18},
uses $n=r_1=12$ in \eqref{eq:arch-lower} and proves
\[
 A+S_{\mathrm{low}}-\log D_M>0.28.
\]
Logarithms are enclosed using the positive atanh series and a geometric
remainder; each bin contribution is rounded down. Thus the positive
margin contradicts \eqref{eq:class-bound-contradiction} and proves the
bound $b$.

The generation argument uses the ordinary class group throughout.
Passing to its cyclic quotient is what forces real places to split
in the corresponding extension and makes the signatures subtract in
exactly the same proportion as the degrees. The same subtraction
removes the pole term, which would otherwise obstruct a useful bound
with this compact test function. Every prime outside the verified
range is charged its worst allowed sign, so neither the distribution
of those primes among ideal classes nor an unproved splitting density
enters the inequality.

Finally \cert{type18_descent/check_factor_base.gp} applies PARI's exact
prime-ideal relation test through $b$ to the field defined by $B$,
and verifies that every prime in its factor base has norm at most
$1000$. This is an exact reduction of the proven generating set,
not an assumption about a provisional class group or unit index;
see \cite{pari}. Hence $1000$ is a proven generating bound.
With this bound, \cert{type18_descent/rank.m} computes class-group
invariants $[2]$. The custom bound is used only after checking both
the discriminant and a field isomorphism with $M$; other fields retain
Minkowski bounds. This is the proof-bound interface of \cite{magma}.
\end{proof}

\subsection{The rank and the locally soluble quartics}

\begin{proposition}\label{prop:marked-rank}
The curve \eqref{eq:marked-curve} has $\operatorname{rank}E(K)=1$.
\end{proposition}
\begin{proof}
Using Lemma~\ref{lem:type18-class-group}, the exact two-descent in
\cert{type18_descent/rank.m} gives
\[
 \dim_{\F_2}\Sel_2(E/K)=3.
\]
Irreducibility of $q$ shows $E(K)[2]=0$. To distinguish rank one
from rank three, we use two explicit everywhere locally soluble
quartics, rather than infer rank from a point search.

The quartics are $V^2=f_i(u)$, with
$f_i(u)=\sum_{j=0}^4 a_{ij}(h)u^j/219$. For a compact specification,
write a four-tuple $(b_3,b_2,b_1,b_0)$ for
$b_3h^3+b_2h^2+b_1h+b_0$. Their coefficient tuples are
\begingroup\small
\begin{align*}
 a_{14}&=3(326191125074,6862578283530,-228375184609562,359306674735384),\\
 a_{13}&=(-68146848515935,-1146954650973021,42423477924819832,-67064446139806428),\\
 a_{12}&=3(532419837520287,9948662909456332,-349661419343905572,551521283298925706),\\
 a_{11}&=(-17476367657300689,-311397084968974479,11197832892950970757,-17680324780047370650),\\
 a_{10}&=(70099333572212947,1276741669345440462,-45426339226399038625,71690263015899809025),\\[2pt]
 a_{24}&=3(21676496499,-813389539308,8232971883903,-11541580840841),\\
 a_{23}&=(2668272288424,-100136656976148,1013669592704180,-1421062099339992),\\
 a_{22}&=3(13684809418338,-513673043233358,5200629586198914,-7290949365165148),\\
 a_{21}&=(280668222680764,-10541234662342296,106777148298199184,-149707943110710504),\\
 a_{20}&=(719625048806446,-27039042667262238,273964615213522625,-384132484185360732).
\end{align*}
\endgroup
Here the common denominator is part of the specification; no change
of squareclass is being made by clearing it. Each displayed tuple is
interpreted in the fixed power basis of $K$ from
\eqref{eq:quartic-K}, so the quartics can be reconstructed without
choosing embeddings or matching numerical approximations to roots.

The exact certificate \cert{type18_descent/pairing.m} checks that each
quartic is squarefree and everywhere locally soluble. It also checks
the Jacobian identifications: for
$f=au^4+bu^3+cu^2+du+e$, put
\[
 I=12ae-3bd+c^2,\qquad
 J=72ace+9bcd-27ad^2-27b^2e-2c^3.
\]
The curve with coefficients $[0,0,0,-27I,-27J]$ is verified to be
$K$-isomorphic to $E$ for both quartics. The same certificate evaluates
their Cassels--Tate pairing as $1\in\F_2$, in Magma's binary convention.
Thus their images generate a subspace of dimension at least two in
$\Sha(E/K)[2]$: an alternating pairing nonzero on two vectors makes
them independent. The Kummer sequence
\[
 0\longrightarrow E(K)/2E(K)\longrightarrow\Sel_2(E/K)
 \longrightarrow\Sha(E/K)[2]\longrightarrow0
\]
therefore gives rank at most one. This does not require finiteness of
$\Sha$. The exact calculation in Lemma~\ref{lem:dyadic-direction}
shows that $P_0$ has infinite order, supplying the reverse inequality.
\end{proof}

\subsection{Excluding the entire dyadic disk}

\begin{lemma}\label{lem:dyadic-direction}
There is a unique place $v$ of $K$ above $2$, with ramification degree
one and residue degree four. In $L=K_v$, normalize $v(2)=1$.
For $T=24P_0$ on the displayed model,
\[
 v(X(T))=-6,\qquad v(Y(T))=-9,\qquad
 \overline{\frac{-X(T)}{8Y(T)}}=\bar h^3+\bar h\notin\F_2.
\]
The point $P_0$ has infinite order.
\end{lemma}
\begin{proof}
These are exact number-field computations in
\cert{type18_descent/logarithm.gp}. The certificate also verifies that
$k$ is a dyadic unit and that the given point lies on $E$.
Modulo two the defining polynomial of $h$ is
$h^4+h^3+h^2+h+1$, irreducible over $\F_2$; its root has degree four.
In particular $\bar h^3+\bar h$ is neither zero nor one.

The model is integral at $v$. Its invariant differential in the
formal parameter $z=-X/Y$ has an integral power series, so
\[
 \log_E(z)=z+\sum_{i\ge2}b_i z^i/i,\qquad b_i\in\mathcal O_L.
\]
If $v(z)\ge2$, then
$iv(z)-v_2(i)>v(z)$ for every $i\ge2$. Thus the series converges,
preserves valuation, and preserves the leading residue.
For $T$ its parameter has valuation three and nonzero logarithm,
so $T$, and hence $P_0$, has infinite order.
\end{proof}

By Proposition~\ref{prop:marked-rank}, every global elliptic logarithm
lies on the $\Q_2$-line generated by $\log_E(T)$. More explicitly,
extend the logarithm from the subgroup $v(z)\ge2$ to $E(L)$ by
multiplication by an integer taking a point into that subgroup, and
then divide by that integer. The subgroup has finite index, so this
extension is well defined and kills torsion. For any $P\in E(K)$,
rank one gives an integral relation $mP=nT+R$ with $m\ne0$ and $R$
torsion. It follows that $\log_E(P)=(n/m)\log_E(T)$.
Every nonzero element of this line has normalized leading residue
$\bar h^3+\bar h$, since a unit of $\Q_2$ reduces to one.

We now compare this direction with the chord through $P_0$. This
comparison uses a convergent logarithm on the whole admissible disk,
not a finite expansion or a search to a specified rational height.

\begin{proposition}\label{prop:marked-isolation}
A rational $S\equiv1\pmod8$ satisfying
\eqref{eq:five-squareclasses} for $a=25$ and with $H(S)$ a nonzero
rational square must equal $25$.
\end{proposition}
\begin{proof}
Suppose $S\ne25$ and set $j=v(S-25)\ge3$.
Since $q(25)=189^2/k$ is a unit and $q$ has integral dyadic
coefficients, $q(S)/q(25)\in1+8\mathcal O_L$.
The convergent binomial series supplies a square root in
$1+4\mathcal O_L$. Choose the sign of $Y$ in \eqref{eq:marked-curve}
so that $Y/(189k)\equiv1\pmod4$. Then
\[
 Q=P(S)-P_0,\qquad m=\frac{Y+189k}{k(S-25)},\qquad v(m)=1-j.
\]
The addition formulas give
\[
 X(Q)=m^2-kc_2-kS-25k,\qquad
 Y(Q)=-Y+m\bigl(kS-X(Q)\bigr).
\]
All original coordinates and coefficients are integral, whereas
$v(m)\le-2$. Consequently
\[
 X(Q)/m^2\equiv1\pmod2,\qquad
 Y(Q)/(-m^3)\equiv1\pmod2.
\]
In particular $Q$ lies in the formal subgroup with
$v(z(Q))=j-1\ge2$, and the leading residue of $z(Q)$ equals that
of $1/m$. But
\[
 \frac{1/m}{2^{j-1}}
 =\frac{(S-25)/2^j}{189}\,
   \frac{2(189k)}{Y+189k}\equiv1\pmod2.
\]
The first factor is a dyadic rational unit, and the second is one
modulo two by the sign choice. The logarithm therefore has normalized
leading residue one and is nonzero. This contradicts the global
logarithm line, whose nonzero elements have normalized residue
$\bar h^3+\bar h\notin\F_2$. Hence $S=25$.
\end{proof}

We can now finish the progression rigidity statement. The denominator
lemma and Proposition~\ref{prop:five-classes} provide exactly the vertex
required by Proposition~\ref{prop:marked-isolation}. Thus
$U=25H(25)=945^2$, and \eqref{eq:split-quintic} gives
$\{y_i^2\}=\mathcal A$ with multiplicity.

Finally, translate and scale a nonconstant ten-term arithmetic
progression to
\[\{\pm1,\pm3,\pm5,\pm7,\pm9\}.\]
If a rational
multiset has the same first nine power sums, Newton's identities make
its monic degree-ten root polynomial differ from the progression
polynomial by a constant. It is therefore even. Its rational roots
can be arranged in opposite pairs with multiplicity, including zero
pairs if present. The first four even moments of one root from each
pair are precisely \eqref{eq:four-even-moments}. The preceding argument
recovers the five seed squares, so the full ten-term multiset is the
progression. This proves Theorem~\ref{thm:progression} and supplies the
progression case needed for the shift classification.

\section{Inverting a shift and classifying its chains}\label{sec:shift}
For a finitely supported function on \(\Q\), write \(T_hu(x)=u(x-h)\) and \(R_cu(x)=u(2c-x)\). A reduced shift of the signed indicator \(u=\mathbf1_X-\mathbf1_Y\) has signed indicator \(\nu=(1-T_h)u\). Reduction keeps the signed multiplicity at every point. The classical generating-polynomial formulation records vanishing moments as divisibility by a power of \(z-1\) \cite[Lemma~1]{blp}. Inverting the shift also recovers multiplicities and support.

\begin{proposition}\label{prop:inverse}
For \(h\ne0\), a finite-support inverse exists precisely when each coset of \(h\Z\) has signed sum zero. It is unique. On the coset of \(r\) it takes the values
\[
u(r+jh)=\sum_{i\le j}\nu(r+ih).
\]
It represents two disjoint sets if and only if all these sums belong to \(\{-1,0,1\}\). If \(\nu\) has zero moments through degree \(k\ge1\), its inverse has zero moments through degree \(k-1\). If \(R_c\nu=\epsilon\nu\) for \(\epsilon=\pm1\), then \(R_{c-h/2}u=-\epsilon u\).
\end{proposition}
\begin{proof}
The cumulative formula solves the first-difference recurrence with initial value zero and has finite support exactly under the coset condition. Writing \(M_i=\sum_x x^iu(x)\), the \(j\)th moment of \(\nu\) is
\[
-\sum_{i=0}^{j-1}\binom ji h^{j-i}M_i .
\]
The successive coefficients \(-jh\) give the moment assertion by induction. Finally \(R_cT_h=T_{-h}R_c\) and \(1-T_{-h}=-T_{-h}(1-T_h)\). Thus \((1-T_h)(R_{c-h/2}u+\epsilon u)=0\), and uniqueness proves the reflection formula.
\end{proof}

\begin{corollary}
For an integer distinct-entry ideal target of size \(n+1\), a single shift of disjoint rational sets of size \(n\) has integer step and integer precursor entries.
\end{corollary}
\begin{proof}
Take \(h=p/q>0\) in lowest terms, reversing orientation if necessary. Only cosets meeting the target contribute, so the precursor lies on \(q^{-1}\Z\). Write \(P(z)=\sum_x\nu(x)z^x\). In scaled coordinates its inverse is the Laurent polynomial \(U(z)=P(z^q)/(1-z^p)\). Taking \(q\)th powers permutes the \(p\)th roots of unity. Hence \(V(z)=P(z)/(1-z^p)\) is also a Laurent polynomial and
\[
 U(z)=(1+z^p+\cdots+z^{(q-1)p})V(z^q).
\]
The summands have disjoint supports modulo \(q\). The coefficients of \(V\) lie in \(\{-1,0,1\}\), and each of its two signs occurs \(n/q\) times. As the inverse for the integer step \(p\), its two supports agree through degree \(n-1\). They are distinct, so Newton's identities give \(n/q\ge n\), proving \(q=1\). For fixed target a positive integer step divides a nonzero cross difference by the coset condition. Thus the target diameter is also a finite bound on the steps to test.
\end{proof}

\subsection{The graph of a minimal precursor}
For a symmetric target of size eleven the reflection sign is \(-1\). Proposition~\ref{prop:inverse} makes each precursor side invariant under reflection about \(c-h/2\). Both sides have the same size since \(M_0=0\). This size is even since two odd invariant sets would both contain the fixed point. Their moments through degree nine agree, so Newton's identities give size at least ten. It remains to rule out size ten. Translate the precursor center to zero and scale to \(h=1\).

Let \(P,Q\) be the two monic degree-ten root polynomials, and put \(P-Q=C\ne0\). Before cancellation the polynomials are \(F=P(x)Q(x-1)\) and \(G=Q(x)P(x-1)\). Their difference equals \(C(Q(x-1)-Q(x))\) of degree nine. Since the original colors are disjoint, their common roots are precisely the monochromatic unit edges, recorded by the larger endpoint. Every such root has multiplicity one in both products. Within one product, a double root would come from a cross-color unit edge and could never cancel. Thus a successful lift has exactly nine monochromatic edges, no cross-color edge and twenty distinct vertices. Conversely these conditions leave eleven simple roots on each side and a nonzero constant difference after cancellation.

The graph is a disjoint union of eleven paths, with consecutive vertices one unit apart. Reflection permutes the paths. A fixed path is centered at zero and has even length, because neither input contains zero. Two such paths would share \(\pm1/2\). The odd path count forces exactly one fixed path. Exchange colors to put this path, of length \(2m\), on side \(X\). Denote the multisets of lengths of the remaining reflected pairs by \(\lambda,\mu\), grouped by color. They satisfy
\begin{equation}\label{eq:parts}
 |\lambda|+|\mu|=5,\qquad \sum\lambda=5-m,\qquad \sum\mu=5,\quad 1\le m\le5 .
\end{equation}
Ordering each list gives precisely the eighteen rows in Table~\ref{tab:parts}. This enumeration runs over integer partitions only, leaving all centers rational and unrestricted.

\begin{table}[htbp]\centering\small
\caption{All solutions of \eqref{eq:parts}. A string such as \(112\) lists three path lengths \(1,1,2\). The dash denotes an empty list.}\label{tab:parts}
\begin{tabular}{rrrr@{\hspace{3em}}rrrr}
\toprule Type&\(m\)&\(\lambda\)&\(\mu\)&Type&\(m\)&\(\lambda\)&\(\mu\)\\\midrule
\hyperref[app:type1]{1}&1&4&1112&\hyperref[app:type10]{10}&2&12&113\\
\hyperref[app:type2]{2}&1&13&113&\hyperref[app:type11]{11}&2&12&122\\
\hyperref[app:type3]{3}&1&13&122&\hyperref[app:type12]{12}&2&111&14\\
\hyperref[app:type4]{4}&1&22&113&\hyperref[app:type13]{13}&2&111&23\\
\hyperref[app:type5]{5}&1&22&122&\hyperref[app:type14]{14}&3&2&1112\\
\hyperref[app:type6]{6}&1&112&14&\hyperref[app:type15]{15}&3&11&113\\
\hyperref[app:type7]{7}&1&112&23&\hyperref[app:type16]{16}&3&11&122\\
\hyperref[app:type8]{8}&1&1111&5&\hyperref[app:type17]{17}&4&1&1112\\
\hyperref[app:type9]{9}&2&3&1112&\hyperref[sec:marked]{18}&5&--&11111\\\bottomrule
\end{tabular}
\end{table}

To recover the equations for a row let \(O_\ell=\{j-(\ell-1)/2:0\le j<\ell\}\) and put
\[
 M_\ell(k,w)=\sum_{i=0}^k\binom{2k}{2i}w^{k-i}\sum_{s\in O_\ell}s^{2i},
 \qquad C_m(k)=\tfrac12\sum_{s\in O_{2m}}s^{2k}.
\]
Half the \(2k\)th moment of the two paths centered at \(\pm z\) is \(M_\ell(k,z^2)\). With centers \(u_i,v_j\) on the two sides the system becomes
\begin{equation}\label{eq:chains}
 C_m(k)+\sum_i M_{\lambda_i}(k,u_i^2)=
 \sum_j M_{\mu_j}(k,v_j^2),\qquad 1\le k\le4 .
\end{equation}
Odd moments vanish identically. In solving \eqref{eq:chains} one must test that its centers give disjoint paths, distinct entries and no cross-color unit edge. A rational solution of an auxiliary curve can fail each of these tests.

\begin{proof}[Proof of Theorem~\ref{thm:shift}]
The preceding argument exhausts the possible size-ten precursors. The individual exclusions for Types 1--9 and 11--17 are in Appendix~\ref{app:cases}, with arithmetic derivations in Appendix~\ref{app:details}. Type~10 is treated in Appendix~\ref{app:type10}. For Type~18, \(X=O_{10}\), so Theorem~\ref{thm:progression} gives \(Y=X\), contrary to disjointness. Since every such precursor has even size, at least twelve entries per side are necessary.
\end{proof}

\section{The ordered distance of an ideal pair}\label{sec:distance}
Let \(P,Q\) be the monic root polynomials of an ideal pair of size \(n\ge2\) with distinct integer entries. Since \(P-Q\) is a nonzero constant, their derivatives coincide. Rolle's theorem puts the \(n-1\) roots of \(P'=Q'\) strictly between consecutive roots of each polynomial. Each of the resulting \(n\) open intervals of monotonicity contains exactly one \(a_i\) and one \(b_i\). The order of \(a_i,b_i\) reverses in successive intervals because the sense of monotonicity reverses and the two polynomials differ by a fixed constant. This is the interlacing property of ideal pairs \cite[Theorem~2.1.3]{caley}.

The unit-step inverse \(u\) from Proposition~\ref{prop:inverse} is therefore supported on \(n\) successive blocks of alternating sign. The \(i\)th block starts at \(\min(a_i,b_i)\) and ends just before \(\max(a_i,b_i)\). Between successive blocks there is at least one unused integer. Each sign occurs at \(L/2\) places by the degree-one moment of the target. The two supports have equal moments through degree \(n-2\). Newton's identities require \(L/2\ge n-1\).

\begin{proposition}\label{prop:equality}
If \(L=2(n-1)\), the two supports of \(u\) are preserved separately by a common reflection for \(n\) odd and exchanged by a common reflection for \(n\) even.
\end{proposition}
\begin{proof}
At equality the inverse is itself an ideal pair of size \(n-1\). Its cumulative signed count stays between \(-1\) and \(1\) by the same interlacing argument (also immediate at size one). Every interior block has at most two entries and the first and last have at most one each. These upper bounds sum to \(2(n-1)\), so each is attained. Translate the first singleton to zero, choose its sign positive and let \(D\) be the last singleton. Let \(U,V\) be the monic polynomials whose roots are the starting integers of the positive and negative interior doubletons, respectively. Set \(R=V/U\) and note that \(U,V\) are coprime.

Suppose first that \(n=2m+1\). Both endpoints are positive, \(\deg U=m-1\) and \(\deg V=m\). The two root polynomials of the inverse are
\[
 x(x-D)U(x)U(x-1),\qquad V(x)V(x-1).
\]
Their constant difference gives \(R(x)R(x-1)=x(x-D)+O(x^{-2m+2})\) in Laurent series at infinity. Write \(R^*(x)=-R(D-1-x)\). Its product satisfies the same estimate on substituting \(D-x\) for \(x\). Both functions start with \(x\). A nonzero difference \(R-R^*\) with leading term \(c x^e\) has \(e\ge-2m+2\) by its denominator degree. The leading term of the difference of products would then be \(2c x^{e+1}\), impossible by the common estimate. Thus \(R=R^*\), and the zeros and poles of \(R\) are each invariant under \(s\mapsto D-1-s\). The map \(t\mapsto D-t\) preserves both colors of entries, including the singletons.

Suppose instead \(n=2m\). The last singleton is negative and \(\deg U=\deg V=m-1\). The root polynomials are \(xU(x)U(x-1)\) and \((x-D)V(x)V(x-1)\), giving
\[
 R(x)R(x-1)=\frac{x}{x-D}+O(x^{-2m+1}).
\]
Now \(R^*(x)=1/R(D-1-x)\) satisfies the same estimate by reflection and inversion of a unit series. For \(m=1\) we have \(U=V=1\) and the two singletons are exchanged directly. For larger \(m\) both functions start with one. If their difference began with \(cx^e\ne0\), again \(e\ge-2m+2\), whereas the product difference would start with \(2cx^e\), a contradiction. Thus reflection exchanges the zeros and poles of \(R\). It exchanges the two colors of entries, including the endpoints.
\end{proof}

\begin{proof}[Proof of Theorem~\ref{thm:distance}]
Taking the unit difference reverses the reflection sign in Proposition~\ref{prop:equality} and moves the center from \(D/2\) to \((D+1)/2\). Thus equality forces the asserted target symmetry. The distance is even, so a nonsymmetric pair has \(L\ge2n\).

For \(n=11\), equality \(L=20\) would give an ideal precursor of size ten. Its blocks consist of two positive singletons at the ends, four positive doubletons and five negative doubletons. The common reflection fixes one negative doubleton and pairs all the other blocks with equal lengths and colors. Exchanging colors to put the central doubleton first gives Type~5 in Table~\ref{tab:parts}. Its exclusion alone rules out \(L=20\). Hence \(L\ge22\). For every symmetric size-eleven target the inverse has even cardinality on each side by Proposition~\ref{prop:inverse}; consequently \(4\mid L\) and \(L\ge24\).
\end{proof}
If the entries all lie on a translate of \(d\Z\) with \(d>0\), translating and rescaling shows that the corresponding lower bounds are multiplied by \(d\). Wright's size-four example in the introduction attains \(2(n-1)=6\); the argument makes no sharpness assertion in size eleven.

\appendix
\section{The sixteen remaining chain models}\label{app:cases}

We record the models, exceptional parameters and completeness arguments for
Types~1--9 and 11--17. The notation $(m,\mathcal A,\mathcal B)$ means that
$A$ has the central chain of length $2m$, and its other chains occur in
reflected pairs with lengths in $\mathcal A$; the reflected chain pairs
of $B$ have lengths in $\mathcal B$. All shifts have length one. A lower-case
center parameter denotes the square of a rational center. We use the moment polynomials from \eqref{eq:chains}.
For three singleton squares whose first power sums are $p_1,p_2,p_3$,
we always use the exact monic polynomial
\[
 K(X)=X^3-p_1X^2+\frac{p_1^2-p_2}{2}X
       -\frac{p_1^3-3p_1p_2+2p_3}{6}.
\]
Thus statements that this polynomial splits into squares include the
individual center conditions. A finite-field test throughout this appendix
retains repeated roots and every possible pole chart unless a stated
valuation argument excludes that chart. An overlap, a zero reflected center,
or a repeated singleton invalidates the proposed input.

For Types~1, 9, 14 and 17, let $a=u^2$ be the noncentral $A$ center
square, $b=v^2$ the $B$ doubleton center square, and set
\[
 p_k=C_m(k)+M_{5-m}(k,a)-M_2(k,b).
\]
The necessary noncollision branches for $m=1,2,3,4$ are respectively
\[
 a=5b+4,\quad a=(10b+9)/4,\quad a=(5b+3)/3,\quad a=(5b+1)/4.
\]
Every other factor of the fourth Newton residual forces an overlap.
The factors $b=0$ give coincident reflected chains. The polynomial identities,
including every discarded factor, are checked in
\cert{check_shift_chains.py} and \cert{check_all_shift_chain_types.py}.

\subsection{Type 1}\label{app:type1}
\noindent\emph{Arithmetic details in Section~\ref{detail:type1}.}\enspace
Here $(m,\mathcal A,\mathcal B)=(1,(4),(1,1,1,2))$ and $a=5b+4$.
The three singleton squares are the roots of $K$ above. Their product gives
a point on an elliptic curve with minimal coefficients
$[1,0,0,-51865,3207225]$. Its rational group is torsion-free of rank four.
The reductions of the full group modulo doubling at $7,67,73,107,257$,
with both integral and pole tests, leave two classes represented by $b_0=0,1$.
At 73 these classes require $b\equiv b_0$. This is a restriction to global
parity classes, rather than a bound on the height of $b$.

The root-field character turns each class into a cover over
$K=\mathbf Q(w)$, where $w^3-w^2-330w-149578=0$ and
$r=(-w^2+20w-5426)/7168$. For $d=5(b_0-r)$ the exact model is
\[
 E_d:Y^2=X(X-dr)(X+4d/5),\qquad X=db.
\]
It has full two-torsion. A split two-descent and two nondegenerate
four-dimensional Cassels pairing minors give rank bounds two and one for
$b_0=0,1$, respectively. In the second case the seed is non-torsion.
For the three embeddings at 73 the primitive logarithm and parameter tangent
are
\[
 L=(30,34,22),\qquad V=(55,61,36),\qquad
 (L_iV_j-L_jV_i)_{i<j}=(33,16,28).
\]
At least two of the three signs agree. The corresponding logarithmic
linear combination has a unit derivative, so its only zero in the entire
disk is $b=1$. The doubleton then shares a vertex $1/2$ with the central
chain.

For $b_0=0$ write $b=v^2$. The center square $5b+4$ and the local character
force $v_2(v)\geq3$ if $v\ne0$. Use the sextic
\[
 h(B)=-800B^6+2625B^5-12050B^4+529B^3+54448B^2+53744B+14400.
\]
The singleton discriminant gives $h(b)$ square. In $F=\mathbf Q(\theta)$,
$h(\theta)=0$, the class of $1-b/\theta$ belongs to the norm-square
$S$-unit space with
$S=\{2,3,5,23,41,2111,11801,144461\}$. The local conditions at these
primes, at 73 and at the real places force this class to be trivial.
A second division at the two dyadic places then excludes the primitive
logarithm direction forced by every nonzero $v$. The computation uses a
certified class group of discriminant $14755078257381100$, complete local
squareclass bases and a four-cover linear system. Thus the only parameter
in this disk is $v=0$, which repeats reflected chains. The odd-prime
linear term estimates and the dyadic primitive direction estimates apply
to the whole corresponding disk.
The group, parity, pairing and local calculations are given in
\cert{check_type1_parity.m}, \cert{type1_descent/verify.gp},
\cert{type1_descent/tangents.gp}, \cert{type1_descent/disk73.m},
\cert{type1_descent/sextic.gp}, \cert{type1_descent/sextic_logs.m} and
\cert{type1_descent/sextic_four.gp}.

\subsection{Type 2}\label{app:type2}
\noindent\emph{Arithmetic details in Section~\ref{detail:type2}.}\enspace
The type is $(1,(1,3),(1,1,3))$. If $u,v$ are the length-three
centers on the two sides, the noncollision factor of the moment equations is
\[
 8u^4-40u^2v^2-58u^2+80v^4+50v^2+95=0.
\]
Put $y=8u^2-20v^2-29$. This gives
\[
 y^2=-240v^4+760v^2+81.
\]
For $v\ne0$, the substitution
$X=9(y+9)/(8v^2)+95/2$, $Y=9X/(2v)$ identifies its smooth projective
model with $E:Y^2=X^3-95X^2+2560X$. The two exceptional points
$(0,\pm9)$ give the identity and $(0,0)$. The two-isogenous curve is
$Y^2=X^3+190X^2-1215X$. For an isogeny descent class $d$, a point
requires a primitive integral solution of
\[
 w^2=du^4+Au^2v^2+(B/d)v^4.
\]
On $E$ the only possible classes are $1,2,5,10$ by positivity; 2 and 5
are excluded modulo 3. On the dual curve, among $\pm1,\pm3,\pm5,\pm15$,
all except 1 and $-15$ are excluded in the primitive charts modulo
16 or 128. Both descent images therefore have order two and the rank is
zero. Good reduction at 7 and 11 bounds torsion by ten, and the ten points
\[
 O,(0,0),(16,\pm144),(40,\pm120),(64,\pm192),(160,\pm1440)
\]
attain this bound. The complete quartic list is
$(v,y)=(0,\pm9),(\pm1/2,\pm16),(\pm3/2,\pm24)$.
At $v=0$ the value of $u^2$ is $19/4$ or $5/2$; neither is square.
At the other values the $B$ tripleton meets the central $A$ chain.
There is no rational point at infinity since $-240$ is nonsquare.
The elementary residue tables and every substitution are certified by
\cert{check_type2_exclusion.py}; \cert{check_type2_exclusion.m} gives
an additional arithmetic check.

\subsection{Type 3}\label{app:type3}
\noindent\emph{Arithmetic details in Section~\ref{app:type3-details}.}\enspace
The type is $(1,(1,3),(1,2,2))$. Let $a,b$ be the singleton and
tripleton center squares on $A$, $c$ the $B$ singleton square, and
$S,T$ the sum and product of its doubleton squares. Moment elimination gives
\[
 1024b^2-1520bc+636b+640c^2-1860c+2025=0,
 \qquad S=(20b-8c+27)/6.
\]
To justify division in the second equality, the exceptional resultant is
\[
 (4c-25)(4c-9)(4c-1)^4
 (80c^2-280c+281)(144c^2-616c+3025)(144c^2-328c+4489).
\]
The three quadratics have negative discriminant. The constant possibilities
$c=1/4,9/4,25/4$ force overlaps or nonsquare $b$.
Set $w=(b-25/4)/(c-25/4)$ and $z=(w-16)/(w-1)$.
The direction $w=1$ gives $T=0$, and $c=25/4$ was already treated.
With
\[
\begin{aligned}
 L&=z^2-57z+1656,& A&=49z^2-393z+39744,\\
 B&=9z^2+487z+14904,& C&=9z^2+247z+39744,\\
 V&=2z^2-14z+1587,
\end{aligned}
\]
one has $(a,b,c)=(A,B,C)/(4L)$ and $T=300zV/L^2$.
Positivity gives $z>0$. The individual square conditions and the
resultant support bounds reduce the squarefree parts of $L$ and the
numerator times denominator of $z$ to
$(d,k)=(1,6),(115,46)$. The projective tests use $2^{13}$ and $3^7$,
and the support bounds $d\mid690$, $k\mid138$ come from resultants,
so no other global squareclass can occur.

The first pair gives a point on
$E_6:Y^2=X(X^2-342X+59616)$ with $X=6z>0$.
This curve has rank zero and torsion of order two, so its only finite
rational abscissa is zero. For the second pair put $t=1656/z$ and
$q(t)=t^2-131t/8+3381$. Then $t$, $10L(t)$ and $15q(t)$ are squares.
The rational rank-one curve $E_{10}$ fixes the class of $t-r$, where
$r=(57+15\sqrt{-15})/2$. Over $\mathbf Q(\sqrt{-15})$ use
\[
 \ell=64(16-r)\,3/5,\qquad
 E:Y^2=X\bigl(X^2+\ell q'(r)X+\ell^2q(r)\bigr),\qquad t=r+X/\ell.
\]
The point with $X=7200\sqrt{-15}-26400$, $Y=360X$, and $(0,0)$
generate a subgroup of finite odd index. This follows from rank one,
the complete two-torsion and nondivisibility of both translates at 19.
Elliptic Chabauty has index obstruction 4 and gives precisely the rational
images $t=16,\infty$. Back-substitution gives only overlapping chains.
The exact elimination and support bounds are in
\cert{check_type3_exclusion.py}; ranks, indices and rational images are in
\cert{check_type3_exclusion.m}.

\subsection{Type 4}\label{app:type4}
\noindent\emph{Arithmetic details in Section~\ref{detail:type4}.}\enspace
The type is $(1,(2,2),(1,1,3))$. Let $c$ be the $B$ tripleton square,
$S,T$ the symmetric functions of the $A$ doubleton squares, and $U,V$
those of the $B$ singleton squares. The unique valid branch is
\[
 S=(20c+25)/6,\quad T=(240c^2+40c+167)/48,
\]\[
 U=(44c+91)/12,\quad V=(256c^2+97c+466)/36.
\]
The exceptional factors $c=1/4,9/4$, the linear branch $S=2c+1/2$
and the quadratic branch in $S$ all give explicit common vertices.
The remaining possible loss of degree has $c=-11/4$, which is not square.
The original centers require $c$ and the three polynomials
\[
 q_T=3(240c^2+40c+167),\quad q_A=-80c^2+220c+31,
 \quad q_B=3(-720c^2+2152c+275)
\]
to be squares. Descent on $Y^2=q_Tq_Aq_B$ gives three fake classes.
The three individual quadratic component norms select the class of
$c_0-\theta$, $c_0=1/4$. In an identity
$c-\theta=\lambda(c_0-\theta)z^2$, unit valuations away from 2,3,5,
the split Newton polygons at 3 and 5, and positivity at the two real
roots of $q_A$ imply $\lambda\in\mathbf Q^{\times2}$.
In particular the common scalar in a fake class is removed using the
original individual squares.

Take $s^2-s-1=0$, $r=(14+27s)/20$ and $\gamma=1/4-r$.
Every valid configuration therefore lies over
\[
 z^2=\gamma(c-r)(720c^2+120c+501).
\]
Put $L=720\gamma$, $u=Lc$, $v=Lz$ to obtain its monic elliptic model.
It has rank one over $\mathbf Q(\sqrt5)$, a non-torsion point above
$c=1/4$, and torsion of order two. Saturating at two and adjoining
torsion gives an odd-index subgroup. Elliptic Chabauty returns index
obstruction 8 and complete images $c=1/4,\infty$. The finite image
makes the paired-center squares 1 and 4 and repeats central vertices.
An original center is finite. All auxiliary proof bounds are Minkowski
bounds. The complete calculation is \cert{check_type4_exclusion.m},
with independent identities and scalar normalization in
\cert{check_type4_scalar_normalization.py}.

\subsection{Type 5}\label{app:type5}
\noindent\emph{Arithmetic details in Section~\ref{detail:type5}.}\enspace
The type is $(1,(2,2),(1,2,2))$. For the singleton square $c$ of $B$,
let $S,T$ and $U,V$ describe the two doubleton squares on $A$ and $B$.
The noncollision branch is
\[
 S=(20c-85)/16,\quad T=(80c^2-840c+589)/256,
\]\[
 U=(12c-83)/16,\quad V=(4c-49)(4c-9)/256.
\]
The exceptional values $c=1/4,9/4,25/4$ give central or cross-color
collisions. The necessary individual squares are $c$ and
\[
 q_1=80c^2-840c+589,\quad q_2=80c^2-40c+4869,
 \quad q_3=80c^2-1064c+5125.
\]
Descent on $Y^2=q_1q_2q_3$ gives the two classes represented by
$c_0=49/4,81/4$. The normalization lemma proves that
$c-\theta=\lambda(c_0-\theta)z^2$ implies
$\lambda\in\mathbf Q^{\times2}$. For odd primes outside 5 choose a
quadratic factor with unit discriminant. Square norms then give even
valuations. At 5, $c$ is integral and the root of $q_3$ of valuation
$-1$ gives a square unit ratio. At 2, the conditions force $v_2(c)=-2$;
the split roots of $q_1$ make corresponding root-difference valuations
odd for both $c$ and $c_0$, so their ratios have even valuation.
Finally $c\geq49/4$ and the two real roots of $q_1$ lie below this,
which fixes the positive sign of $\lambda$.

For $c_0=49/4$, take $K=\mathbf Q(r)$, $q_1(r)=0$, and
$g(c)=(c-r)q_3(c)/80$, $\gamma=g(49/4)$.
The cover $z^2=\gamma g(c)$ is elliptic of rank one and torsion order
two over $K=\mathbf Q(\sqrt{505})$. An odd-index subgroup obtained
by two-saturation gives complete rational images $49/4,\infty$,
with index obstruction 16. At $49/4$, $V=0$, so a reflected doubleton
pair repeats itself. For $c_0=81/4$, let $q_2(r)=0$ and
$\gamma=c_0-r$. Use the extra square $c$ to obtain
\[
 z^2=\gamma c(c-r)q_3(c)
\]
over $K=\mathbf Q(\sqrt{-95})$. Inversion in $c$ gives a rank-one
elliptic curve, again with torsion order two. The complete rational
images are $0,81/4$ and the index obstruction is 8. The value zero
has $S<0$; at $81/4$ both inputs are the ten half-integers from
$-9/2$ to $9/2$. The cover has no additional rational infinity
for the original sextic, whose leading coefficient is nonsquare.
The certificates are \cert{check_type5_scalar_normalization.py} and
\cert{check_type5_exclusion.m}. These checks retain the scalar in the
fake descent until its squarehood has been proved.

\subsection{Type 6}\label{app:type6}
\noindent\emph{Arithmetic details in Section~\ref{detail:type6}.}\enspace
The type is $(1,(1,1,2),(1,4))$. Let $a$ be the $A$ doubleton square,
$S,T$ its singleton symmetric functions, $b$ the $B$ singleton square
and $c$ its quadrupleton square. Moment elimination gives the smooth cubic
\[
 H_6=a^3-5a^2c-9a^2+10ac^2+15ac+24a-10c^3+10c^2-20c-16=0,
\]
and
\[
 b=(5a-c+5)/4,\quad S=(-3a+15c+22)/4,
 \quad T=(8a^2-40ac-67a+80c^2+95c+149)/16.
\]
The eliminated factors $a-c$ and $(a-c-1)^2-4c$ are overlaps.
The possible vanishing coefficient
$E_6=3a^2-8ac-7a+6c^2+2c+4$ has resultant proportional to
$c(c-4)^2(c-1)^3$; these values repeat reflected chains or meet the
central chain. Thus the formulas apply to every valid input.
The square cover $c=v^2$, resolved at $(a,c)=(4,0)$ by $a=4+vw$,
has plane-quartic equation
\[
 vw^3-5v^2w^2+10v^3w-10v^4+3w^2-25vw+50v^2-40=0.
\]
The origin $(1,0)$ on $H_6$ gives minimal model
$y^2+xy+y=x^3+x^2-2913x+80031$, of rank two and torsion order two.
Exact division at 2 and 3 and simultaneous reduction modulo 216 leave
exactly the classes represented by $(1,1),(9,1),(16,4)$.
The pole charts and the extra branch condition $3w^2=40$ over $(4,0)$
are retained in this reduction.

For a common Prym construction put
\[
 t=(a-1)/c,\quad A=t^3-5t^2+10t-10,\quad B=-6t^2+5t+20,
 \quad Y=2Ac+B.
\]
Then $Y^2=5t(32t^2-155t+192)$. In $K=\mathbf Q(r)$,
$r^2=-551$, let $e=(155+r)/64$ and $h=-B(e)/(2A(e))$.
The character of $q_1=a-1-ec$ is fixed in each class modulo 216,
with representatives $-e,8-e,-e$. Adjoining its square root gives
an unramified double cover of the quartic and a genus-two Prym.
On one twist two independent collision divisors span the global
logarithm space; their tangent determinants isolate the classes of
$(1,1)$ and $(16,4)$.

For the remaining twist, affine descent on divisor classes of degree one
is needed. Its global affine space has dimension 14 after homogenization
with parity, the even subspace has dimension 13, and scalars account
for 11 dimensions. Two explicit classes $D,T$ span the quotient.
The verifier constructs complete local images, including dimensions
$8,6$ at the two dyadic places, and proves the requisite divisor
representability via local Brauer obstructions. On the disk of
$(9,1)$ the displacement, normalized by $(v-1)/8$, has primitive
row $(1,0,0,0)$ outside the global span. Opposite local sheets add
$(0,0,1,0)$, also outside that span. The integrated power-series
bounds apply to the whole disk, not just its first coefficients.
The only remaining value is $(a,c)=(9,1)$, a collision.
Rank, class and exact-division data are in \cert{check_type6_classes.m}.
The two Prym twists and their complete affine descent are in
\cert{type6_descent/class_bounds.gp}, \cert{type6_descent/prym.m},
\cert{type6_descent/beta_selmer.m} and \cert{type6_descent/beta_logs.m}.

\subsection{Type 7}\label{app:type7}
\noindent\emph{Arithmetic details in Section~\ref{detail:type7}.}\enspace
The type is $(1,(1,1,2),(2,3))$. Let $a,S,T$ describe the doubleton
and singleton squares on $A$, and $b,c$ the doubleton and tripleton
squares on $B$. Every valid configuration lies on
\[
\begin{split}
 H_7={}&256a^3-1280a^2c-64a^2+2560ac^2-2560ac-2976a\\
       &-960c^3+3920c^2+780c+3159=0,
\end{split}
\]
with
\[
 b=(10a-4c+15)/6,\quad S=(16a+20c+81)/12,
\]\[
 T=-(128a^2-640ac+48a+80c^2-1380c-1701)/144.
\]
The discarded factors are $4c-1,4c-9$, the overlap quadratic
$16a^2-32ac-8a+16c^2-8c+1$, and $8a-12c-5$.
On the last branch the next equation forces $c=-11/4$ unless there
is already a collision. The subresultant coefficient can also vanish
only at the two collision values, a negative square, or roots of
quadratics with negative discriminant. This treats the exceptional loci.

With origin $(1,1/4)$ the cubic has rank two and no torsion.
Saturation at 2 and 3 followed by reductions modulo 216 leaves
\[
 (1,1/4),\ (9,9/4),\ (4,49/4),\ (4,1/4),\ (1,25/4).
\]
The square cover $b=v^2$ has genus four and the rational branch point
$(1,25/4)$. Put $u=a-1$, $t=b/u$ and
$c=(25+(10-6t)u)/4$. After removal of the branch factor its equation is
\[
 (45t^3-145t^2+135t-27)u^2-8(45t^2-80t+21)u+144(5t-2)=0,
\]
so $y^2=5F(t)$ with $F=27t^3-79t^2+57t-9$.
The omitted origin is the pole $t=\infty$; the other point at $a=1$
has $t=2/5$. The five classes give respectively
\[
 (t_0,d)=(\infty,15),(2,2),(1/3,3),(3,3),(2/5,10),\qquad t=dz^2.
\]
The support bound $d\mid30$ follows from the valuations at zeros and poles.
For a root $e$ of $F$ over its cubic field, let
$F(t)=27(t-e)(t^2+g_1t+g_0)$ and $k=135d(t_0-e)$, with the
origin character used in the first row. The elliptic curves
\[
 Y^2=X(X^2+kg_1X+k^2g_0),\qquad X=kt
\]
have ranks $0,3,3,2,2$. The first is exhausted by its two-torsion.
For the rank-two rows at 61 and 389, the determinants of two global
logarithms and the parameter tangent are units, with residues 31 and
216. A second descent in the sextic algebra treats the $d=3$ row
at $t_0=1/3$ by a rank-one cover and transversality at 83.

For the final row $d=2,t_0=2$, the same descent produces an elliptic
curve over the sextic field $L=\mathbf Q(h)$,
$h^6-96h^4+144h^2-54=0$. Its model has rational two-torsion and a
two-isogenous curve. The paired isogeny Selmer spaces have dimensions
4 and 1. Locally at the unique dyadic place, the kernel on the
four-dimensional space has one nonzero class $s_*$. A conic lift of
$s_*$ supplies the second descent character. Its two possible local
halves both lie outside the global image. The power-series estimate
on $z=1+4\xi$ bounds every higher coefficient, and division by the
maximal local power of two forces one of those excluded halves unless
$\xi=0$. The resulting point $(9,9/4)$ is a collision.
All the other anchors give repeated or cross-color vertices.
The initial group and branch checks are \cert{check_type7_classes.m};
the complete cover, logarithm and divisibility data are
\cert{type7_descent/classes.m}, \cert{type7_descent/covers.m},
\cert{type7_descent/odd_logs.gp}, \cert{type7_descent/sextic.gp} and
\cert{type7_descent/dyadic.m}.

\subsection{Type 8}\label{app:type8}
\noindent\emph{Arithmetic details in Section~\ref{detail:type8}.}\enspace
The type is $(1,(1,1,1,1),(5))$. Scale the squared $B$ center and
singleton squares by four, writing $B=4b$. Define
\[
 P_B(X)=(X-B)((X-B-4)^2-16B)((X-B-16)^2-64B),
 \qquad F(B,X)=\frac{P_B(X)-P_B(1)}{X-1}.
\]
Every valid input requires the four roots of this monic quartic to be
rational squares. Its constant term is $5f(B)$, where
$f(T)=T^4-18T^3+114T^2-97T+405$. Denominator valuations imply
that odd denominator primes are 1 modulo 20 and that $B$ is integral
at two. Indeed the denominator reduction has roots $1-\zeta^i$;
splitting forces fifth roots of unity, and square ratios force $-1$
to be square. The exceptional prime 5 is ruled out by the constant
coefficient valuation, and at two the denominator polynomial has no root.
The genus-two two-cover descent fixes one of the classes represented by
$B_0=1,9,25$.

At 23, the class $B_0=1$ forces $B\equiv1$, and
$F(1,X)=(X-1)(X-9)^2(X-25)$. Write $U=B-1,V=X-9$.
Its quadratic tangent cone is $-64(2V^2-15U^2)$. Since 30 is
nonsquare modulo 23, its valuation at a nonzero pair is exactly twice
the minimum valuation, while higher terms have larger valuation.
Thus the entire singular disk forces $B=1$, a repeated singleton.

For $B_0=9,25$, let $r$ be a root of $f$, $q(T)=f(T)/(T-r)$,
$g(T)=T^3q(1/T)/q(0)$, $t=1/B$, $t_0=1/B_0$, $k=g(t_0)$.
The reciprocal cover has monic model
\[
 y^2=x^3+kg_2x^2+k^2g_1x+k^3g_0,\qquad x=kt.
\]
Use $K=\mathbf Q(h)$, $h^4-h^3+28h^2-54h+231=0$, with
$r=(-2h^3+h^2-33h+129)/15$. The ordered-pair root field has degree
12, signature $(0,6)$, discriminant $5528109^5$ and class group of
order three. The explicit formula, with support $2^{22}$ and bound
$2^{17}$, has margin greater than $0.862$; exact ideal reduction lowers
the bound to 1000. Both reciprocal Selmer groups have dimension three.
Three explicit point logarithms establish the matching rank three.

For $B_0=9$ the primes $17,131,197,269$ fix the center residue.
Two- and three-saturation give all global coefficient classes modulo
16 and 9. Their simultaneous reductions make the three local signs
at 131 equal. The three-row global logarithm matrix has a unit
transverse tangent determinant. For $B_0=25$ the unique unramified
quartic dyadic place suffices, and its odd-square parameter lies in
one disk modulo eight. Here the common valuation of the three-row
minors is one, and adjoining the rational tangent does not raise it.
In each case the resulting annihilating form has a unit linear term;
the integrated binomial expansion makes every nonlinear coefficient
strictly smaller throughout the disk. Hence $B=9$ or 25. At both
values $P_B(1)=0$, so the two sets share a vertex square.
The certificates are \cert{check_type8_product.m},
\cert{type8_descent/node.m}, \cert{type8_descent/factor_base.gp},
\cert{type8_descent/covers.m}, \cert{type8_descent/logarithms.m} and
\cert{type8_descent/signs.m}.

\subsection{Type 9}\label{app:type9}
\noindent\emph{Arithmetic details in Section~\ref{app:type9-details}.}\enspace
Here $(m,\mathcal A,\mathcal B)=(2,(3),(1,1,1,2))$,
$a=(10b+9)/4$. In homogeneous coordinates $[R:B:W]$, the singleton
cubic is
\[
\begin{split}
 64R^3-(352B+688W)R^2+(432B^2+1840BW+2124W^2)R\\
 \qquad+216B^3-492B^2W-1974BW^2-2025W^3=0.
\end{split}
\]
The matrix
\[
 \begin{pmatrix}-44&-84&347\\-1088&3102&4064\\16&6&-28\end{pmatrix}
\]
carries it to $E:y^2+y=x^3+x^2+295x+1109$. The full rational
group is generated by $G=(31,202)$ and has no torsion.
Three roots in a common finite fiber give collinear points $n_iG$,
so $n_1+n_2+n_3=0$. Put $M=\max|n_i|$. One index has absolute
value at most $M/2$. Explicit denominator and duplication bounds
compare the three naive heights in a common fiber and give
$-46\leq h(2P)-4h(P)\leq17$. Telescoping gives bounds for the
canonical height, and the proportional index gap gives $M\leq25$.
The resulting finite group-index list is therefore an exhaustive
consequence of a global height argument. Exact substitution in the
pencil leaves only $b=4$, or $b=0$ when repeated roots are allowed.
The first is the coincident half-integer configuration; the second
has a zero reflected center. The group computation is
\cert{check_type9_group.m}; the height constants, duplication identities
and complete fiber test are \cert{check_type9_exclusion.py}.

\subsection{Type 11}\label{app:type11}
\noindent\emph{Arithmetic details in Section~\ref{app:type11-details}.}\enspace
The type is $(2,(1,2),(1,2,2))$. Let $a,b$ be the $A$ singleton
and doubleton squares, $c$ the $B$ singleton square and $S,T$ the
symmetric functions of the two $B$ doubleton squares. The necessary
noncollision conic is
\[
 16b^2+80bc-140b-80c^2+280c-281=0,
 \qquad S=(5b-4c+16)/3.
\]
The second formula can fail only at $c=1/4,9/4,49/4$.
The first two values meet the central chain. At $49/4$ the only square
choice is $b=9$, and the two residual possibilities for $S$ repeat
central $A$ vertices. Put $w=(b-9)/(c-49/4)=m/n$ in lowest terms.
The square conditions force
\[
 m^2+5mn-5n^2=-5k^2,\qquad m:n\equiv0:1\text{ or }4:1\pmod{37}.
\]
To obtain this conclusion the resultant support restricts the squarefree
part to $-1$ or $-5$; the two residues for the $-1$ class fail the
remaining product-square test at 37. The denominator is prime to 37.

Use $r=(47w-40)/(4w-5)$. The excluded value $4w=5$ would give $T=0$.
The genus-three curve needed for the last class is
\[
 C:Y^2=r(r^2-11r-1)(7r^2-2r+18)(9r^2-69r-49).
\]
At 37 only the disks of $(0,0),(8,\pm1500)$ are possible.
The Jacobian has unconditional rank at most two. The two rational
classes
\[
 P=[(8,1500)-\infty],\quad
 D=\left(r^3-\tfrac{51}{4}r^2+\tfrac{23}{4}r+8,
          \tfrac{415}{8}r^2+\tfrac{785}{8}r+20\right)
\]
have independent logarithms, since
\[
 \log(58624P)/37\equiv(25,28,33),\qquad
 \log(58624D)/37\equiv(3,30,12)\pmod{37}.
\]
Their cross product is $(12,21,0)$. The corresponding annihilating
differential has unit constant coefficient in all three disks, with
values 32 at $r=8$ and $12/31$ in the $Y$-parameter at zero.
Its integral is injective on each whole disk by the bound
$k-1-v_{37}(k)>0$ for every nonlinear degree $k\geq2$.
Thus the three rational anchors exhaust the eligible disks.
At zero $T=0$; at $r=8$ one has $c=1/4$, a central collision.
The exact rank bounds and local arithmetic are
\cert{check_type11_genus3_rank_bounds.m} and
\cert{check_type11_thirty_seven_adic.py}; the global squareclass reduction
is included in \cert{check_all_shift_chain_types.py}.

\subsection{Type 12}\label{app:type12}
\noindent\emph{Arithmetic details in Section~\ref{detail:type12}.}\enspace
The type is $(2,(1,1,1),(1,4))$. With $b$ the quadrupleton square
and $c$ the singleton square on $B$, the fourth Newton residual is
$-4b(b-4)(b-1)(b+4c-10)$. The first three factors repeat reflected
chains or meet the central chain, so $c=(10-b)/4$. The singleton
polynomial is defined exactly by
$p_k=M_4(k,b)+((10-b)/4)^k-C_2(k)$, $k=1,2,3$.
Its product and square conditions reduce the global parity classes to
$b_0=1,9$, with $(b-r)/(b_0-r)$ square in the cubic field
\[
 K=\mathbf Q(w),\quad w^3+38w-504=0,\qquad r=(4-w)/8.
\]
Put $q(B)=B^2+(r-3/2)B+r^2-3r/2+43/32$.
The product polynomial is $f(B)=160(B-r)q(B)$.
For $b_0=1$ take $d=-160(1-r)$ and
\[
 E^{(1)}:y^2=x\bigl(x^2+dq'(10)x+d^2q(10)\bigr),\qquad x=d(b-10).
\]
For $b_0=9$ let $D=160(9-r)$, $c_0=10Dq(0)$, and
\[
 F(U)=(U-c_0/10)(U^2+10Dq_1U+10Dc_0).
\]
The model is $y^2=F(x+c_0/10)$ with $x=c_0/b-c_0/10$.
Since $b$ and $10-b$ are squares, $b\ne0$ and no genuine point is
lost in inversion. Both curves have rank one over $K$. Their Selmer
groups have order 16 and rational two-torsion order two, and the
Cassels--Tate pairing on the two-dimensional quotient by the seed
and torsion is nonzero. The first descent uses a certified sextic
field bound of 100 after field identification, and the second uses
Minkowski bounds throughout.

At 907 no pole is possible, since $-1$ is nonsquare. The three
embeddings and the full singleton splitting tests leave only
$b\equiv b_0$. The logarithm/tangent minors for the three pairs
of embeddings are respectively $(698,516,760)$ and $(238,499,300)$.
All are nonzero modulo 907. Two signs among three agree, so one
minor supplies a function with a unit linear coefficient and integral
derivative. Integration bounds prove its unique zero in the whole
disk is $b=b_0$. Both values are collisions.
The certificates are \cert{check_type12_center_parity.m},
\cert{check_type12_fields.gp} and \cert{check_type12_exclusion.m}.

\subsection{Type 13}\label{app:type13}
\noindent\emph{Arithmetic details in Section~\ref{app:type13-details}.}\enspace
The type is $(2,(1,1,1),(2,3))$. The $B$ doubleton and tripleton
squares $b,c$ satisfy $6b+4c=25$ after removal of the collision
factors $(b-1)(4c-9)(4c-1)$. The three singleton squares on $A$
have polynomial $G_c$ defined by
$p_k=M_2(k,(25-4c)/6)+M_3(k,c)-C_2(k)$.
The product elliptic curve has two global Kummer classes, denoted
$\alpha$ and $\beta$, with $\alpha=25/4$ and $\beta=1/4$.
The cubic field may be written
\[
 K=\mathbf Q(\theta),\qquad
 \theta^3-\tfrac{20}{3}\theta^2+\tfrac{13}{12}\theta-\tfrac{95}{48}=0.
\]
For the $\alpha$ class the two-isogenous model is
$Y^2=X(X^2+aX+e^2)$, where
$\delta=\alpha-\theta$, $d=-6\delta$, $e=d\alpha$ and
$a=2d(\alpha-2\theta)$. Its center map is
\[
 c=\frac{\alpha(X^2+aX+e^2)}{(X-e)^2},\qquad
 \frac{\alpha}{c}=1+24\delta^2t^2,\quad t=-X/Y.
\]
A valid lift has $XY(X-e)\ne0$. Over the unique ramified cubic
dyadic completion, $\delta t=\pm v/(4w)$ where $c=w^2$ and
$4w^2+6v^2=25$. The local divisibility character is injective on
the global Selmer image. Repeated halving to the primitive point
then requires one of two local halves that are both outside that
image. This excludes the $\alpha$ class at all heights.
The $\beta$ class similarly requires $v_2(v)=1$ and
$c-\beta\in16\mathbf Z_2$; at 23 one has $c-\beta\in23\mathbf Z_{23}$.

For the last class retain signed sums of the three singleton centers.
If $t$ is the square of such a sum, their quartic is
\[
 q_2(c,t,1)^2=64t e_3(c),\quad
 q_2=t^2-\tfrac{10c+50}{3}t+5c^2-30c+19.
\]
Write it as $q_2^2=q_1q_3$, with
$q_1=t(c-\theta)$ and $q_3=64e_3(c)/(c-\theta)$.
The double cover adjoining the square root of
$q_1/(\beta-\theta)$ is unramified, and its Prym is the Jacobian
of
\[
 Y^2=-(\beta-\theta)\det(M_1+2xM_2+x^2M_3),
\]
where $M_i$ are the symmetric matrices of the homogeneous quadratics.
This exact determinant specifies the twist as well as the curve.
The quartic-root field has absolute degree twelve, signature $(2,5)$
and absolute discriminant
$172707117717588564754287819295488000000$.
The explicit formula with support $2^{31}$ proves the generating
bound $48000001$, with margin above $0.080$; exact ideal relations
reduce it to 1000. The ordinary class group is $(\mathbf Z/2)^2$.
The Prym two-descent has Selmer dimension six and two-torsion
dimension one, so its rank is at most five.

A valid split triple supplies the ordered signed sums
$-a-b+d,a-b+d,-a+b+d,a+b+d$. At $c=\beta$ these are
$1/2,3/2,7/2,9/2$. Three fixed collision classes and three moving
classes have independent logarithms. For equal signs at the three
places above 23, the determinant after removing the three displacement
factors has leading term $7(c-\beta)^2$ modulo 23. The full coefficient
bounds prevent cancellation throughout the nonzero disk. For mixed
signs the simultaneous reductions at the places above 1181 give
empty intersections for the remaining three relation tests. This
forces six independent global classes, contradicting rank at most five.
The anchor $c=\beta$ itself is a collision.
The complete data are \cert{check_type13_product_covers.py},
\cert{check_type13_product_covers.m}, \cert{check_type13_alpha_selmer.m},
\cert{check_type13_beta_tail.m}, \cert{check_type13_beta_prym_rank.m}
and \cert{check_type13_beta_prym.m}. The unconditional ideal bound is
reproduced by \cert{type13_prym_class/prime_counts.cpp},
\cert{type13_prym_class/check_bound.py} and
\cert{type13_prym_class/check_factor_base.gp}.

\subsection{Type 14}\label{app:type14}
\noindent\emph{Arithmetic details in Section~\ref{detail:type14}.}\enspace
Here $(m,\mathcal A,\mathcal B)=(3,(2),(1,1,1,2))$ and
$a=(5b+3)/3$. With $K$ defined at the beginning of the appendix,
write its discriminant as $-4b h_3(b)/243$. The necessary smooth
curve is $C:y^2=-3h_3(b)$, of genus two. Its points
$\infty,(9,\pm3240)$ occupy all three residue disks at seven.
The Jacobian has rank one, certified by two-descent with Minkowski
bounds, and $D=[(9,3240)-\infty]$ is non-torsion.
On the monic model $X=-2700b$, $Y=2700^2y$, exact Cantor
arithmetic for $25D$ gives
\[
 U\equiv(X-5)^2,\quad v_7(\operatorname{disc}U)=2,
 \quad \operatorname{disc}(U)/49\equiv5,\quad7V\equiv5X+3\pmod7.
\]
The nonsquare discriminant unit resolves the two support points in
the unramified quadratic extension. Their displacement integrals
produce the rank-one logarithm and an annihilating differential
with unit coefficient in each of the three rational disks.
The linear term dominates throughout $7\mathbf Z_7$, so each disk
contains only its anchor. The finite parameter is $b=9$, and then
\[
 a=16,\qquad K(X)=(4X-1)(4X-9)(4X-81)/64.
\]
Both inputs are the same ten half-integers. The complete rank and
point calculations are \cert{check_type14_exclusion.m} and
\cert{check_type14_seven_adic.py}.

\subsection{Type 15}\label{app:type15}
\noindent\emph{Arithmetic details in Section~\ref{detail:type15}.}\enspace
The type is $(3,(1,1),(1,1,3))$. Let $c$ be the $B$ tripleton
square and $S,T$, $U,V$ the symmetric functions of the $A$ and
$B$ singleton squares. After removal of $c=1/4,9/4,25/4$, all
of which give collisions, the exact branch is
\[
 S=(20c+15)/8,\quad T=(160c^2-20c+49)/96,
 \quad U=(69-4c)/8,\quad V=(2c-23)(4c-31)/48.
\]
In addition to $c$, the following must be rational squares:
\[
\begin{aligned}
 q_T&=6(160c^2-20c+49),&q_A&=3(-80c^2+1960c+283),\\
 q_B&=3(-80c^2+808c+2875),&q_V&=3(2c-23)(4c-31).
\end{aligned}
\]
Put $K=\mathbf Q(\sqrt{15})$, $r=(245+64\sqrt{15})/20$ and
$R=(c-r)/(1/4-r)$. Its norm is a rational square. The valuations
outside 2,3,5 are even, and it is totally positive. The fundamental
unit $4+\sqrt{15}$ and the local square test at 5 show that $R$
has class 1 or 6. The identity
$(3+\sqrt{15})^2=6(4+\sqrt{15})$ identifies the second class with
the representative $c_0=49/4$. Thus only $c_0=1/4,49/4$ occur.

Use $z^2=(c_0-r)(c-r)q(c)$, taking $q=q_B$ for $1/4$ and
$q=q_V$ for $49/4$. Both cubics are rank-one elliptic curves over
$K$, with torsion orders two and four respectively. The seed point,
after two-saturation and addition of torsion, gives finite odd index.
Elliptic Chabauty returns index obstruction 4 in both cases and the
complete rational images
\[
 \{1/4,\infty\},\qquad\{31/4,23/2,49/4,\infty\}.
\]
The two extra finite values are nonsquares. At $1/4$ an $A$
singleton meets its central chain; at $49/4$ both inputs coincide.
The normalization is elementary in $K$ and every auxiliary class
proof uses a Minkowski bound. See
\cert{check_type15_unit_normalization.py} and
\cert{check_type15_exclusion.m}.

\subsection{Type 16}\label{app:type16}
\noindent\emph{Arithmetic details in Section~\ref{detail:type16}.}\enspace
The type is $(3,(1,1),(1,2,2))$. For its $B$ singleton square $c$,
the unique noncollision branch is
\[
 S=(195-20c)/12,\quad T=(160c^2-1335c+2304)/36,
 \quad U=4(9-c)/3,\quad V=2(c-6)(4c-9)/3.
\]
The vanishing elimination denominator gives only $c=1/4,25/4$
or $21/4$. The first two collide and the third is nonsquare.
The additional branches $S=c+j/4$, $j=1,9,49$, give a singleton
square equal to $c$ and hence cross-color repetitions.
Retaining $c=v^2$ and $V=y^2$ gives
\[
 y^2=\tfrac23(v^2-6)(4v^2-9).
\]
At $t=v-3/2$ the substitution $X=-10-30/t$, $Y=-30y/t^2$
gives $E:Y^2=X^3+44X^2+100X$; the exceptional branch point
$(3/2,0)$ is the identity. Its dual is $Y^2=X^3-88X^2+1536X$.
Primitive isogeny charts modulo 9,16 and 5 give descent images
$\{1,10\}$ and $\{1,6\}$, so the rank is zero. The good
reduction orders 4 and 12 at 7 and 11 bound rational torsion by
four. The points $O,(0,0),(10,\pm80)$ exhaust it. Consequently
the quartic has exactly $(v,y)=(\pm3/2,0),(0,\pm6)$ and no
rational infinity, since $8/3$ is nonsquare. At $c=0$ the
$A$ singleton discriminant is $129/16$, a nonsquare. At $c=9/4$
an $A$ singleton square is $1/4$, already in the central chain.
Every local chart, map, group count and collision is checked by
\cert{check_type16_exclusion.py}.

\subsection{Type 17}\label{app:type17}
\noindent\emph{Arithmetic details in Section~\ref{app:type17-details}.}\enspace
Here $(m,\mathcal A,\mathcal B)=(4,(1),(1,1,1,2))$ and
$4u^2=5b+1$. Put $x=1/b>0$. Both $x$ and $x+5$ are squares,
and the singleton product makes
$11025x^3-4131x^2+472x-16$ square. Thus
\[
 E_0:Y^2=X^3-4131X^2+11025\cdot472X-16\cdot11025^2,
 \qquad X=11025x
\]
contains the product point. Its rational group is infinite cyclic,
generated by $G=(11025/16,165375/64)$.
For the cubic field $K_0=\mathbf Q(w)$,
$w^3+w^2-240w+1120=0$, a root of the product polynomial is
$r=(2w^2+5w+140)/3675$. On odd multiples $x-r$ has class
$d=1/16-r$. The curve $Y^2=X(X+5d)(X-rd)$ has rank one,
full two-torsion, and an explicitly two-saturated non-torsion point.
Elliptic Chabauty returns index obstruction 16 and images
$\infty,0,-5,1/16$. Positivity and finiteness leave $b=16$,
whose doubleton meets the central octupleton. This handles all
odd multiples.

For an even multiple write the product point as $2Q$, set $T=X(Q)$,
$a=11025$ and
\[
\begin{aligned}
 g(T)&=T^3-4131T^2+472aT-16a^2,\\
 A(T)&=T^4-944aT^2+128a^2T+a^2(472^2-4\cdot4131\cdot16),\\
 B(T)&=A(T)+20ag(T).
\end{aligned}
\]
Duplication gives $x=A/(4ag)$ and $x+5=B/(4ag)$; all three
of $g,A,B$ are nonzero squares. The two-cover descents on
$y^2=gA$ and $y^2=gB$ have sizes two and one. Compatibility
in their common cubic factor selects the identity in both, so
$T-\theta$ is square for a root $\theta$ of $A$, and
$Q\in2E_0(\mathbf Q)$. Its formal parameter $s=-T/Y(Q)$ then
satisfies $k=v_2(s)\geq3$ and $v_2(T)=-2k$.

Write $K=\mathbf Q(w)$,
$w^4-194w^2-96w+8302=0$, and
$\theta=3(2w^3+4w^2-240w-532)$. The monic elliptic curve
$Y'^2=h(T)=A(T)/(T-\theta)$ contains the putative point.
Its two-division root field has degree twelve, is totally real,
and has discriminant $58160039859394517078535559761100800000$.
The explicit formula at support $2^{20}$ has margin greater than
$0.262$. Exact ideal reduction to norm 1000 certifies class group
$\mathbf Z/2$ and Selmer dimension two.
At the unique totally ramified quartic dyadic place, put
$m=105\cdot2^{12}$. Complete local Kummer images show that
$m\log E(K_v)=2^{11}\mathfrak p$. In its basis
$(2^{12},2^{11}\pi,2^{11}\pi^2,2^{11}\pi^3)$ the global Selmer
image has dimension two and does not contain $(1,0,0,0)$.
The ratio of the two forced formal parameters satisfies
$(z/s)^2=g(T)/h(T)\in1+2^{2k}\mathcal O_{K_v}$.
The normalized global logarithm therefore has precisely that excluded
primitive vector. Dividing by the largest global power of two,
using injectivity of the localized Selmer image, proves the
contradiction at every height.
The odd calculation is \cert{check_type17_odd_branch.m}; the
unconditional field bound and the complete dyadic argument are
\cert{check_type17_factor_base.gp} and
\cert{check_type17_even_branch.m}, together with the
\texttt{--type17} option of \cert{type13_prym_class/check_bound.py}.

\section{The quadratic Chabauty calculation for Type 10}\label{app:type10}

Type 10 has a centered chain of length four on the first side, one
reflected pair of chains of length two and one reflected pair of
singletons on that side, and a reflected pair of chains of length three
and two reflected pairs of singletons on the second side. We prove that
its simultaneous rational center conditions force a collision. The
argument uses two genus-three covers, a normalized height at 13 and a
sieve on their common elliptic quotient. The sieve uses the full group
of that quotient. A full Mordell--Weil basis of either genus-three
Jacobian is unnecessary for the final step.

\subsection{Reduction to simultaneous center squares}

Let the positive centers of the first side be $u,w$, and let $v$ be the
positive center of the length-three chain pair on the second side. Put
$a=u^2$, $b=v^2$, $c=w^2$. For a chain of length $L$, write
$O_L=\{j-(L-1)/2:0\leq j<L\}$. The two remaining squared singleton
centers on the second side have power sums
\[
 p_k=\frac12\sum_{d\in O_4}d^{2k}
       +\sum_{d\in O_2}(u+d)^{2k}+c^k
       -\sum_{d\in O_3}(v+d)^{2k},\qquad 1\leq k\leq4.
\]
With $e_1=p_1$, $e_2=(p_1^2-p_2)/2$, the necessary recurrence relations
are $F=p_3-e_1p_2+e_2p_1=0$ and
$G=p_4-e_1p_3+e_2p_2=0$. Define
\begin{align*}
 D_0&=8a^2-48ab+4a+48b^2+24b-9,\\
 N_0&=48a^2b-8a^2-192ab^2-168ab+50a
                   +160b^3+352b^2+10b-27,\\
 J_0&=16a^2-32ab-8a+16b^2-8b+1,\\
 H_0&=24a^2-80ab-100a+80b^2+40b+81.
\end{align*}
Direct expansion gives
\[
 F=-\frac3{16}(N_0-2D_0c),\qquad
 \operatorname{Res}_a(N_0,2D_0)
       =-128(4b-25)(4b-9)^2(4b-1)^3.
\]
Thus $D_0=0$ forces $b=1/4,9/4,25/4$, each of which puts a
second-side chain vertex in the centered first-side chain. For a valid
configuration $D_0\ne0$ and $c=N_0/(2D_0)$. Substitution gives
\[
 G=\frac{(4b-9)(4b-1)J_0H_0}{64D_0}.
\]
The first two factors force the same collisions. Moreover
\[
 J_0=16\bigl(u^2-(v+1/2)^2\bigr)
              \bigl(u^2-(v-1/2)^2\bigr),
\]
so its vanishing makes chains of different colors intersect. Hence
$H_0=0$. A parametrization of this conic in $a,b$ is
\[
 d=3s^2-10s+10,\quad p=6s^2-5s+5,\quad
 q=3s^2+26s+10,\quad r=3s^2-10s+100,
\]
\[
 a=\frac{2p}{d},\qquad b=\frac q{4d},\qquad c=\frac r{4d}.
\]
The denominator $d$ has no rational zero. The exceptional parameters
$s=0,4,4/3$ are colliding configurations, as is the parameter at
infinity. With $z=3s-5$ define
\begin{equation}\label{eq:t10-q}
 q_0=z^2+5,\quad q_1=4z^2+30z+80,\quad
 q_2=z^2+36z+185,\quad q_3=z^2+275.
\end{equation}
The three square conditions become $a=q_1/q_0$, $4b=q_2/q_0$,
$4c=q_3/q_0$. In particular every valid finite configuration makes
all four $q_i(z)$ have the same nonzero rational squareclass.
The polynomial identities and all exceptional cases are checked by
\cert{check_sixth_mode.py} and \cert{check_sixth_center_lift.py}.

Write $z=A/B$ with coprime integers $A,B$, and homogenize $q_i$ to
$Q_i(A,B)$. All four quadratics are positive over the reals. Their
common squarefree part is therefore a positive squarefree integer
$\delta$. A prime dividing $\delta$ divides every $Q_i$, and
$Q_3-Q_0=270B^2$ implies that it lies in $\{2,3,5\}$. Such a prime
cannot divide $B$. If $A,B$ are odd, then $Q_0\equiv6\pmod8$ and
$Q_3\equiv4\pmod8$, so their squareclasses differ. The other
primitive parity cases give $Q_0\equiv1\pmod4$. This removes
$\delta=3,15$ and leaves
\begin{equation}\label{eq:t10-delta}
 \delta\in\{1,5\},\qquad q_i(z)=\delta c_i^2\quad(0\leq i\leq3).
\end{equation}
In fact $A$ is odd and $B$ is even, since the other mixed parity
would give $Q_3\equiv3$ or $7\pmod8$. This is a squareclass reduction
for the original center conditions, including $q_2$.

\subsection{The elliptic quotient and its full rational group}

Put $b_0^2=q_0q_3$. The genus-one quartic is birational to
\begin{equation}\label{eq:t10-E}
 E:\quad Y^2=X(X^2-560X+72900),\qquad
 X=2b_0+2z^2+280,\quad Y=2zX.
\end{equation}
For $X\ne0$, its inverse is $z=Y/(2X)$ and
$b_0=(X-2z^2-280)/2$. These maps extend over the smooth projective
models, so the points where the affine formulas have poles are retained.
Let $P=(180,900)$ and $T=(0,0)$. We use
\begin{equation}\label{eq:t10-full-group}
 E(\mathbf Q)=\mathbf Z P\oplus\langle T\rangle,\qquad 2T=0.
\end{equation}
Here is the arithmetic justification of completeness. The two-isogeny
descent images on $E$ and its partner
$Y^2=X(X^2+1120X+22000)$ are
\[
 \{1,5\},\qquad \{1,-5,-11,55\}.
\]
For each signed squarefree divisor $d$ of the relevant cubic coefficient,
the primitive covering quartic is tested on the two primitive local
patches modulo 256 and 3. The surviving classes are all attained by
explicit rational points. The isogeny index formula gives rank one.
Good reduction at 7 and 13 bounds the torsion order by four. There is
one nonzero rational two-torsion point, and a half of $T$ would have
$X=\pm270$, where the cubic takes a negative value. Thus the torsion
is precisely $\langle T\rangle$.

To establish that $P$ generates the free quotient, use $x/270=a/b$
and the duplication map
\[
 [a:b]\longmapsto
 [27(a^2-b^2)^2:4ab(27a^2-56ab+27b^2)].
\]
For coprime $a,b$, the coordinate gcd divides $11880$, and the maximum
of the coordinate absolute values lies between
$(121/48)\max(|a|,|b|)^4$ and $440\max(|a|,|b|)^4$.
Telescoping gives, for the doubled canonical height $h_c$ and the
naive height $h=\log\max(|a|,|b|)$,
\[
 h-\tfrac13\log(51840/11)\leq h_c
               \leq h+\tfrac13\log440.
\]
The two-isogeny class of $P$ is 5 and that of $T$ is 1, so the index
of $P$ in the free quotient is odd. If its absolute value were at
least three, a generator would have $h<\log24$. Exhaustion of
$|a|,b\leq24$ gives only the abscissae
$0,2/3,3/2,15/8,8/15$, corresponding to $T$, $\pm P$,
$\pm(P+T)$, $\pm2P$, and $\pm(2P+T)$. This contradicts a proper
odd index. The cutoff is certified by the integer inequality
$27\cdot440\cdot51840^9<24^{27}\cdot11^9$.
The rank and saturation arguments are recorded in
\cert{check_sixth_center_lift.py} and \cert{check_rank_one_generator.py}.

Write every elliptic image as $nP+\epsilon T$ with $\epsilon\in\{0,1\}$.
The common squareclass in \eqref{eq:t10-delta} is 1 for even $n$ and
5 for odd $n$. Good reductions at $13,37,59,113$, where the orders of
$P$ are $6,40,16,120$, respectively, give the necessary restrictions
\begin{equation}\label{eq:t10-mod240}
 \begin{array}{c|c}
 \epsilon& n\pmod{240}\\ \hline
 0&0,1,120\\
 1&0,120,239
 \end{array}.
\end{equation}
At undefined affine points the local test retains the point. Thus these
restrictions lose no candidate at infinity or at $X=0$. The remaining
small indices $0,\pm1$ give precisely the known collision configurations.

\subsection{Genus-three targets and rank hypotheses}

For $\delta=1,5$, let $X_\delta$ be the smooth projective curve
\[
 \delta a_0^2=q_1(z),\qquad b_0^2=q_0(z)q_3(z).
\]
Every finite simultaneous center lift maps to $X_\delta$. The quotients
by the independent signs are a conic, the elliptic quartic above and
\[
 C_\delta:\quad \delta y^2=q_0q_1q_3,\qquad y=a_0b_0.
\]
There are six simple branch points for the degree-four cover of the
$z$-line, so $X_\delta$ has genus three, and
$J(X_\delta)\sim_{\mathbf Q}E\times J(C_\delta)$.
The following conic parameters give the even-degree model
\begin{equation}\label{eq:t10-model}
 W^2=F_\delta(t)=(N_\delta^2+5D_\delta^2)
                         (N_\delta^2+275D_\delta^2),\qquad
 z=N_\delta/D_\delta,\quad a_0=M_\delta/D_\delta,\quad
 b_0=W/D_\delta^2,
\end{equation}
\[
\begin{array}{c|lll}
\delta&N_\delta&D_\delta&M_\delta\\ \hline
1&80-t^2&4t-30&2t^2-30t+160\\
5&25t^2-120t+80&10t^2-8&-60t^2+100t-48
\end{array}.
\]
They satisfy $\delta M_\delta^2=4N_\delta^2+30N_\delta D_\delta+
80D_\delta^2$. The degree-eight polynomials are squarefree and have
good ordinary reduction at 13.

A Richelot isogeny sends $J(C_\delta)$ to the Jacobian of
\[
 D_k:\quad v^2=kz(z-1)(z+5)(z^2-68z-275),\qquad k=-15/\delta.
\]
On the monic model $y^2=x(x-k)(x+5k)(x^2-68kx-275k^2)$,
full two-descent has étale algebra
$\mathbf Q^3\times\mathbf Q(\sqrt{159})$. The global norm kernel
unramified outside $\{2,3,5,11,19,53\}$ has dimension 25 over
$\mathbf F_2$. Its dimension after successive local conditions is
\[
\begin{array}{c|rrrrrrrr}
 k&\mathrm{norm}&\mathbf R&2&3&5&11&19&53\\ \hline
 -15&25&23&18&15&11&8&5&5\\
 -3 &25&23&18&15&11&8&6&5
\end{array}.
\]
The local image dimensions are $2,5,3,4,4,4,3$ at the seven listed
places. Explicit local divisor classes attain each dimension, so these
are complete local images. The required squareclass basis in
$\mathbf Q(\sqrt{159})$ follows from Minkowski's bound $<13$ and
norm relations among the primes over $2,3,5,11$. Its independence is
checked by local characters. This supplies unconditional Selmer dimension
five. Three dimensions are rational two-torsion, and the following
classes $A_k,B_k$ supply the other two dimensions. On the monic model
use
\[
\begin{array}{c|c|ll}
 k&A_k=[(x,y)-\infty]&u_{B_k}(x)&v_{B_k}(x)\\ \hline
 -15&(225/4,320625/32)&x^2+1020x+81225&-13590x-1564650\\
 -3 &(33,10692)&x^2+60x+3249&-1026x-40014
\end{array}.
\]
The divisor convention for $B_k$ is $[\operatorname{div}(u,v)-2\infty]$.
Exact divisibility $u\mid f_k-v^2$ and the independent Kummer classes
prove rank two. These calculations implement the full odd-degree descent
of Schaefer and Poonen--Schaefer~\cite{Schaefer1995,PoonenSchaefer}, and are reproduced by
\cert{check_richelot_two_selmer.py}. Thus $J(X_\delta)$ has rank three.

At 13, both genus-two reductions have 176 Jacobian points. For
$\eta=(dx/(2y),x\,dx/(2y))$, the rows of
$(\log A_k/13,\log B_k/13)$ modulo 13 are
\[
 \begin{array}{c|cc|c}
 k&\log A_k/13&\log B_k/13&\det\\ \hline
 -15&(3,1)&(1,11)&6\\
 -3 &(9,1)&(10,9)&6
 \end{array}.
\]
The exact calculation adds 176 times each divisor to an auxiliary
unramified degree-two divisor, stays in charts with unit leading
coefficients, and integrates in its two residue disks. The omitted tiny
integral terms have valuation at least $m-v_{13}(m)\geq2$ for $m\geq2$.
Since $176$ is prime to 13, the formal logarithm lattice is
$13\mathbf Z_{13}^2$. The determinant proves that the subgroup generated
by $A_k,B_k$ is 13-saturated. An independent reduction proof uses the
coefficient rows $(1,10),(0,1)$ for $k=-15$ and $(1,7),(1,3)$ for
$k=-3$ at 181 and 199. Each pair has nonzero determinant modulo 13.
The certificates are \cert{check_thirteen_logarithms.py} and
\cert{check_thirteen_saturation.py}.

For the elliptic generator, exact addition gives
$\log(P)/13\equiv1\pmod{13}$. The quotient-product and Richelot
isogenies have two-primary kernels. They therefore preserve logarithm
rank and saturation at 13. Both $J(X_\delta)(\mathbf Q)$ have logarithm
rank three. Their rational Néron--Severi rank is two. Indeed, the
Rosati-fixed fields of the genus-two reductions at 17 and 23 are
$\mathbf Q(\sqrt2)$ and $\mathbf Q(\sqrt{17})$, forcing the rational
Rosati-fixed algebra of $J(C_\delta)$ to be $\mathbf Q$. The two
nonisogenous factors of dimensions one and two give the additional
class. Consequently the quadratic Chabauty inequality is
$3<3+2-1$. We use the Mordell--Weil restricted height construction of
\cite{BalakrishnanDogra2018} and
make no assumption on the finiteness of the 13-primary Tate--Shafarevich
group.

\subsection{The correspondence and its normalized height}

Let $\iota(t,W)=(t,-W)$ and let the elliptic involution be
\[
 \sigma_1(t,W)=\left(\frac{15t-160}{2t-15},
                         \frac{9025W}{(2t-15)^4}\right),\qquad
 \sigma_5(t,W)=\left(\frac{24-25t}{25-30t},
                         \frac{9025W}{(25-30t)^4}\right).
\]
It fixes $z,b_0$ and changes the sign of $a_0$. The trace-zero
self-adjoint endomorphism is $Z=1+3\sigma$, with eigenvalues 4 and
$-2$ on the elliptic and genus-two factors. Put
\[
 B_1=(15/2,9025/16),\qquad B_5=(5/6,9025/144),\qquad
 H=\infty^++\infty^-.
\]
In both cases $\sigma(B_\delta)=\infty^+$. The four-point orbit of
$B_\delta$ under $\iota,\sigma$ is rational.

Use the proper graph correspondence
$\Gamma=\Gamma_\iota-3\Gamma_\sigma$. Its fiber degrees are $-2$,
and its diagonal intersection has degree $8-3\cdot4=-4$, giving
trace-zero degree after subtracting the two fibers. If $R_\tau$ denotes
the fixed divisor of an involution, the associated degree-zero divisor is
\begin{equation}\label{eq:t10-graphdiv}
 D_\Gamma(B,Q)=R_\iota-3R_\sigma-\iota Q+3\sigma Q
                                      -\iota B+3\sigma B.
\end{equation}
This divisor has class $Z[Q-B]+2Z[B-H/2]$. For
$j_H(Q)=[Q-H/2]$ and $\mathfrak q(A)=h_X(A,ZA)$, bilinearity gives
\[
 h_X([Q-B],D_\Gamma(B,Q))
      =\mathfrak q(j_H(Q))-\mathfrak q(j_H(B)).
\]
This is the centered height identity used below. The correspondence and finite local value theorem of
Balakrishnan--Dogra~\cite[Definition~6.2 and Theorem~6.3]{BalakrishnanDogra2018}
apply to this divisor.

Here all heights use $\log_{13}(13)=0$ and a compatible isotropic
complement to the holomorphic differentials. The elliptic complement is
the specified algebraic complement spanned by its second-kind quotient
form. The genus-two complement is transported through the Richelot
correspondence. Its pullback matrix on second-kind cohomology, in the
quotient basis, is
\[
 \begin{pmatrix}
 2&0&0&0\\0&2&0&0\\0&-32/3&-8/15&0\\
 -2200/9&-2464/9&-1114/45&-16/45
 \end{pmatrix}.
\]
Exact primitives verify the two nonholomorphic rows, and the cup matrices
satisfy the polarization identity with multiplier two. In particular the
transport fixes the height splitting, including its sign.
See \cert{check_richelot_de_rham_transport.py}.

For the two genus-two classes above, the global height matrices and the
scaled logarithm matrices are, modulo $13^3$ and $13^2$ respectively,
\[
\begin{array}{c|c|c}
 k&\mathcal H_k&L_k\\ \hline
 -15&\begin{pmatrix}923&1586\\1586&689\end{pmatrix}
     &\begin{pmatrix}68&92\\105&89\end{pmatrix}\\[3pt]
 -3&\begin{pmatrix}2093&676\\676&1027\end{pmatrix}
     &\begin{pmatrix}126&23\\105&22\end{pmatrix}
\end{array}.
\]
Thus $\mathcal A_k=L_k^{-t}(\mathcal H_k/13)L_k^{-1}$ is
\begin{equation}\label{eq:t10-heightmat}
 \mathcal A_{-15}=\begin{pmatrix}37&131\\131&135\end{pmatrix},\qquad
 \mathcal A_{-3}=\begin{pmatrix}90&114\\114&36\end{pmatrix}
                  \pmod{169}.
\end{equation}
The complete finite-prime sums for the three doubled mixed divisors
$(D_A,D_B)$, $(D_A,D_{A-B})$, $(D_B,D_{A-B})$ are
\[
\begin{array}{c|lll}
-15&-2\log(17)-2\log(593)&-2\log(5)&
                      2\log(7)+2\log(19)+2\log(313)\\
-3&-\log(3)&2\log(19)&-2\log(29)-2\log(31)
\end{array}.
\]
The dyadic terms vanish by smooth charts and signed intersection
cancellation. The finite terms depend on the chosen representatives.
For the shared-disk pairing on $D_{-3}$ one instead moves to
$4A-B$, computes its own finite intersections, and uses global
bilinearity. In that calculation the away sum is
$\log(3)-2\log(7)-2\log(61)+2\log(569)+2\log(1188727)$.
The raw Coleman integrals are corrected by their holomorphic de Rham
projections before adding these finite terms. The certificates
\cert{check_remaining_qr_heights13.py} and
\cert{check_complete_height_blocks13.py} reconstruct the displayed
matrices with explicit guard precision. This also records the factor
of two between a Jacobian class and its doubled anti-divisor.

In raw quotient logarithms $L=(L_e,L_{w,1},L_{w,2})$, write the
global quadratic expression as $L^t\mathcal J_kL/13$, with
\[
 \mathcal J_k=\operatorname{diag}(144,-\mathcal A_k/2).
\]
Adjunction and the Richelot polarization give
\[
 h_X(D,ZD)=2h_E(f_*D,f_*D)
                    -\tfrac12h_{D_k}(\phi g_*D,\phi g_*D).
\]
The elliptic coefficient satisfies
$13h_E(P,P)/\log_E(P)^2\equiv72\pmod{169}$.
The corresponding differential coordinates are
\[
 \frac{dz}{2b_0},\qquad \frac{dz}{ka_0b_0},\qquad
                              \frac{z\,dz}{a_0b_0}.
\]
Their pullback identities determine the normalization uniquely.

\subsection{The exact away-from-13 target sets}

Define the analytic function
\begin{equation}\label{eq:t10-rho}
 \rho_B(Q)=h_{13}(Q-B,D_\Gamma(B,Q))
               -\mathfrak q(j_H(Q))+\mathfrak q(j_H(B)).
\end{equation}
For rational $Q$ its value is the negative sum of the away local
heights. The complete target sets are
\begin{equation}\label{eq:t10-targets}
 \Omega_1=\{0,-1/6,-2/3\}\log_{13}(3),\qquad
 \Omega_5=\{0,1/2,2/3,-4/3\}\log_{13}(3).
\end{equation}
We give the local geometry that establishes exhaustiveness.

At every good prime other than 13, the normalized curve-level height
vanishes. The model discriminants restrict the remaining primes to
$2,3,5,11,19,53$. At 11 and 53 there is one geometric component with
ordinary nodes of thickness one. A local section cannot meet a node
with equation $xy=\ell$ over an unramified extension, since both
coordinates would have positive integral valuation. Every section
therefore retracts to the single vertex. At 19 the inner chart
$t=-2+19u$ for $\delta=1$, or $t=4+19u$ for $\delta=5$, and the
outer normalization give two elliptic components joined by two
thickness-one nodes. The involution $\sigma$ exchanges the components.
It commutes with $Z$, so the graph potential has the same value at
both vertices, and normalization makes that value zero.

At 5, pass to $\mathbf Q_5(\pi)$ with $\pi^2=5$ and set
$t=\pi u$, $W=25V$ on $X_1$. The reduction has nodes at $u=\pm1$,
each of thickness two. The main component is fixed and the two
exceptional components introduced by resolution are exchanged by
$\sigma$. A $\mathbf Q_5$ point of $X_1$ has
$v_\pi(u)=2v_5(t)-1$, so cannot reduce to either node and meets the
main component. For $X_5$, the isomorphism $a_1=\pi a_5$ gives
$u=a_5-2z/\pi$. Its conic equation forces $v_5(z)\geq1$ and
$a_5^2\equiv1\pmod5$, so all sections meet the two exceptional
components. Each twist has just one accessible orbit of components,
and its normalized local height is zero.

At 2, over any finite unramified extension, $t=2u$ on $X_1$ would give
$F_1(2u)/16\equiv3\pmod4$, which is impossible for a square. Thus
$s=1/t$ is integral for all local points, including infinity. Put
\[
 N_\infty=80s^2-1,\quad D_\infty=4s-30s^2,\quad
 S_\infty=N_\infty^2+140D_\infty^2,\quad
 V_\infty=(s^4W+S_\infty)/2.
\]
Then
\[
 V_\infty^2-S_\infty V_\infty=-18225D_\infty^4/4.
\]
Here $S_\infty\equiv1\pmod2$, so this equation splits into two
analytic disks, exchanged by $\iota$. In a semistable vertex set
containing their boundary points, further refinements within these
disks are genus-zero hanging trees. They have neither graph homology
nor positive-genus component contribution to the correspondence
measure. The potential is constant on each disk, and $\iota$
symmetry makes the constants equal. Normalization makes them zero.
The curve $X_5$ becomes $X_1$ over the unramified extension
$\mathbf Q_2(\sqrt5)$, preserving $\iota,\sigma,Z$. Vanishing thus
holds for both twists. This argument uses analytic disks and the
semistable vertex-set refinement theorem, without assuming a computed
full wild regular model.

At 3, the cluster picture has a four-root cluster near $-1$ of depth
$1/2$ containing two twins of depth $5/2$, and a four-root cluster near
4 of depth $3/2$ containing a twin of depth 2. After the quadratic
ramified base change, the charts reduce to
$V^2=(u^2-1)^2$ and $V^2=u^2(u^2-1)$. The cluster-to-graph rules give
two edges $A^+\to A^-$ of length 4, bridges $A^+\to B$ and
$B\to A^-$ of length 2, and one loop at $B$ of length 1, with lengths
normalized by $v_3(3)=1$. The involution exchanges the long edges,
fixes the bridges and reverses the loop. Cycle vectors
$(1,-1,0,0,0)$, $(1,1,-2,-2,0)$, $(0,0,0,0,1)$ have weighted
Gram matrix $\operatorname{diag}(8,24,1)$, while $Z$ acts by
$(-2,4,-2)$. The edge densities of the graph potential are
$(-1/12,-1/12,2/3,2/3,-2)$.

Normalized at $B_1$, its values along a bridge, a long edge and the
loop are respectively
\[
 -\frac1{12}+\frac s3-\frac{s^2}3,\qquad
 \frac{x^2}{24}-\frac x6-\frac1{12},\qquad
 x^2-x-\frac34.
\]
Integer valuations leave accessible only $A^\pm$, bridge coordinates
$1/2,3/2$, and the loop midpoint. The midpoint is absent for $X_1$,
since $F_1(4+9u)/3^7\equiv2u^2+2u+1\pmod3$ has no residue root
and has odd valuation. Local witnesses attain the other positions.
For $X_5$ all four positions occur. After its basepoint normalization,
the exact potential sets are
\[
 \{-1/12,0,-1/3\},\qquad \{1/4,1/3,0,-2/3\}.
\]
An intersection calculation calibrates the scalar and sign. At $t=-1$
on $X_1$, choose $W/3\equiv-1\pmod3$ and take the opposite divisor
$-D_\Gamma$. Its principal term has valuation 2, its horizontal
intersection is $3/2$, and its vertical graph correction is $1/3$.
Its height is $(2-3/2-1/3)\chi_3(3)=-\log_{13}(3)/6$.
Returning to $D_\Gamma$ gives scalar $-2\log_{13}(3)$ for the
displayed potential. This gives \eqref{eq:t10-targets} with the
negative-away convention in \eqref{eq:t10-rho}.
The graph potential is given by \cite[Theorem~3.2]{BDMTV}. The
cluster construction uses \cite[Theorems~7.1 and~8.6]{DDMM}, and the
vertex refinement uses \cite[Proposition~3.13(3)]{BPR}. All polynomial charts, accessibility witnesses, graph
Laplacians and the scalar calibration are checked by
\cert{check_qc_normalization_and_local_targets.py} and
\cert{check_complete_away_targets.py}.

\subsection{Moving normalization and its precision}

Use the monic model $Y^2=f(t)=F_\delta(t)/c^2$, with $c=1,5625$
for $\delta=1,5$, and put $\omega_j=t^jdt/(2Y)$, $\rho=\omega_3$.
If $r_j$ is the coefficient of $u^j$ in
$(u^8f(1/u))^{-1/2}$, a residue-zero basis is
\[
 \eta=(\omega_0,\omega_1,\omega_2,
             \omega_4-r_1\rho,\omega_5-r_2\rho,\omega_6-r_3\rho).
\]
The triples $(r_1,r_2,r_3)$ are $(0,-2080,33600)$ and
$(208/75,6752/1875,-386816/84375)$. The compatible complement is Lagrangian. Its lower-block determinants
are $3/760$ and $4271484375/9728$. Both are 13-adic units, so holomorphic projection
loses no precision.

For the Coleman--Gross de Rham projection $\Psi$ and holomorphic
projection $\lambda$ along this complement, the infinity input is
\[
\begin{array}{c|c|c}
\delta&\Psi\rho\pmod{2197}&\lambda\rho\pmod{2197}\\ \hline
1&(1433,392,1539,1586,728,1248)&(289,925,590)\\
5&(1194,1844,368,1690,533,1898)&(427,1168,628)
\end{array}.
\]
To compute it, expand Frobenius with
$R=f(t^{13})-f(t)^{13}$ as
\[
 13t^{13(i+1)-1}\sum_{n\geq0}\binom{-1/2}{n}
           R^n\frac{dt}{2Y^{13+26n}}.
\]
The retained terms $0\leq n\leq4$ are reduced modulo $13^{25}$ with
an auxiliary scale $13^8$. At most seven denominator divisions lose
digits. For the omitted terms, removing finite branch principal parts
loses at most $\lfloor\log_{13}(11+26n)\rfloor$ digits and horizontal
reduction loses at most one. The resulting valuation bound
\[
 n+1-\lfloor\log_{13}(11+26n)\rfloor-1\geq4\quad(n\geq5)
\]
tends to infinity. The omitted infinite tail is therefore
$O(13^4)$. The residue row of affine Frobenius is
$(0,0,0,0,0,0,13)$. Its proper block $F_{\rm dR}$ and last column $v$
satisfy $(F_{\rm dR}-13I)\Psi\rho=v$. The determinant has valuation
three, but the inverse entries lose at most one digit, so precision
$13^4$ gives $\Psi\rho$ modulo $13^3$. This is the even-degree
Frobenius projection formula of \cite[Proposition~5.3]{GajovicMuller}. The complete exact implementation is
\cert{check_infinity_projection.py}.

For a finite source $Q=(a,y_a)$ use the third-kind kernel
\[
 \nu_Q=\frac12\left(1+\frac{y_a}{Y}\right)\frac{dt}{t-a},
\]
whose residue divisor is $Q-H/2$. At a source tending to infinity set
$h=1/a$, $v=h^4y_a$ and use instead
\[
 \nu_\infty(h,v)=-\frac{h\,dt}{2(1-ht)}
                       -\frac{v t^3dt}{2Y(1-ht)}.
\]
It differs from the finite kernel by a holomorphic differential and
has the same normalized third-kind form. It specializes to
$-\rho$ at $\infty^+$.
For a moving finite source, differentiation modulo an exact form gives
\[
 \partial_a\nu_Q\equiv
 \frac{(f'(a)+f'(t))(t-a)-2(f(t)-f(a))}{(t-a)^2}
                       \frac{dt}{4y_aY}.
\]
Its numerator quotient is polynomial of degree at most six. The
analogous expression in the infinity chart is also polynomial and
residue zero. The projection commutes with this differentiation by
its global-symbol characterization: differentiate the analytic
Coleman primitives against fixed second-kind test forms and take
residues coefficientwise. Nondegeneracy of the cup product then
identifies the derivative of the projection.

A residue-correct differential for \eqref{eq:t10-graphdiv} is
\[
 d\log(W/M^3)+3\nu_{\sigma B}-\nu_{\iota B}
                         +3\nu_{\sigma Q}-\nu_{\iota Q}.
\]
Subtracting its holomorphic projection defines the local height.
For $q=t(Q)-t(B)$, integration on the basepoint disks yields
\[
 h_{13,1}=991q+2083q^2,\qquad
 h_{13,5}=1627q+81q^2\pmod{2197}.
\]
For the opposite-sheet kernel one uses its nonsingular expression
\[
 \nu_{\iota Q}=
 \frac{f(t)-f(t(Q))}{t-t(Q)}
              \frac{dt}{2Y(t)(Y(t)+Y(Q))}.
\]
All denominators are units on the relevant bidisk. The integral of a
term of total order $n\geq3$ has valuation at least
$n-v_{13}(n)\geq3$. For the twice-integrated moving normalizer the
bound is $n-v_{13}(j)-v_{13}(i+1)\geq3$, where $n=i+j+1$.
These estimates tend to infinity and control analytic tails.

The centered basepoint logarithms in the three quotient coordinates
are $(0,1313,832)$ for $\delta=1$ and $(1339,1339,1079)$ for
$\delta=5$, modulo 2197. They are determined from the signed cup
matrix and $\Psi\nu_B$, with
$\langle\Psi(\partial_a\nu_Q),\omega_j\rangle=-a^j/(2y_a)$.
If the raw logarithm increment is $d_1q+d_2q^2+\cdots$, the global
change through precision $13^3$ is
\[
 \frac{2L_B^t\mathcal J_kd_1q+
 (d_1^t\mathcal J_kd_1+2L_B^t\mathcal J_kd_2)q^2+
 2d_1^t\mathcal J_kd_2q^3}{13}.
\]
The cubic cross term is retained. Every omitted product of increment
degrees $a,b$ with $a+b\geq4$ has valuation at least
$a+b-v_{13}(a)-v_{13}(b)-1\geq3$.
Putting $q=13u$ now gives
\begin{align}\label{eq:t10-rhoseries}
 \rho_{B_1}/13&=149u+136u^2\pmod{169},\\
 \rho_{B_5}/13&=129u+110u^2+117u^3\pmod{169}.\nonumber
\end{align}
The entire analytic functions have integral coefficients on
$\mathbf Z_{13}$. All terms of degree at least three vanish modulo
13. Thus their simple residue roots exhaust the exact analytic roots,
by Hensel's lemma. The coefficients and both tail estimates are part
of the proof, while \cert{check_sixth_mode_qc_roots13.py} checks the
finite algebra and the congruences.

\subsection{The 48 local roots and the final elliptic sieve}

The mod-240 condition and squareclass parity retain four of the twelve
residue disks on each curve. On $X_1$ their reductions are
$(1,\pm1)$ and the two infinity branches with $W/t^4=\pm1$.
On $X_5$ they are $(3,\pm3)$ and the infinity branches with
$W/t^4=\pm9$. Each set is the free orbit of its basepoint disk under
$\{1,\iota,\sigma,\iota\sigma\}$.
Since $\log_{13}(3)/13\equiv5\pmod{13}$, reducing the exact target
equations from \eqref{eq:t10-targets} gives
\[
\begin{array}{c|c|c|c}
\delta&\text{target}/\log_{13}(3)&u\pmod{13}&
        \text{derivatives modulo }13\\ \hline
1&0&0,12&6,7\\
1&-1/6&2,10&4,9\\
1&-2/3&\varnothing&\\ \hline
5&0&0,11&12,1\\
5&1/2&2,9&10,3\\
5&2/3&3,8&9,4\\
5&-4/3&5,6&7,6
\end{array}.
\]
Every displayed residue gives exactly one exact local root. Thus there
are four roots on the $B_1$ disk and eight on the $B_5$ disk.

For completeness of transport, put
$Z_{\rm act}=-\iota_*+3\sigma_*$. Direct subtraction of correspondence
divisors gives
\[
 D_\Gamma(A,Q)-D_\Gamma(B,Q)=Z_{\rm act}(A-B),\qquad
 D_\Gamma(B,Q)-D_\Gamma(B,A)=Z_{\rm act}(Q-A).
\]
Local-height symmetry and self-adjointness cancel the cross terms,
so $h_A(Q)=h_B(Q)-h_B(A)$. Equivariance under $\iota,\sigma$
then shows that $h_B(\tau Q)-h_B(Q)$ is constant. Applying
$\tau$ twice makes it zero. The argument is first made with all
supports disjoint and extends by the normalized analytic kernel.
The centered global term has the same invariance. Transporting the
base-disk equations therefore gives exactly $16+32=48$ local roots
on the eight retained disks. Eight are the known basepoint-orbit
roots. The other forty are local roots to be tested, without an
assertion that they are rational points.

At each of the good primes $61,101,107$, enumerate the complete cycle
of $P$ and its translate by $T$. Their orders are $26,13,104$.
At a finite point with $X\ne0$, form $z=Y/(2X)$ and require the four
$q_i(z)$ in \eqref{eq:t10-q} to have the same nonzero quadratic
character. A zero value, the point at infinity, and $X=0$ are retained
conservatively. Projecting each permitted set to the gcd of the point
order with 3120 and combining with \eqref{eq:t10-mod240} leaves
\[
\begin{array}{c|c}
\epsilon&n\pmod{3120}\\ \hline
0&0,1,1560\\
1&0,1560,3119
\end{array}.
\]
This finite group calculation is used on the isolated height roots.
On the $B_1$ disk the elliptic reduction is the identity, the torsion
index is $\epsilon=0$, and $n$ is even. Hence the sieve forces
$n\equiv0\pmod{13}$. On the $B_5$ disk the reduction is $-P+T$,
$\epsilon=1$ and $n$ is odd, giving $n\equiv12\pmod{13}$.

The elliptic quotient differential is $dz/(2b_0)$ and its
$t$-derivative at both basepoints is 11 modulo 13. With
$\log(P)/13\equiv1\pmod{13}$, this gives
\[
 n\equiv11u\pmod{13}\quad(B_1),\qquad
 n\equiv12+11u\pmod{13}\quad(B_5).
\]
The logarithm calibration is independently checked by exact rational
addition: the formal parameter of $6P$ is $78\pmod{169}$, and its
formal logarithm is the same modulo 169. Only $u=0$ survives on
either basepoint disk. The exact zero-target equation has the root
$B_\delta$ there, and simple-root uniqueness excludes a second root
in that subdisk. The original center-square conditions are invariant
under both involutions, so the same conclusion holds on every
retained disk. All simultaneous center-lift candidates are among the
eight rational orbit points. Their elliptic indices are $0,\pm1$,
and the exact collision tests exclude them. This proves the Type 10
exclusion at every rational height.

\subsection{Arithmetic dependencies and reproducibility}

The numerical inputs in this appendix are obtained from different
certificates. Their roles are summarized here to make the proof
reproducible from the stated models.
\begin{longtable}{@{}p{.23\textwidth}p{.70\textwidth}@{}}
\toprule
Input & Certificate and mathematical role\\ \midrule
Center reduction & \cert{check_sixth_mode.py},
\cert{check_sixth_center_lift.py},
\cert{check_sixth_genus_two_covers.py}. Moment elimination,
exceptional collision parameters and coverage by $\delta=1,5$.\\
Elliptic group & \cert{check_rank_one_generator.py},
\cert{check_lift_congruences.py}. The height cutoff proves full
generation, and finite reductions give the mod-240 conditions.\\
Genus-two ranks & \cert{check_richelot_two_selmer.py}. Full local
images and an explicit global squareclass basis prove rank two
unconditionally.\\
Logarithms and index & \cert{check_thirteen_logarithms.py},
\cert{check_thirteen_saturation.py}. Unit normalized determinant and
independent finite-field tests prove full logarithm rank and
13-saturation.\\
Compatible splitting & \cert{check_genus_three_qc_de_rham.py},
\cert{check_richelot_de_rham_transport.py}. Exact cup products,
correspondence action and full second-kind pullbacks fix conventions.\\
Global heights & \cert{check_elliptic_height13.py},
\cert{check_complete_height_blocks13.py},
\cert{check_remaining_qr_heights13.py}. Normalized Coleman integrals
plus the complete representative-specific finite terms give
\eqref{eq:t10-heightmat}.\\
Curve-level local values & \cert{check_qc_normalization_and_local_targets.py},
\cert{check_complete_away_targets.py}. Regular charts, metric graph
and intersection calibration give exactly \eqref{eq:t10-targets}.\\
Infinity projection & \cert{check_infinity_projection.py}. Affine
Frobenius and the proved tail bound give $\Psi\rho$ and $\lambda\rho$.\\
Moving roots & \cert{check_sixth_mode_qc_roots13.py}. It recomputes
the derivative reductions, normalized local heights and centered
quadratic term, and checks simple-root residue lists.\\
Final candidate sieve & \cert{check_sixth_mode_candidate_sieve.py}.
Complete elliptic cycles at three primes and independent logarithm
calibration remove every non-basepoint local height root.\\ \bottomrule
\end{longtable}

The root certificate uses the global height matrices, the infinity
projection and the rank statements as inputs established by the earlier
certificates. It verifies the root series using the coefficients it
actually computes. Its Hensel refinements modulo 169 are reproducibility
outputs. The infinite-tail bounds, metric-graph interpretation and
height-cocycle argument are mathematical parts of the proof.
The final sieve requires no saturation at untested odd primes of the
genus-three Jacobians because it uses \eqref{eq:t10-full-group}.
The conclusion classifies the simultaneous-center-lift candidates
needed by Type 10. The full rational-point sets $X_1(\mathbf Q)$ and
$X_5(\mathbf Q)$ can contain points failing the remaining center-square
conditions and are outside this assertion.

\section{Descent, rational images and local isolation}\label{app:details}
Appendix~\ref{app:cases} gives the map from each chain to its arithmetic model and explains what happens at the exceptional parameters. Here we give the arithmetic within those models. For an elliptic curve over a number field, a rank-one logarithm condition can use several embeddings into \(\Q_p\) or a single extension of \(\Q_p\). In either situation only rational chain parameters are under consideration. The image of an individual curve can contain other points. For descent arguments, all class-group bounds below refer to the ordinary class group unless otherwise specified. The labels in the certificate scripts follow the chain numbering in Table~\ref{tab:parts}.
\subsection{An elementary isogeny descent for Type 16}\label{detail:type16}
We give the complete descent behind the Type~16 exclusion, including its
exceptional charts. In this type $c$ is the squared singleton center on
$B$, and $S,T$ and $U,V$ are respectively the sum and product of the
singleton squares on $A$ and the doubleton squares on $B$. The valid
moment branch is
\[
 S=\frac{195-20c}{12},\quad T=\frac{160c^2-1335c+2304}{36},
 \quad U=\frac{4(9-c)}3,\quad V=\frac{2(c-6)(4c-9)}3.
\]
Here are the divisions used to reach that branch. The first two moments give
\[
 U=\frac{4S-4c+31}{8},\quad
 V=\frac{-16S^2-32Sc+296S+64T+48c^2-296c-73}{128}.
\]
Writing $D=4S+4c-51$, the third moment determines $T$ when $D\ne0$.
The fourth residual then factors as a nonzero constant times
\[
 (4c-25)(4c-1)\prod_{j\in\{1,9,49\}}(-4S+4c+j)
 (12S+20c-195)/D^2.
\]
At $D=0$ the third residual is $-3(4c-25)(4c-21)(4c-1)/64$.
The values $c=1/4,25/4$ are central overlaps, and $21/4$ is nonsquare.
For $S=c+j/4$ the $A$ singleton polynomial is $(X-c)(X-j/4)$,
so the two colors share the singleton center. Thus every valid input
belongs to the displayed branch.

Since $c=v^2$ and $V$ is square, it gives a rational point on
\[
 \mathcal B:\quad y^2=\frac23(v^2-6)(4v^2-9).
\]
We claim that its complete projective rational-point set is
\[
 \mathcal B(\Q)=\{(3/2,0),(-3/2,0),(0,6),(0,-6)\}.
\]
There are no rational infinities, because the leading coefficient $8/3$
is nonsquare. With $t=v-3/2$ the equation becomes
$y^2=-30t+14t^2+16t^3+(8/3)t^4$. The substitutions
\[
 X=-10-30/t,\quad Y=-30y/t^2,
 \qquad v=3/2-\frac{30}{X+10},\quad y=-\frac{30Y}{(X+10)^2}
\]
give the elliptic curve
\[
 E:Y^2=X^3+44X^2+100X,
 \qquad E':Y^2=X^3-88X^2+1536X.
\]
The omitted point $(3/2,0)$ corresponds to the identity. A point of $E$
with $X=-10$ would require $Y^2=2400$, which is impossible.

For $E_{A,B}:Y^2=X(X^2+AX+B)$ the two-isogeny descent map has values
\[
 \alpha(O)=1,\qquad \alpha((0,0))=B,\qquad
 \alpha((X,Y))=X\quad\text{in }\Q^\times/\Q^{\times2}.
\]
Its kernel is the image of the dual isogeny. Taking valuations shows that
each image has a signed squarefree representative $d\mid B$. A point
with $X=du^2/v^2$ in lowest terms gives
\begin{equation}\label{eq:early-isogeny-quartic}
 w^2=du^4+Au^2v^2+(B/d)v^4,\qquad u,v,w\in\Z,
 \quad\gcd(u,v)=1.
\end{equation}
This follows by putting $Y=duw/v^3$ in the elliptic equation.
The square of a rational number is integral only if that number is
integral. The isogeny exact sequences and the Mordell--Weil theorem give
\begin{equation}\label{eq:early-isogeny-rank}
 2^{\operatorname{rank}E(\Q)}
 =\frac{|\alpha(E(\Q))|\,|\alpha'(E'(\Q))|}{4}.
\end{equation}
Both isogeny kernels have order two. This is the standard two-isogeny
rank formula~\cite[\S2]{FisherHigher}.

For $A=44,B=100$, the classes are $\pm1,\pm2,\pm5,\pm10$.
The following are the complete right-side residue sets in the primitive
charts for \eqref{eq:early-isogeny-quartic}.
\begin{center}
\begin{tabular}{ccl}
\toprule
$d$ & Modulus & Residues\\
\midrule
$-10,-1$ & $9$ & $6,8$\\
$2,5$ & $9$ & $2,5,6$\\
$-5$ & $16$ & $3,11,12$\\
$-2$ & $16$ & $8,14$\\
\bottomrule
\end{tabular}
\end{center}
For $p^e$ the charts are $(u,1)$ for all residues $u$, and $(1,v)$
with $p\mid v$. They exhaust primitive pairs after a unit rescaling.
Such a rescaling multiplies the right side by a fourth power, preserving
squarehood. The table avoids the square sets $\{0,1,4,7\}$ modulo 9 and
$\{0,1,4,9\}$ modulo 16. The remaining classes 1 and 10 occur at
$O$ and $(10,80)$, so $\alpha(E(\Q))=\{1,10\}$.
On $E'$ the classes are $\pm1,\pm2,\pm3,\pm6$. If $d<0$, all three
coefficients in \eqref{eq:early-isogeny-quartic} are negative, so there
is no nonzero real solution. For $d=2,3$ the complete primitive
residue set modulo 5 is $\{2,3\}$. The remaining classes 1 and 6
occur at the isogeny-kernel points, since $1536=6\cdot16^2$.
Formula~\eqref{eq:early-isogeny-rank} proves rank zero.

The discriminant of $E$ is $2^{17}\cdot3\cdot5^4$. Direct counts give
$\#E(\F_7)=4$ and $\#E(\F_{11})=12$. Torsion of
order prime to the reduction characteristic injects at good reduction.
Using both primes also removes the exceptional 7- and 11-primary
possibilities, so the full rational torsion order divides four.
The four points $O,(0,0),(10,\pm80)$ attain this bound, with
$2(10,80)=(0,0)$. Their inverse images are exactly the four stated
points of $\mathcal B$.

At $c=0$ one has $S=65/4,T=64$, and the singleton discriminant is
$S^2-4T=129/16$, a nonsquare. At $c=9/4$ the singleton squares are
$49/4,1/4$; the latter repeats vertices of the central chain of length
six. Hence no rational point gives a valid Type~16 input.
\cert{check_type16_exclusion.py} verifies the maps, residue tables,
reduction counts, doubling identity and final substitutions.

\subsection{Scalar normalization and the two elliptic images for Type 5}\label{detail:type5}
Here the valid moment branch is
\[
 S=\frac{20c-85}{16},\quad T=\frac{80c^2-840c+589}{256},
 \quad U=\frac{12c-83}{16},\quad V=\frac{(4c-49)(4c-9)}{256}.
\]
The factor $c=1/4$ gives a central overlap. At $c=9/4$ the $A$
paired-center polynomial has root 1 and repeats central vertices.
At $c=25/4$ it has root 4, making an $A$ vertex equal to the $B$
singleton. These are all constant factors removed when determining
$S,T,U,V$.
Retain the separate center conditions, which say that $c$ and
\[
 q_1(c)=80c^2-840c+589,\quad q_2(c)=80c^2-40c+4869,
 \quad q_3(c)=80c^2-1064c+5125
\]
are rational squares. The three quadratics retain the product and
both paired-center discriminants separately. Their product defines a
smooth genus-two curve. Put $f=q_1q_2q_3$ and
$\mathcal A=\Q[\theta]/(f(\theta))$.
The exact fake two-cover descent has two surviving classes represented
by $c_0-\theta$ for $c_0=49/4,81/4$.

We first show that a genuine input in either class satisfies
\begin{equation}\label{eq:type5-normalized-class}
 c-\theta=(c_0-\theta)z^2\quad\text{in }\mathcal A,
\end{equation}
without a rational nonsquare multiplier. Suppose initially that the
right side is $\lambda(c_0-\theta)z^2$, where
$\lambda\in\Q^\times$. Nonnegativity of $U,V$ gives
$c\geq83/12$ and $(4c-49)(4c-9)\geq0$, hence $c\geq49/4$.
The values of $(q_1,q_2,q_3)$ at the two anchors are
\[
 (48^2,128^2,64^2),\qquad(128^2,192^2,128^2).
\]
Each quadratic is irreducible, so all root differences used below are
nonzero at rational $c$ and $c_0$.

For an odd $p\ne5$, at least one of the three discriminants
\[
 2^{10}\cdot505,\qquad-2^{14}\cdot95,\qquad-2^{14}\cdot31
\]
is a $p$-unit. The corresponding quadratic, with leading coefficient
80, defines an unramified quadratic algebra at $p$. Write $r$ for
its root. If $v_p(c)<0$, the valuation is even because $c$ is square,
and every valuation of $c-r$ equals it. If $c$ is integral and the
factor is inert, $c-r$ is a unit. If it splits, its two roots have
distinct reductions, so at most one root difference has positive
valuation. Their product is $q_i(c)/80$, of even valuation, so both
valuations are even. Apply the same reasoning to $c_0$.
The squareclass identity now forces $v_p(\lambda)$ even.

At 5, a negative $v_5(c)$ would give $80c^2$ uniquely least valuation
in $q_1(c)$, namely $1+2v_5(c)$, which is odd. Thus $c$ is integral
at 5. The two Newton segments of $q_3$ give roots of valuations 3
and $-1$ in $\Q_5$. At the latter root,
\[
 \frac{c-r}{c_0-r}=\frac{1-c/r}{1-c_0/r}\in1+5\Z_5.
\]
This is a square, and the identity makes $\lambda$ a local square.

At 2 one must have $v_2(c)=-2$. A nonnegative valuation makes
$q_1(c)\equiv5\pmod8$. If the even valuation were at most $-4$,
the ratio $q_1(c)/(80c^2)$ would be 1 modulo 8: its two other terms
have relative valuations at least 3 and 4. The normalized leading
unit is again 5 modulo 8, a contradiction. Since 505 is a dyadic
square, $q_1$ splits with roots
\[
 r_\pm=\frac{105\pm4\sqrt{505}}{20},\quad
 v_2(r_\pm)=-2,\quad v_2(r_+-r_-)=1,\quad4r_\pm\equiv1\pmod8.
\]
For the square $c$, $4c\equiv1\pmod8$. Both $c-r_\pm$ have
valuation at least one, and at least one has valuation exactly one.
Their product $q_1(c)/80$ has even valuation, so both valuations
are odd. The corresponding assertion for $c_0$ makes both root
ratios have even valuation and forces $v_2(\lambda)$ even.
Finally both real roots of $q_1$ are below $49/4$, since
$\sqrt{505}<35$. The ratios are positive at both embeddings, so
$\lambda>0$. A positive rational with even valuation at every
prime is square. Absorbing its square root proves
\eqref{eq:type5-normalized-class}.

The normalization uses individual original squares. Projection to an
odd-degree factor before normalization would retain the multiplier
$\lambda$ to an odd power and would not define a fixed elliptic twist.
After normalization the following fixed twists are justified.
For $c_0=49/4$, take $K=\Q(r)$, $q_1(r)=0$, and
\[
 g(c)=\frac{(c-r)q_3(c)}{80},\qquad
 \gamma=g(49/4)=\frac{256}{5}(49/4-r).
\]
The norm of \eqref{eq:type5-normalized-class} on the cubic factor
gives $g(c)/\gamma$ square in $K$, hence $z^2=\gamma g(c)$.
If $g=c^3+g_2c^2+g_1c+g_0$, the monic model and parameter map are
\[
 E:v^2=u^3+\gamma g_2u^2+\gamma^2g_1u+\gamma^3g_0,
 \qquad (u,v)=(\gamma c,\gamma z),\quad \phi(u,v)=(u:\gamma).
\]
The point $P=(49\gamma/4,\gamma^2)$ has infinite order.
The arithmetic certificate gives rank one over $K=\Q(\sqrt{505})$
and full torsion order two. Two-saturate $P$ modulo torsion, then adjoin
the entire torsion subgroup. The resulting subgroup $G$ has finite odd
index in $E(K)$. Elliptic Chabauty with input index bound 2 returns
three elements, with image set $\{49/4,\infty\}$ and index obstruction
16. The actual odd index is coprime to 16, so these images are complete
for $E(K)$. The original parameter is finite, and the sextic also has
no rational infinity because its leading coefficient 512000 is nonsquare.
At $c=49/4$ one has $U=4,V=0$, so a reflected doubleton pair repeats.

For $c_0=81/4$, take $q_2(r)=0$, $K=\Q(\sqrt{-95})$ and
$\gamma=c_0-r$. Put $h(c)=c(c-r)q_3(c)$. The normalized cubic
norm and the additional square $c/c_0$ show $h(c)/h(c_0)$ square.
Since $q_3(c_0)=128^2$ and $c_0=(9/2)^2$, $h(c_0)$ has class
$\gamma$. Thus $z^2=\gamma h(c)$ and its seed ordinate is
$z_0=\gamma(9/2)128$. With $b=-1064,d=5125,L=-\gamma rd$,
inversion followed by scaling gives
\[
 E:v^2=u^3+\gamma(d-rb)u^2+L\gamma(b-80r)u+80L^2\gamma,
 \qquad u=L/c,\quad v=Lz/c^2,\quad\phi(u,v)=(L:u).
\]
The real bound $c\geq49/4$ ensures the inversion loses no input.
The seed $(L/c_0,Lz_0/c_0^2)$ is non-torsion, the rank is one,
and torsion has order two. The same construction of a finite odd-index
subgroup gives three Chabauty elements, obstruction 8, and complete
images $\{0,81/4\}$. At zero $S=-85/16$ is impossible. At $81/4$
the paired-center squares are $16,4$ on $A$ and $9,1$ on $B$,
with singleton centers $\pm9/2$ on $B$. Both inputs are the same ten
half-integers. This completes both classes.
All class-group proof maps are actual Minkowski bounds, including
fields introduced inside the elliptic calculations.
\cert{check_type5_exclusion.m} carries out the descents, rank and
index computations. \cert{check_type5_scalar_normalization.py}
checks the polynomial, local root and Newton identities.

\subsection{The ideal and unit calculation for Type 15}\label{detail:type15}
For Type~15 the noncollision moment formulas are
\[
 S=\frac{20c+15}{8},\quad T=\frac{160c^2-20c+49}{96},
 \quad U=\frac{69-4c}{8},\quad V=\frac{(2c-23)(4c-31)}{48}.
\]
The removed values $c=1/4,9/4,25/4$ meet opposite-color chains.
The rational quantities $c$ and
\[
\begin{aligned}
 q_T&=6(160c^2-20c+49),&q_A&=3(-80c^2+1960c+283),\\
 q_B&=3(-80c^2+808c+2875),&q_V&=3(2c-23)(4c-31)
\end{aligned}
\]
are squares. The latter are respectively $576T$, $576(S^2-4T)$,
$576(U^2-4V)$ and $144V$.
Put
\[
 K=\Q(\sqrt{15}),\quad r=\frac{245+64\sqrt{15}}{20},
 \quad c_0=1/4,\quad R=\frac{c-r}{c_0-r}.
\]
The polynomial $q_A$ is irreducible. Its positivity places $c$ between
its real roots, as is $c_0$ because $q_A(c_0)=48^2$.
Thus $R$ is totally positive, and
$N_{K/\Q}(R)=q_A(c)/q_A(c_0)$ is a rational square.

Outside 2,3,5 the polynomial $q_A$ has unit leading coefficient and
discriminant. At a pole of $c$ its even valuation gives an even
valuation of both root differences. For integral $c$ an inert root
difference is a unit; in the split case at most one root difference
has positive valuation, which must be even because the rational norm
is square. The same holds for $c_0$. At 2,3,5 the field has one
prime of residue degree one, so the square rational norm directly
forces its valuation of $R$ to be even. Consequently $(R)=I^2$.

The ring of integers is $\Z[\sqrt{15}]$ and its discriminant
is 60. Its Minkowski bound is $\sqrt{15}<4$, so the class group
is generated by the unique primes $\mathfrak p_2,\mathfrak p_3$
of norms two and three. Since
$N(3+\sqrt{15})=-6$, one has $(3+\sqrt{15})=\mathfrak p_2\mathfrak p_3$,
and $\mathfrak p_2^2=(2)$. The prime $\mathfrak p_2$ is
nonprincipal: norm $\pm2$ would require a rational integer square
to be 2 or 3 modulo 5. The class group therefore has order two.
It follows that
\[
 R=2^d\epsilon^e z^2,\qquad d\in\{0,1\},\quad
 \epsilon=4+\sqrt{15},\quad z\in K^\times.
\]
Here the totally positive units are exactly the powers of $\epsilon$.
There is no norm $-1$ unit by reduction modulo 3. If a norm-one unit
$v=x+y\sqrt{15}$ satisfied $1<v<\epsilon$, then
$y=(v-v^{-1})/(2\sqrt{15})$ would be an integer strictly between zero
and one. Thus $\epsilon$ is fundamental.

At 5, a negative even valuation of $c$ would make the leading term
of $q_T(c)$ have the unique least valuation $1+2v_5(c)$, an odd
number. Hence $c$ is integral. The Newton polygon of $q_A$ has one
segment of length two and its roots have valuation $-1/2$ at 5.
Therefore $R\in1+\mathfrak p_5$, which is square in the completion.
The residue of $\epsilon$ is 4 in $\F_5$, whereas 2 is
nonsquare, forcing $d=0$. Finally
\[
 (3+\sqrt{15})^2=6\epsilon,\qquad
 \frac{49/4-r}{6(1/4-r)}=
 \left(\frac{2}{3+\sqrt{15}}\right)^2.
\]
It follows that either $(c-r)/(1/4-r)$ or $(c-r)/(49/4-r)$ is square
in $K$. This is a complete squareclass conclusion, with no height
bound on $c$.

For the two classes use respectively $(c_0,q,s_0)=(1/4,q_B,96)$ and
$(49/4,q_V,9)$, so $q(c_0)=s_0^2$. If $\gamma=c_0-r$, the original
squares give $z^2=\gamma(c-r)q(c)$. Write $q(c)=ac^2+bc+d$ and
$L=a\gamma$. Then
\[
 E:v^2=u^3+\gamma(b-ar)u^2+L\gamma(d-br)u-L^2\gamma dr,
 \qquad (u,v)=(Lc,Lz),\qquad\phi(u,v)=(u:L).
\]
The seed $(Lc_0,L\gamma s_0)$ has infinite order. Each curve has
rank one over $K$. Two-saturation of the seed and addition of the full
torsion subgroup, of orders two and four, gives a subgroup of finite
odd index. Elliptic Chabauty with input index bound 2 returns
\begin{center}
\begin{tabular}{cccc}
\toprule
$c_0$ & Number of group elements & Index obstruction & Rational images\\
\midrule
$1/4$ & $3$ & $4$ & $1/4,\infty$\\
$49/4$ & $5$ & $4$ & $31/4,23/2,49/4,\infty$\\
\bottomrule
\end{tabular}
\end{center}
The odd subgroup indices are coprime to the returned obstructions, so the
lists are complete. Every class-group proof bound in these computations
is a Minkowski bound. The images $31/4,23/2$ are nonsquares, and a
center parameter is finite. At $1/4$ an $A$ singleton is $1/2$,
already on its central chain. At $49/4$ the $A$ singleton absolute
centers are $9/2,7/2$, the $B$ singleton absolute centers are $3/2,1/2$,
and the $B$ tripleton centers are $\pm7/2$. Both inputs are the ten
half-integers from $-9/2$ to $9/2$. This excludes the two classes.
The complete arithmetic and exact unit identities are
\cert{check_type15_exclusion.m} and
\cert{check_type15_unit_normalization.py}.

\subsection{Component norms and the scalar for Type 4}\label{detail:type4}
For Type~4 the valid moment branch has
\[
 S=\frac{20c+25}{6},\quad T=\frac{240c^2+40c+167}{48},
 \quad U=\frac{44c+91}{12},\quad V=\frac{256c^2+97c+466}{36}.
\]
To check the exceptional factors directly, the first two moments give
\[
 U=2S-3c-3/4,\quad
 V=(8S^2-48Sc-24S+16T+48c^2+66c+9)/8.
\]
The resultant of the remaining two equations in $T$ is proportional to
\[
\begin{split}
 &(4c-9)(4c-1)(-6S+20c+25)(-2S+4c+1)\\
 &\qquad\cdot(4S^2-16Sc-20S+16c^2+24c+25).
\end{split}
\]
The first two factors are central overlaps. On $S=2c+1/2$, the third
moment gives $T=(4c-1)^2/16$. With $c=v^2$ the paired centers are
$\pm(v-1/2),\pm(v+1/2)$, whose doubletons overlap the $B$ tripleton.
On the quadratic branch,
\[
 S=2v^2+5/2\pm2v,\qquad
 T=(2v\mp1)^2(2v\pm3)^2/16,
\]
and the same overlap occurs. The coefficients used to determine $T$
can vanish only at $c=-11/4,1/4,9/4$; the negative value is not a
rational square and the other two were excluded. Thus the stated
branch loses no valid input.
The separate square conditions are $c$ and
\[
 q_T=3(240c^2+40c+167),\qquad
 q_A=-80c^2+220c+31,\qquad
 q_B=3(-720c^2+2152c+275).
\]
These are $144T$, $9(S^2-4T)$ and $144(U^2-4V)$.
Let $f=q_Tq_Aq_B$, $\mathcal A=\Q[\theta]/(f(\theta))$,
and $c_0=1/4$. The values of the three quadratics at $c_0$ are
$24^2,9^2,48^2$.

The fake two-cover descent of $Y^2=f(c)$ returns three classes.
For a representative $\delta$, let $\delta_i$ be its projection
to the quadratic factor field $K_i$ of $q_i$. Necessarily
\[
 \operatorname{lc}(q_i)N_{K_i/\Q}(\delta_i)
       \in\Q^{\times2}\qquad(i=T,A,B).
\]
The condition is independent of the fake representative, since a
rational scalar has square norm on a quadratic component.
For an actual point it is exactly the original square condition
$q_i(c)\in\Q^{\times2}$. Exactly one computed class passes
all three tests and it is represented by $c_0-\theta$.
Thus a genuine input has
$c-\theta=\lambda(c_0-\theta)z^2$ in $\mathcal A$.

The three discriminants are
\[
 -2^{10}3^2\cdot5\cdot31,\quad 2^4 3^6\cdot5,
 \quad 2^{14}3^2\cdot331.
\]
For each $p\notin\{2,3,5\}$ choose a component with unit leading
coefficient and discriminant. The same root-difference argument used
in Type~5 gives even valuations of both corresponding ratios. It
forces $v_p(\lambda)$ even.
If $v_3(c)<0$, it is even and at most $-2$. In
$q_T(c)=720c^2+120c+501$ the leading term is uniquely minimal,
of valuation $2+2v_3(c)$ and normalized unit $80\equiv2\pmod3$.
The other terms have relative valuations at least one and three,
so the polynomial is not a local square. Hence $c$ is integral at 3.
The polynomial $-720c^2+2152c+275$ has roots of valuations $-2,0$
in $\Q_3$. At the first root the ratio of differences lies
in $1+9\Z_3$ and is square. Thus $\lambda$ is a
$\Q_3$-square.

At 5 a negative even $v_5(c)$ would make $720c^2$ the unique term
of lowest, odd valuation in $q_T(c)$. Hence $c$ is integral there.
The same quadratic has $\Q_5$ roots of valuations $-1,2$.
The ratio at the first root lies in $1+5\Z_5$, so
$\lambda$ is a $\Q_5$-square. Since $q_A(c)>0$ and
$q_A(c_0)>0$, both parameters lie strictly between its two real
roots. Their root ratios are positive at both embeddings, forcing
$\lambda>0$. The known valuations leave only rational classes 1
and 2. The second is nonsquare in $\Q_3$. Thus $\lambda$
is a rational square and can be absorbed into $z$.

Let $s^2-s-1=0$, $r=(14+27s)/20$ and $\gamma=1/4-r$.
The normalized identity projected to $K=\Q(s)$ makes
$(c-r)/\gamma$ square, so
\[
 z^2=\gamma(c-r)(720c^2+120c+501).
\]
With $a=720,b=120,d=501,L=a\gamma$, its monic model is
\[
 E_0:v^2=u^3+\gamma(b-ar)u^2+L\gamma(d-br)u-L^2\gamma dr,
 \qquad(u,v)=(Lc,Lz),\quad\phi_0(u,v)=(u:L).
\]
The point $P_0=(L/4,24L\gamma)$ is non-torsion. The computation
transports both this point and the rational function through an exact
minimal-model isomorphism $\iota:E_0\to E$, so the function on $E$
is $\phi_0\circ\iota^{-1}$. It has the original $c$-coordinate as
its image. The rank over $K$ is one and torsion has order two.
Two-saturation and torsion give finite odd index; elliptic Chabauty
returns three elements, index obstruction 8 and images
$\{1/4,\infty\}$. The coprime-index condition proves completeness.
At $1/4$ the two paired-center squares are 1 and 4, whose chains
repeat $\pm1/2,\pm3/2$, including the central chain. Infinity is
not a center parameter. All auxiliary class-group proof bounds are
Minkowski bounds. The certificates are
\cert{check_type4_exclusion.m} and
\cert{check_type4_scalar_normalization.py}.

\subsection{The complete elliptic quotient for Type 2}\label{detail:type2}
The necessary Type~2 relation for the rational tripleton centers $u,v$ is
\[
 8u^4-40u^2v^2-58u^2+80v^4+50v^2+95=0.
\]
Consequently $y=8u^2-20v^2-29$ gives a point on
\[
 \mathcal D:y^2=-240v^4+760v^2+81.
\]
The smooth projective model has no rational infinities since $-240$
is nonsquare. For $v\ne0$, the maps
\[
 X=\frac{9(y+9)}{8v^2}+\frac{95}{2},\quad Y=\frac{9X}{2v},
 \qquad v=\frac{9X}{2Y},\quad
 y=\frac89(X-95/2)v^2-9
\]
identify it with $E:Y^2=X^3-95X^2+2560X$.
The omitted points $(0,\pm9)$ map to $O,(0,0)$ respectively.
The maps extend to the smooth projective models.
The dual curve is $E':Y^2=X^3+190X^2-1215X$.

Use \eqref{eq:early-isogeny-quartic} and
\eqref{eq:early-isogeny-rank}. On $E$ the quadratic
$X^2-95X+2560$ is positive on the real line, since its discriminant
is $-1215$. Thus any nonzero real abscissa is positive and the
candidate classes are $1,2,5,10$. The classes 1 and 10 occur at
the two isogeny-kernel points. On $E'$ the classes are
$\pm1,\pm3,\pm5,\pm15$, with 1 and $-15$ occurring at its
kernel points. Every other class is excluded by this complete table
of primitive-chart right-side residues.
\begin{center}
\begin{tabular}{ccll}
\toprule
Curve & $d$ & Modulus & Residues\\
\midrule
$E$ & $2,5$ & $3$ & $2$\\
$E'$ & $-1,15$ & $16$ & $7,12,15$\\
$E'$ & $3,-5$ & $16$ & $3,11,12$\\
$E'$ & $-3$ & $128$ & $5,21,29,69,80,93,117,125$\\
$E'$ & $5$ & $128$ & $5,13,21,77,80,85,101,109$\\
\bottomrule
\end{tabular}
\end{center}
For each modulus $p^e$, the charts $(a,1)$ and $(1,b)$, $p\mid b$,
exhaust all primitive pairs up to a unit. Every listed residue is
nonsquare. Both descent images have order two, and the rank formula
therefore proves rank zero.

The integral discriminant is $-2^{22}3^5 5^3$. Good reduction gives
$\#E(\F_7)=\#E(\F_{11})=10$, bounding all rational
torsion by ten, including the two possibly exceptional primary parts.
The ten rational points
\[
 O,(0,0),(16,\pm144),(40,\pm120),(64,\pm192),(160,\pm1440)
\]
attain this bound and exhaust $E(\Q)$. The inverse formulas
give exactly
\[
 (v,y)=(0,\pm9),\quad(\pm1/2,\pm16),\quad(\pm3/2,\pm24),
\]
with the two signs independent. At $v=0$, the recovered value
$u^2=(y+20v^2+29)/8$ is $19/4$ or $5/2$, both nonsquares.
At $v=\pm1/2,\pm3/2$, a $B$ tripleton contains a central $A$
vertex $\pm1/2$. Thus no point of the complete quotient lifts.
The independent checker \cert{check_type2_exclusion.py} verifies the
maps, complete residue tables, point counts, point list and substitutions.
The additional check \cert{check_type2_exclusion.m} is independent of
the elementary rank proof.

\subsection{The three seven-adic disks for Type 14}\label{detail:type14}
For Type~14 the moment branch is $a=(5b+3)/3$, where $a$ is the
$A$ doubleton center square and $b$ the $B$ doubleton center square.
The other three $B$ squared centers have power sums
$p_k=C_3(k)+M_2(k,a)-M_2(k,b)$, where
$M_L(k,z)=\sum_{d\in O_L}(\sqrt z+d)^{2k}$ and
$C_3(k)=\frac12\sum_{d\in O_6}d^{2k}$.
Odd powers of $\sqrt z$ cancel, so these are rational polynomials.
The monic singleton cubic is determined by Newton's formulas
$e_1=p_1$, $e_2=(p_1^2-p_2)/2$,
$e_3=(p_1^3-3p_1p_2+2p_3)/6$.
Its discriminant is
$-4bh_3(b)/243$, where
\[
 h_3(b)=900b^5-7100b^4-21675b^3+71658b^2+3321b-93312.
\]
Since $b=v^2\ne0$ and the singleton cubic splits rationally, a valid
input gives a finite point on the smooth genus-two curve
\[
 C:y^2=-3h_3(b).
\]
The points $\infty,(9,\pm3240)$ are rational. The Jacobian
has rank at most one by the unconditional two-descent in
\cert{check_type14_exclusion.m}. The class-group proof callback is
the actual Minkowski bound for every encountered order. In particular
the quintic root field has discriminant $-7745836638000$ and proof
bound 136074. We establish rank one and the complete list of eligible
points by one local calculation.

Put $a_0=-2700$, $X=a_0b$, $Y=a_0^2y$ and
$F(X)=a_0^4(-3h_3(X/a_0))$. The monic model $Y^2=F(X)$ has smooth
reduction at seven, with only $\infty,(4,\pm3)$ as residue points.
The three rational points above already occupy those three disks.
Let $D=[(9,3240)-\infty]$, viewed on the monic model. Exact Cantor
arithmetic for the reduced Mumford representative $(U,V)$ of $25D$
gives
\[
\begin{gathered}
 \deg U=2,\quad U\equiv(X-5)^2\pmod7,\quad
 v_7(\operatorname{disc}U)=2,\quad
 \operatorname{disc}(U)/7^2\equiv5\pmod7,\\
 7V\equiv5X+3\pmod7,\qquad F(5)\equiv5\pmod7.
\end{gathered}
\]
These congruences are independently recomputed by
\cert{check_type14_seven_adic.py}. They include sufficient precision
to determine the logarithm direction.

Because the discriminant unit 5 is nonsquare, the two support
abscissas $x_1,x_2$ are in the unramified quadratic extension of
$\Q_7$, have difference of valuation one and both reduce
to 5. The ordinates $y_i=V(x_i)$ are units. The leading coefficient
of $V$ has valuation $-1$, so $y_1-y_2$ is a unit; its two reductions
are the opposite roots of 5. Hence $P_1=(x_1,y_1)$ and
$\iota P_2=(x_2,-y_2)$ lie in the same disk, and
$25D=[P_1-\iota P_2]$. For
$\omega_0=dX/(2Y)$ and $\omega_1=X\,dX/(2Y)$, the local integrals are
\[
 \left(\int_{\iota P_2}^{P_1}\omega_0,
       \int_{\iota P_2}^{P_1}\omega_1\right)
 \equiv\frac{x_1-x_2}{-2y_2}(1,x_2)\pmod{7^2}.
\]
Indeed $Y$ is a unit and its expansion in $X-x_2$ is integral.
For every primitive term of degree $k\geq2$, the valuation is at
least $k-v_7(k)\geq2$. The first logarithm coordinate has valuation
one and is nonzero, proving $D$ non-torsion and the rank exactly one.
If $\lambda$ is the second coordinate divided by the first, then
$\lambda\in\Z_7$ and $\lambda\equiv5\pmod7$.
The differential $\omega_1-\lambda\omega_0$ therefore annihilates
the entire rational Mordell--Weil group.

At each finite residue point $X=4\ne5$, this differential has unit
coefficient in the local coordinate $X-4$. At infinity $\omega_1$
is a unit and $\omega_0$ vanishes. Thus in all three rational residue
disks it has integral coefficients and unit constant coefficient.
Its primitive is injective on $7\Z_7$. To see this for any
two distinct parameters $t_1,t_2$, the linear term in their primitive
difference has valuation $v_7(t_1-t_2)$, whereas each degree-$k$
term, $k\geq2$, has valuation at least
\[
 v_7(t_1-t_2)+k-1-v_7(k)>v_7(t_1-t_2).
\]
For two rational points in a disk the primitive difference must vanish
because their Abel--Jacobi difference is a rational Jacobian point.
Each disk therefore contains at most its known rational anchor.
This proves $C(\Q)=\{\infty,(9,\pm3240)\}$.
The rank-one Chabauty intrinsic in the Magma certificate independently
returns the same set.

For a genuine input $b$ is finite, so $b=9$ and $a=16$. The
singleton polynomial becomes
\[
 K(X)=\frac{(4X-1)(4X-9)(4X-81)}{64}.
\]
The $A$ doubleton center has absolute value 4, the $B$ doubleton
center has absolute value 3, and its singleton absolute centers are
$1/2,3/2,9/2$. With the prescribed offsets both inputs are exactly
\[
 \{\pm1/2,\pm3/2,\pm5/2,\pm7/2,\pm9/2\},
\]
contradicting disjointness. The exclusion covers every rational height.

\subsection{Type 11 and the three disks at 37}\label{app:type11-details}

We give the global squareclass argument and the Chabauty calculation
underlying Appendix~\ref{app:type11}. Retain its center-square variables
$a,b,c$ and the sum and product $S,T$ of the two remaining squared
centers. In the two higher moment equations the linear subresultant
in $S$, modulo the conic of that appendix, is a nonzero multiple of
\[
 \bigl(16bc+4b-16c^2+56c-85\bigr)
                 \left(S-\frac{5b-4c+16}{3}\right).
\]
The resultant of the first factor with the conic in $b$ is
$16(4c-49)(4c-9)^2(4c-1)$. The exceptional values $1/4,9/4$
meet the central chain. At $c=49/4$ the conic gives $b=9,-123/2$.
The only nonnegative square choice $b=9$ gives $S=4,5$ in the two
uncanceled equations, respectively forcing $a=1/4,9/4$. These also
repeat a central vertex. This proves the displayed formula for $S$
for every valid input and justifies the slope through $(b,c)=(9,49/4)$.

For $w=(b-9)/(c-49/4)$ put
\[
\begin{aligned}
 D&=w^2+5w-5,& F_A&=25w^2-53w+27,\\
 F_B&=41w^2-70w+30,&F_C&=49w^2-37w-5,\\
 F_E&=(4w-5)(47w-40).
\end{aligned}
\]
The necessary center data are
\begin{equation}\label{eq:md11-param}
 a=-\frac{15F_A}{4D},\quad b=-\frac{3F_B}{2D},\quad
 c=\frac{F_C}{4D},\quad T=\frac{F_E(5w-4)^2}{2D^2}.
\end{equation}
The quadratic $F_B$ is strictly positive on the real line, so $b>0$
forces $D<0$. Write $w=m/n$ primitively with $n>0$, and homogenize
all quadratic forms to the pair $(m,n)$. A prime $p\notin\{2,3,5\}$
dividing $D_h$ does not divide $n$. Since
$\operatorname{Res}(D,F_A)=-36$, it does not divide $F_{A,h}$,
so $v_p(D_h)$ is even by the square condition on $a$. Primitivity
makes $D_h$ odd. If $3\mid D_h$, then $3\nmid n$ and
\[
 4D_h=9\left(\left(\frac{2m+5n}{3}\right)^2-5n^2\right).
\]
The parenthesis is a 3-adic unit because 5 is nonsquare modulo 3.
Consequently $v_3(D_h)=2$. The negative squarefree part of $D_h$
is therefore $d=-1$ or $-5$.

The four integer forms
$dD_h,-15dF_{A,h},-6dF_{B,h},dF_{C,h}$ are squares.
The complete projective test modulo 37 leaves slopes $23,30$ for
$d=-1$ and slopes $0,4$ for $d=-5$. The condition $T\ne0$
requires $2F_{E,h}$ to be square. At the first two slopes its squareclass
is represented by 19 and 8, both nonsquares modulo 37. Hence
\begin{equation}\label{eq:md11-class}
 D_h=-5k^2,\qquad w\equiv0,4\pmod{37}.
\end{equation}
The projective test also excludes $n\equiv0\pmod{37}$, which proves
integrality of $w$ at the prime used below.

Set $r=(47w-40)/(4w-5)$. The denominator cannot vanish for a valid
input because it would make $T=0$. The squareclasses in
\eqref{eq:md11-param} give a point on
\[
 C:Y^2=f(r)=r(r^2-11r-1)(7r^2-2r+18)(9r^2-69r-49).
\]
The residual slopes in \eqref{eq:md11-class} require $r\equiv8,0$
at 37. The curve has good reduction there, so precisely three disks
are eligible, containing $(0,0),(8,1500),(8,-1500)$.

A full two-Selmer computation with Minkowski proof bounds gives
dimension five. The four irreducible factors of degrees $1,2,2,2$
contribute three independent rational two-torsion classes. Thus
$\operatorname{rank}J(C)(\mathbf Q)\leq2$.
Define
\[
 P=[(8,1500)-\infty],\qquad
 D=\left(r^3-\frac{51}{4}r^2+\frac{23}{4}r+8,
              \frac{415}{8}r^2+\frac{785}{8}r+20\right).
\]
The second expression is a Mumford divisor with $Y=v(r)$ on its
support. Its support polynomial is squarefree, coprime to $f$ and
divides $f-v^2$. For $\omega_j=r^jdr/(2Y)$, $j=0,1,2$, exact
arithmetic modulo $37^2$ gives
\[
 \frac{\log(58624P)}{37}=(25,28,33),\qquad
 \frac{\log(58624D)}{37}=(3,30,12)\pmod{37}.
\]
Their cross product is $(12,21,0)$. Both global logarithms are
independent, so their span is the full rational logarithm space.
There is an annihilating differential $\sum\lambda_j\omega_j$ with
$(\lambda_0,\lambda_1,\lambda_2)\equiv(12,21,0)\pmod{37}$.

To specify the local calculation, translate each multiple by a
nonspecial effective divisor with three distinct abscissas and unit
ordinates. The following table lists its residue coordinates and the
Mumford endpoint of the translated multiple. Polynomial coefficients
are in descending order modulo $1369$.
\[
\begin{array}{c|c|c|c}
R&\text{auxiliary }(r,Y)&u_{B+58624R}&v_{B+58624R}\\ \hline
P&(1,-7),(3,8),(6,3)&(1,693,101,1129)&(1324,576,276)\\
D&(1,7),(6,-3),(8,-17)&(1,540,876,1173)&(808,1113,417)
\end{array}.
\]
Hensel lifting at each fixed abscissa supplies the auxiliary ordinate.
Unit-pivot Cantor composition and reduction~\cite{Cantor} certify the endpoints.
For doubling, the ordinate lift is $v+uh$ with
$h\equiv((f-v^2)/u)(2v)^{-1}\pmod u$. The endpoint support roots
satisfy $r_i'-r_i\equiv-u(r_i)/u'(r_i)\pmod{37^2}$ and lie on
the same sheets. Therefore the tiny integrals are
$(r_i'-r_i)r_i^j/(2Y_i)\pmod{37^2}$. All nonlinear primitive
terms have valuation at least two, which proves the stated logarithms.

On the two disks over $r=8$, the numerator of the annihilator is
$12+21\cdot8=32\pmod{37}$, a unit. On the Weierstrass disk use $Y$
as parameter, so $dr/(2Y)=dY/f'(r)$ and the constant is
$12/f'(0)=12/31\pmod{37}$, also a unit. Each primitive is injective
on its whole disk. Indeed a nonlinear term of degree $k\geq2$
in the difference at parameters $t_1,t_2\in37\mathbf Z_{37}$ has
valuation at least
\[
 v_{37}(t_1-t_2)+k-1-v_{37}(k)>v_{37}(t_1-t_2).
\]
The unit linear term cannot cancel. Integrals between rational points
in one disk vanish, so the three anchors exhaust the eligible disks.
At $r=0$, $w=40/47$ gives $F_E=0$ and $T=0$. At $r=8$, $w=0$
gives $c=1/4$. All three are invalid. The computations are
\cert{check_type11_genus3_rank_bounds.m} and
\cert{check_type11_thirty_seven_adic.py}, with the preceding elimination
and squareclass tables in \cert{check_all_shift_chain_types.py}.

\subsection{Type 13 and six forced Prym directions}\label{app:type13-details}

We expand the global rank contradiction used in
Appendix~\ref{app:type13}. Put
\[
 N(c)=-240c^3+1600c^2-260c+475,\quad
 \alpha=25/4,\quad\beta=1/4,
\]
and let $K=\mathbf Q(\theta)$ with $N(\theta)=0$. The product of
the three first-side singleton squares is $N(c)/144$. The two other
center squares satisfy $b=(25-4c)/6$, $c=w^2$, so every input gives
a point on the genus-two curve
\[
 H:y^2=6c(25-4c)N(c).
\]
Its two-Selmer group bounds the rank by two. The four classes
\[
 P=[(\beta,135)-\infty],\quad
 Q=[(19/4,2565/2)-\infty],\quad
 T_0=[(0,0)-\infty],\quad T_\alpha=[(\alpha,0)-\infty]
\]
are independent modulo doubling. To check this without assuming a
point-search basis, pass to the monic model by $X=5760c$,
$Y=5760^2y$. The two rational-root components of the Kummer images
of these four classes are
\[
 (10,-15),\quad(190,-15),\quad(285,-10),\quad(10,-15).
\]
Valuations at 2,3,19 leave only $P+T_\alpha$ as a possible nonzero
relation. In the cubic component its image is
$(\beta-\theta)/(\alpha-\theta)$. At the simple residue root
$\theta=6\pmod{11}$ the ratio is the nonsquare 10. Thus the four
classes are independent and form a basis modulo doubling. A nonzero
square value of $c$ forces the first rational Kummer coordinate to
have class 10. The same valuation test reduces its class to
$P$ or $T_\alpha$, giving precisely
\begin{equation}\label{eq:md13-twists}
 (c-\theta)/\delta\in K^{\times2},\qquad
          \delta\in\{\alpha-\theta,\beta-\theta\}.
\end{equation}
Their second rational coordinate is $-15$, and hence the condition
that $b$ be square is preserved. The group bounds and Kummer character
calculations are \cert{check_type13_product_covers.m} and
\cert{check_type13_product_covers.py}.

\subsubsection{The alpha class and divisibility at the ramified dyadic place}

The alpha cover is an intersection of quadrics over $K$,
$w^2-\delta z^2=\theta$, $4w^2+6v^2=25$. Taking the kernel over
$d\theta$, with $d=-6\delta$, on its associated elliptic curve gives
\[
 E':Y^2=X(X^2+aX+e^2),\quad
 a=2d(\alpha-2\theta),\quad e=d\alpha,
\]
with rational center map
\[
 c=\frac{\alpha(X^2+aX+e^2)}{(X-e)^2},\qquad
 \frac{\alpha}{c}=1+24\delta^2t^2,\quad t=-X/Y,
 \qquad \delta t=\pm\frac{v}{4w}.
\]
A valid lift makes $XY(X-e)\ne0$.
Use $K=\mathbf Q(u)$ with $u^3+u^2+315u+1983=0$ and
$\theta=(u^2-3u+342)/60$. Its unique dyadic place has uniformizer
$\pi=(u+1)^2/2$ satisfying
$\pi^3+314\pi^2+26632\pi-347778=0$. Normalize $v_\pi(2)=3$.
The valuations in $4w^2+6v^2=25$ give
$v_2(w)=-1$, $m=v_2(v)\geq1$. Also
$v_2(\delta)=-2$ and $v_2(e)=-3$. Put $s=X/e$.
Since $c-\alpha=4\delta s/(s-1)^2$ has valuation $2m-1>0$,
$s$ is not a unit. Translation by $(0,0)$ replaces $s$ by $1/s$
without changing $c$, so arrange $v_2(s)=-(2m-1)$.

On the integral model $x'=\pi^{10}X$, $y'=\pi^{15}Y$, the
coefficients of $x'^2,x'$ have valuations 4 and 2. The point is
formal, with
\[
 t'=-x'/y'=\kappa2^{m-1}h,\qquad
 \kappa=(\delta\pi^5)^{-1},\quad v_\pi(\kappa)=1,
 \quad h\in\mathbf Z_2^\times.
\]
The formal logarithm identifies this formal group $G$ with
$\pi\mathcal O$, and $\log(t')\equiv t'\pmod{4t'\mathcal O}$.
For a direct bound, set $-1/y'=t'^3U$. Then
$U=1+a't'^2U+b't'^4U^2$, and the invariant differential is
$(1+t'U'/(2U))dt'$. Its nonconstant terms lie in the ideal generated
by $a't'^2,b't'^4$, both of valuation at least six. Their powers of
$t'$ are even, so integration divides only by odd integers. The
logarithm also has an integral compositional inverse.

The global two-Selmer group of $E'$ has order four and injects on
localizing at this dyadic place. Its two local torsion points are
global: the other two would require a square root of
$\theta(\theta-\alpha)$, whose norm is
$35^2\cdot19/2^{10}$, a dyadic nonsquare. At the reference parameter
$t'=\kappa$, there is a point $P_*$ with trivial local Kummer
class, but every local half has Kummer class outside the localized
global Selmer image. To specify this assertion, choose
$z_*\equiv1\pmod{\pi^6}$ with $z_*^2=1-6/\delta$, and put
\[
 X_*=(5\delta(1+z_*)/2)^2,\qquad Y_*=-\delta X_*.
\]
It has $\delta(-X_*/Y_*)=1$. Its halves are calculated by the
identity, on $y^2=x(x^2+a_2x+a_4)$,
\[
 r^2=x,\quad p=-y/r,\quad s^2=2x+a_2+2p,\quad
 x_h=x+rs+p,\quad y_h=-y-(r+s)(rs+p).
\]
The square roots are certified by Hensel margins. On the integral
model $x=60^2X$, $y=60^3Y$, coordinate accuracy exceeding
$2v_\pi(y)+6$ preserves all Kummer squareclasses, since the three
factor valuations are nonnegative and sum to $2v_\pi(y)$.
The resulting injectivity and the two reference classes are checked
by \cert{check_type13_alpha_selmer.m}.

It follows that every global point locally divisible by two is
globally divisible by two. Every chosen local half of such a point
is global because all local two-torsion is global. Divide the
putative formal point $m-1$ times within $G$, obtaining a global
point $R$ whose logarithm is $\kappa h\pmod{4\pi\mathcal O}$.
Choose an odd integer $j\equiv h\pmod4$. Then $R-jP_*\in4G$,
so $R$ is locally divisible by two and has a global half. Its local
half class is that of $j$ times a half of $P_*$, up to rational
two-torsion. Both possible classes lie outside the global Selmer
image, a contradiction. This excludes the alpha class on its whole
rational parameter range.

\subsubsection{The signed-sum quartic and its divisor map}

For the beta class set $\delta=\beta-\theta$ and write the singleton
centers as $a,b,d$. Their squared sum $t=(a+b+d)^2$ satisfies
\[
 ((t-e_1)^2-4e_2)^2=64te_3,
\]
where $e_i$ are the elementary symmetric functions of $a^2,b^2,d^2$.
The homogeneous quadratic in $(c,t,Z)$ is
\[
 q_2=t^2-\frac{10c+50Z}{3}t+5c^2-30cZ+19Z^2.
\]
Let $q(c,Z)$ be the degree-two homogenization of
$(-N(c)/240)/(c-\theta)$ and put
$q_1=t(c-\theta Z)$, $q_3=-(320/3)q(c,Z)$.
The smooth plane quartic is $q_2^2=q_1q_3$. The two lines defining
$q_1$ cut twice distinct degree-two divisors. Equality of those
classes would supply a degree-two pencil on a smooth plane quartic,
which is impossible. Adjoining $\sqrt{q_1/(\delta Z^2)}$ therefore
gives a nontrivial étale double cover. Every beta lift maps to it by
\eqref{eq:md13-twists}. If $M_i$ is the symmetric matrix of $q_i$,
its Prym is the Jacobian $J/K$ of
\begin{equation}\label{eq:md13-prym}
 Y^2=-\delta\det(M_1+2XM_2+X^2M_3).
\end{equation}
This is Bruin's determinant construction~\cite[Theorem~5.1]{BruinPrym}, with twist obtained from
$q_1/\delta,q_2,\delta q_3$ and the change $X\mapsto\delta X$.
The twist and smoothness are checked in
\cert{check_type13_product_covers.py}.

Write $H(c,t)=q_2(c,t,1)^2-4tN(c)/9$ and $h=c-\beta$.
The necessary local reduction at 23 forces $h\in23\mathbf Z_{23}$.
For $h=0$ the input collides. The three distinct nonzero singleton
squares at that point allow their rational roots to be labeled and
signed so that $(a,b,d)\equiv(1/2,3/2,5/2)\pmod{23}$.
Use the ordered signed sums
\[
 v_1=-a-b+d,\quad v_2=a-b+d,\quad
 v_3=-a+b+d,\quad v_4=a+b+d.
\]
For every sum take $t=v_i^2$, $r=v_i z$, $s=q_2/r$, with the
same global square root $z^2=(c-\theta)/\delta$. At $h=0,z=1$
the sums are $1/2,3/2,7/2,9/2$. Let $D_i$ be their four Prym
images at this collision, and $A_i$ the images of a putative input.
We shall show that $D_1,D_2,D_3,A_1,A_2,A_3$ are independent
modulo torsion whenever $h\ne0$.

Here are formulas fixing both the divisor images and all signs.
Primes mean differentiation along $H(c,t)=0$. Put
\[
 a_0=t'/\delta-(r')^2,\quad
 b_0=(t')^2-10t'/3+5-r's',\quad
 k_0=-320\delta/3-(s')^2.
\]
The monic Mumford support on \eqref{eq:md13-prym} is
\[
 U=X^2+2\delta b_0X/k_0+\delta^2a_0/k_0.
\]
Its ordinate is the remainder modulo $U$ of
\begin{multline}\label{eq:md13-ordinate}
 \delta^2\left[
 (r'+Xs'/\delta)\left(\frac{c-\theta}{2\delta}
                      -\frac X\delta\left(\frac{10c+50}{3}-2t\right)\right)\right.\\
 \left. +(r+Xs/\delta)\left(-\frac1{2\delta}
                         -\frac X\delta(2t'-10/3)\right)\right].
\end{multline}
To derive it, take $V=(c,t,1)^t$, $W=(1,t',0)^t$ in the pencil
$Q=M_1/\delta-2xM_2+\delta x^2M_3$. Its tangent condition is
$a_0-2b_0x+k_0x^2=0$. Put $q_v=r-xs$, $q_w=r'-xs'$ and
$I=q_wV-q_vW$. At a tangent root, $QI=\lambda(V\times W)$.
The adjugate identity, together with the vector triple product, gives
$\lambda^2=-\det Q$. The second coordinate of $V\times W$ is
one, so it fixes the signed ruling coordinate. Changing to
$X=-\delta x$ and scaling this coordinate by $\delta^2$ gives
\eqref{eq:md13-ordinate}. This is one rational divisor map over $K$.
Changing $z$ to $-z$ negates the image.

\subsubsection{Equal signs and the analytic determinant}

Order the three split places over 23 by
$\theta\equiv11,13,21$. The reduction orders of $J$ are
$n_1=640,n_2=720,n_3=408$. In each two-coordinate block use
$\log(n_jR)/23$ in the differential basis
$dX/(2Y),X\,dX/(2Y)$. The three fixed collision rows form a
matrix $B$ with
\[
 \bar B=-\begin{pmatrix}
 4&11&0&15&7&4\\22&17&2&8&3&11\\20&4&21&20&12&22
 \end{pmatrix}.
\]
These rows are independent. Exact Jacobian arithmetic also verifies
$D_1=D_2+D_3-D_4$. To compute the rows, translate $n_jD_i$ to
an effective degree-two divisor with distinct integral support roots
and unit ordinates. A root displacement is $-U(x_0)/U'(x_0)$
modulo $23^2$, so first-order tiny integrals give the displayed
six-coordinate logarithms.

For the branch with $z(0)=1$ at every place the moving rows have the
convergent expansion
\[
 B+(h/23)(C_1+hC_2+h^2C_3+\cdots).
\]
The differential pullbacks along the signed-sum branch are
$(-2s/\delta,2r)dc/H_t$. They can be verified on the canonical
intersection of three quadrics: the Jacobian determinant of the
cover equations in $t,r,s$ is $\pm2H_t$. Hence the anti-invariant
canonical forms are spanned by $r\,dc/H_t,s\,dc/H_t$.
Tracing the support of the signed divisor formula at two collision
points with independent $(r,s)/H_t$ fixes the constants and signs.
Here $H_t,r,\delta$ are units on the disks. Thus the unintegrated
series are integral and $nC_n$ has integral coefficients for every
$n\geq1$.

Subtract the fixed rows and remove three factors $h/23$ from the
six-by-six determinant. The reduced determinant has an exact zero
of order at least two at $h=0$. For this assertion, write the three
derivative rows as $u_i(h)R(h)+v_i(h)S(h)$, where the scalars
$u_i,v_i$ are the same rational functions across the embeddings.
One can use $u_i=\sqrt{t_i}/H_t$ and
$v_i=q_2/(\sqrt{t_i}H_t)$, putting the dependence on $z,1/z,\delta$
and $n_j$ in the common embedding vectors $R,S$. A nonzero constant
combination kills both $u_i(0)$ and $v_i(0)$. After integration and
division by $h$, this combination has zero constant term and linear
term in the original two-plane. The other two derivative rows span
that plane at zero. Both the constant and linear determinant terms
therefore vanish. The first nonzero term is
\[
 \det\begin{pmatrix}\bar B\\
       \bar C_1+h\bar C_2+h^2\bar C_3\end{pmatrix}
        =7h^2\pmod{(23,h^3)}.
\]
It controls the full punctured disk. A coefficient of degree $k\geq3$
in this determinant has denominator valuation at most
$\lfloor k/22\rfloor$, since it is built from three integration
indices satisfying $\sum(n_i-1)=k$ and
$v_{23}(n_i)\leq(n_i-1)/22$. At $h\in23\mathbf Z_{23}\setminus\{0\}$
its valuation is strictly greater than $2v_{23}(h)$. The leading
term cannot cancel. This proves independence for equal local signs,
including their common global negative.

\subsubsection{Mixed signs and the two places above 1181}

Up to changing the global sign of $z$, a mixed pattern has one
negative two-coordinate block. Write $S_j$ for its block sign matrix
and $u=h/23$. The determinant of the six rows reduces to
$\det\binom{\bar B}{(\bar B+u\bar C_1)S_j}$. For the three choices
it is
\[
 u(11+20u),\qquad u(21+18u),\qquad u(10+3u)\pmod{23}.
\]
The exact constant term vanishes because at $u=0$ the bottom rows
differ from the top only within one two-dimensional block. Its
quotient by $u$ is an integral convergent series on $\mathbf Z_{23}$,
as each integrated coefficient carries $23^{n-1}/n$. Its constant
is a unit. Thus a dependence with $h\ne0$ requires
$u\equiv19,18,12$ respectively. In each case the six-row matrix
has rank five modulo 23, with normalized left null vector
\[
 (8,7,8,15,4,1),\quad(12,5,22,10,3,1),\quad(19,1,21,13,8,1)
\]
in the order $D_1,D_2,D_3,A_1,A_2,A_3$.

A dependence modulo torsion gives a primitive integer relation, whose
reduction is a nonzero multiple of one of these vectors. The torsion
has trivial 23-primary part because at a split place the good
reduction order is prime to 23 and the formal logarithm is injective.
The relation must therefore vanish in every reduction quotient of
order 23. Use the places over 1181 with $\theta=99,198$. A rational
pole of $c$ is impossible: the three ratios of singleton squares to
$c$ would split the squarefree residue cubic
$144X^3-240X^2-80X+240$, which has only one root there. The
Jacobian reduction orders are $1383036$ and $1338324$, each divisible
exactly once by 23. Multiplication by the order divided by 23 projects
to its order-23 subgroup. Relative to the projection of $D_1$, the
fixed rows have coordinates $(1,11,21)$ and $(1,16,17)$.

At each place enumerate every residue $c$ for which
$c,(25-4c)/6,(c-\theta)/\delta$ are squares and the singleton
cubic splits into three squares, with multiplicity allowed. Enumerate
all ordered signed roots $(a,b,d)$, evaluate the moving divisor
formula, and allow both signs of $z$ independently at each place.
The number of distinct residue tuples $(c,a,b,d)$ satisfying each
possible relation is
\[
\begin{array}{c|rrr}
 &\text{first}&\text{second}&\text{third}\\ \hline
 \theta=99&8&20&16\\
 \theta=198&18&16&22\\
 \text{intersection}&0&0&0
\end{array}.
\]
All the divisors in this finite calculation remain in the nonspecial
Mumford chart. Since $c,a,b,d$ are rational, their residue tuples
must agree at both places. No remaining relation is possible. This
proves that every noncollision beta lift would force six independent
global Prym classes. All divisor, logarithm, determinant and joint
reduction data are checked by \cert{check_type13_beta_prym.m}.

\subsubsection{The unconditional upper bound five}

The quartic factor in the Prym quintic defines an absolute field $L$
with polynomial
\[
\begin{split}
 P(X)={}&X^{12}-4X^{11}-68X^{10}-18698X^9+336998X^8-500324X^7\\
       &+7039932X^6-194015520X^5+1041300636X^4-3497537936X^3\\
       &+20028904184X^2-109223576168X+269574001992.
\end{split}
\]
Its signature is $(2,5)$ and its absolute discriminant is
\[ D_L=172707117717588564754287819295488000000. \]
The unconditional explicit-formula argument for the class group uses
the general-signature formula in \cite[Section~4]{miller}, with
support $X_0=2^{31}$ and
\[
 F(x)=\frac{\max(1-|x|/\ell,0)}{\cosh(x/2)},\qquad\ell=\log X_0.
\]
If the primes of norm at most $b=48000001$ generated a proper
subgroup, a prime-order cyclic quotient would give an everywhere
unramified extension split at all these primes and at the real
places. Subtracting its zeta explicit formula from that of $L$, and
dividing by the degree minus one, gives
\[
 \log D_L\geq A_{\rm arch}+S.
\]
Here each zero contributes nonnegatively throughout the critical
strip. The triangle has nonnegative Fourier transform, and so does
$\cosh(ax)/\cosh(x/2)$ for $|a|\leq1/2$, which proves this without
a hypothesis on zeros. A lower archimedean bound is
\[
 A_{\rm arch}\geq12\left(\gamma+\log(4\pi)
                  -\frac{\pi^2}{4\ell}\right)
                     +2\left(1-\frac{2\log2}{\ell}\right).
\]
A prime-power norm $t=q^m\leq X_0$ has unsigned weight
$4\log(t)(1-\log(t)/\ell)/(m(t+1))$. Its character multiplier
is one for $t\leq b$, and at least $-1$ elsewhere. Counting prime
ideal powers and taking rational interval lower bounds gives
$S>50.333$ and $A_{\rm arch}+S-\log D_L>0.080$, a contradiction.
The computation groups norms in 256 subdivisions of each dyadic
interval. For rational primes above $\sqrt{X_0}$ outside the
polynomial index, only degree-one primes contribute, and their number
is the number of simple polynomial roots. All other primes are
factored in the maximal order.

PARI's exact ideal relation test reduces every prime ideal of norm
at most $b$ to a factor base with norms at most 1000. Only the
outer rational-prime loop is divided among parallel workers. This
proves 1000 as the required generating bound independently of a
provisional class group. The bound is applied only after identifying
the absolute field with $L$. Other fields use their Minkowski bound,
which is 419 for $K$. The ordinary class groups are
$\operatorname{Cl}(L)=(\mathbf Z/2)^2$ and
$\operatorname{Cl}(K)=0$. The Prym two-Selmer group has dimension
six and its rational two-torsion dimension one. Therefore
$\operatorname{rank}J(K)\leq5$, contradicting the six independent
classes forced above. This completes Type 13.
The field and rank checks are \cert{check_type13_beta_prym_rank.m}.
The explicit-formula and ideal-relation computations are
\cert{type13_prym_class/prime_counts.cpp},
\cert{type13_prym_class/check_bound.py} and
\cert{type13_prym_class/cover.py}.

\subsection{Type 9 and a height comparison within a split fiber}\label{app:type9-details}

The singleton-root cubic in Appendix~\ref{app:type9} has a projective
linear identification with
\[
 E:y^2+y=x^3+x^2+295x+1109.
\]
For coordinates $[R:B:W]$ on that cubic and standard projective
coordinates on $E$, forward and inverse matrices are
\[
 L=\begin{pmatrix}-44&-84&347\\-1088&3102&4064\\16&6&-28\end{pmatrix},
 \qquad J=\begin{pmatrix}2884&7&36757\\-896&112&5152\\1456&28&5908\end{pmatrix}.
\]
Their product is a nonzero scalar matrix. The full group is
$E(\mathbf Q)=\mathbf ZG$ with $G=(31,202)$ and trivial torsion,
certified with unconditional class-group bounds by
\cert{check_type9_group.m}.

Three distinct rational roots at fixed finite $b=B/W$ are three
collinear points of the plane cubic. Their images are collinear on
the Weierstrass model, so they have indices $n_1+n_2+n_3=0$.
If $M=\max|n_i|$, one of the other absolute indices is at most
$M/2$. The point at infinity on $E$ maps to $[7:112:28]$, so the
fiber containing it is exactly $b=4$. We treat it separately.

Write $b=N/D$ primitively with $D>0$, $H=\max(|N|,D)$, and let
$H_C$ denote the maximum primitive coordinate of a point on the
original cubic. Its integral equation for $Dr$ has leading
coefficient 64, so $64Dr$ is an integer. After scaling by $H$, the
remaining coefficients are bounded by $1040/64,4396/64,4707/64$.
The root bound gives $|Dr|<80H$. The coordinates
$[64Dr:64N:64D]$ and their ratio imply
\[
 H\leq H_C\leq5120H.
\]
The maximum absolute row sums of $L,J$ are less than 9000 and
40000. Thus the primitive projective height $H_E$ on $E$ satisfies
$H_E\leq9000H_C$ and $H_C\leq40000H_E$.

For a finite elliptic point write $x=a/t^2$, $y=c/t^3$ in reduced
form, $t>0$, and $U=\max(\sqrt{|a|},t)$. The primitive projective
coordinates are $[at:c:t^3]$. The equation bounds
$|c|^2\leq |c|U^3+1406U^6$, giving the upper estimate below.
For the lower estimate, $|x|\leq100$ gives $U^3\leq1000t^3$.
The range $x<-100$ is impossible because the cubic is less than
$-1/4$, while $x>100$ implies $|y|>x^{3/2}/2$. Hence
\[
 U^3/1000\leq H_E\leq40U^3.
\]
Put $h(R)=\log_2\max(|a|,t^2)$. For any two points in one fiber,
these estimates show that
\[
 2^{3(h(R_i)-h(R_j))/2}
   \leq1000\cdot40\cdot9000\cdot40000\cdot5120<2^{60}.
\]
Their naive heights differ by less than 40, uniformly in $b$.

For the doubled canonical height
$\lambda(R)=\lim4^{-k}h(2^kR)$, the explicit duplication forms are
\[
\begin{aligned}
 F&=A^4-590A^2B^2-8874AB^3+82588B^4,\\
 D_2&=4A^3B+4A^2B^2+1180AB^3+4437B^4.
\end{aligned}
\]
There are integral Bézout identities
$dV^7=S_VF+T_VD_2$ for $V=A,B$, with
$d=2079016875$ and the sum of absolute coefficients of each pair
$(S_V,T_V)$ less than $2^{46}$. They imply
$\gcd(F,D_2)\mid d$ for primitive $A,B$. Taking the larger of
$|A|,|B|$ bounds the unreduced maximum below by
$d\max(|A|,|B|)^4/2^{46}$. The coefficient sums of $F,D_2$ are
less than $2^{17}$. Removing the gcd and telescoping yields
\[
 -46/3\leq\lambda(R)-h(R)\leq17/3.
\]
The numerator or denominator of $x(32G)$ has bit length 149, so
$\lambda(G)\geq(148-46/3)/32^2>1/8$.
By canonical-height quadraticity~\cite[Chapter~VIII, \S9]{Silverman},
heights of points within the fiber differ by
less than $40+(17+46)/3=61$. Comparing indices of maximum and
minimum absolute value gives
$3M^2\lambda(G)/4<61$. Since $3\cdot26^2/32>61$, all three
indices are bounded by 25. This is an all-height bound on any
split fiber.

Exact projection of the 51 multiples from $-25G$ to $25G$ gives
only two finite parameters occurring more than once,
\[
\begin{array}{c|l}
 b& (n,R/W)\\ \hline
0&(-2,25/4),(1,9/4)\\
4&(-3,49/4),(0,1/4),(3,81/4).
\end{array}
\]
At zero the cubic is $(4R-9)^2(4R-25)$, while at four it is
$(4R-1)(4R-49)(4R-81)$. These are also the only repeated-root
possibilities. Indeed the remaining factor of the discriminant has
unit leading coefficient at 7 and nonzero residues
$2,1,6,1,3,2,6$ on $\mathbf F_7$. A rational root would be
7-integral and reduce to a root. The only distinct split fiber is
therefore $b=4$, where the doubleton centered at 2 shares $3/2$
with the central chain. The Bézout forms, heights, exact multiples
and discriminant test are \cert{check_type9_exclusion.py}.

The ordered-root curve is the complete intersection in
$[r:s:B:W]$ defined by $f(r,B,W)=0$ and
$(f(r,B,W)-f(s,B,W))/(r-s)=0$. It is smooth of genus four.
At a double root the base derivative of $f$ and the derivative
$f''/2$ of the divided difference supply independent tangent
conditions. Above infinity the discriminant is $-675/4$, so there
are no rational points there. The two finite fibers above give all
nine rational points, including the three orderings at the repeated
fiber. This adds a complete rational-point consequence to the
height argument.

\subsection{Type 3 and a complete rational-image calculation}\label{app:type3-details}

In Appendix~\ref{app:type3}, the last center constraint has
$S_0=(20b-8c+27)/6$. If $S\ne S_0$, taking the divided residuals,
their resultant in $S$, and its resultant with the conic in $b$
leaves the factors
\[
\begin{split}
 &(4c-25)(4c-9)(4c-1)^4(80c^2-280c+281)\\
 &\hspace{12mm}\cdot(144c^2-616c+3025)(144c^2-328c+4489).
\end{split}
\]
The three quadratics have negative discriminants. The linear cases
are collisions or the nonsquare values $b=15/32,77/32$. Thus the
rational parameterization with $S=S_0$ covers every valid input.
For its parameter $z>0$, use the polynomials $L,A,B,C,V,D$ of
Appendix~\ref{app:type3}. If $z=N/M$ primitively with $M>0$, the
squarefree parts $d,k$ of $L(N,M)$ and $NM$ satisfy
$d\mid690$, $k\mid138$. The complete conditions are that
\[
 dL,\ dA,\ dB,\ dC,\ kNM,\ 3kV,\ D
\]
are integer squares. The resultant of $L$ and $A$ has prime support
$\{2,3,5,23\}$, and the outer coefficients of $V$ have support
$\{2,3,23\}$, which proves those finite squareclass ranges.
Both projective patches are included in the local test. Modulo
$2^{13}$ the remaining pairs are
$(1,6),(3,46),(115,46),(345,6)$. Imposing $3^7$ leaves just
$(1,6),(115,46)$.

For the first pair $6z$ and $L(z)$ are squares, giving a nonzero
point on
$E_6:Y^2=X(X^2-342X+59616)$. Its rank-zero group has only the
finite point $(0,0)$, contradicting $z>0$.
For the second pair, put $t=1656/z$. It is a nonzero square, and
\[
 10(t^2-57t+1656)\in\mathbf Q^{\times2},\qquad
 15(t^2-131t/8+3381)\in\mathbf Q^{\times2}.
\]
On $E_{10}:Y^2=X(X^2-570X+165600)$ there is therefore a point
with $X=10t$. The rational group has rank one and torsion generated
by $(0,0)$. The point $(160,4000)$ has two-isogeny class 10,
independent of the torsion class 46. Thus the abscissa class 10
forces its difference from this point to be twice a rational point.

Over $K=\mathbf Q(\sqrt{-15})$, put
$r=(57+15\sqrt{-15})/2$. The abscissa character at $(10r,0)$
implies $(t-r)/(16-r)\in K^{\times2}$. One can see the character
homomorphism directly: for three collinear points the product of
their $X-10r$ factors is the square of the line ordinate at $10r$.
Let $q(t)=t^2-131t/8+3381$ and $l=64(16-r)3/5$.
The two square conditions now give a point with $X=l(t-r)$ on
\[
 E:Y^2=X\bigl(X^2+lq'(r)X+l^2q(r)\bigr).
\]
Writing $s=\sqrt{-15}$, its two coefficients are
$45300-18900s$ and $488160000s-4186080000$.
The two-Selmer group has order four and torsion is generated by
$(0,0)$. The point $P=(X_0,360X_0)$,
$X_0=7200s-26400$, is non-torsion, so the rank is one.

The subgroup generated by $P$ and torsion has odd index. At the
good place $(19,s-17)$, both $X(P)$ and $X(P+(0,0))$ reduce to
nonsquares, so neither translate is a double. Elliptic Chabauty for
the map $t=r+X/l$ returns three subgroup points with rational
image, and an index obstruction equal to a power of two (4 in the
stated run). The obstruction is coprime to the subgroup index, so
the rational images are exhausted. They are $\infty$ and 16.
This uses the rational-image index conclusion of elliptic Chabauty,
so an unproved full Mordell--Weil basis is unnecessary.
Every rank and class-group computation uses unconditional proof
bounds. At the only finite image $t=16$, one has
$z=207/2$ and $(a,b,c)=(81/4,25/4,25/4)$. The singleton at
$5/2$ meets the center of the tripleton on the other side.
The elimination and squareclass tables are
\cert{check_type3_exclusion.py}, and the complete rational-image
certificate is \cert{check_type3_exclusion.m}.

\subsection{Type 17 and the primitive dyadic direction}\label{app:type17-details}

The product elliptic curve $E_0$ of Appendix~\ref{app:type17} has
full rational group generated by $G=(11025/16,165375/64)$ and no
torsion. On odd multiples, the cubic-root abscissa character fixes
$d=1/16-r$, where
\[
 w^3+w^2-240w+1120=0,\qquad
 r=(2w^2+5w+140)/3675.
\]
The associated curve $Y^2=X(X+5d)(X-rd)$ has full two-torsion,
Selmer order eight and rank one, with non-torsion point
$P=(d/16,9d^2/16)$. At the two good places above 17 with
$w=6,11$, none of the four translates of $P$ by two-torsion is a
double. The generated subgroup consequently has odd index.
Elliptic Chabauty with index obstruction 16 gives all rational
images $X/d$, namely $\infty,0,-5,1/16$. Only the last is finite
and positive, giving $b=16$, a collision. This exhausts the entire
odd-multiple class.

For the even class, write the product point as $2Q$, put $a=11025$
and $T=X(Q)$. Retain the cubics and quartics $g,A,B$ of
Appendix~\ref{app:type17}, so
\[
 x=\frac{A(T)}{4ag(T)},\qquad x+5=\frac{B(T)}{4ag(T)}.
\]
All three of $g(T),A(T),B(T)$ are nonzero squares. The odd-degree two-cover descents~\cite{BruinStoll} on $y^2=gA$ and $y^2=gB$ have sizes two
and one respectively, each containing the identity. The nonidentity
class in the first set projects to a nonsquare in the cubic field
of $g$. A representative in its cubic-times-quartic étale algebra
is divided by its total norm before projection. For an actual
abscissa class that norm is square. Multiplication of the original
representative by a rational scalar changes the normalized class
by a sixth power, which is square. Thus the normalization is
well defined. The two descents share the cubic-root character
of $T-r$. The second descent makes it trivial, selecting the
identity in the first as well. This proves that $T-\theta$ is
square for a root $\theta$ of $A$, and that $Q$ is an even
multiple of $G$.

On the integral equation of $E_0$, the generator has formal
parameter $-4/15$. The formal subgroup on $4\mathbf Z_2$ is
closed under addition and duplication starts with twice the
parameter. Therefore for $s=-T/Y(Q)$,
\[
 k=v_2(s)\geq3,\qquad v_2(T)=-2k.
\]
Take $K=\mathbf Q(w)$, $w^4-194w^2-96w+8302=0$, and
$\theta=3(2w^3+4w^2-240w-532)$. The monic cubic
$h(U)=A(U)/(U-\theta)$ is integral and irreducible over $K$.
The two square conditions give a global point $R=(T,Y')$ on
$E:Y'^2=h(T)$.

Its two-division root field has absolute polynomial
\[
 Z^{12}-472Z^{10}+75282Z^8-4450612Z^6+
                  64570537Z^4-277288596Z^2+234027000.
\]
It is totally real, with discriminant
$58160039859394517078535559761100800000$. The unconditional
explicit-formula estimate used above, now with support and
norm bound $2^{20}$, has margin greater than $0.262$.
Every prime-power term in the support has multiplier one in the
hypothetical cyclic quotient. PARI's exact ideal relations reduce
the generating set to norm 1000. After field identification this
gives class group $\mathbf Z/2$ and
$\dim_{\mathbf F_2}\operatorname{Sel}_2(E/K)=2$.

There is one dyadic place of $K$, totally ramified of degree four.
Write $D=K_v$, let $\pi$ be a uniformizer and $\mathfrak p$ its
maximal ideal. The integral basis is $1,\pi,\pi^2,\pi^3$.
The cubic $h$ has exactly one root in $D$, so
$\dim E(D)/2E(D)=5$, four from the local Lie group and one from
local two-torsion. Put $m=105\cdot2^{12}$ and
$\Lambda=m\log E(D)$. Five explicit points with independent local
Kummer images span this quotient. Their scaled logarithms span
$\Lambda/2\Lambda$, since a torsion point has logarithm zero.
For each point, the formal parameter of its $m$-multiple has
$v_\pi\geq44$. On the displayed integral model the differential
has integral even coefficients, so its logarithm is
$z+O(z^3)$ with dyadically integral coefficients. Its leading binary
coordinates may therefore be obtained from the formal parameter.
The five rows in the basis
\[
 (2^{12},2^{11}\pi,2^{11}\pi^2,2^{11}\pi^3)
\]
are integral with residue rank four. It follows that
$\Lambda=2^{11}\mathfrak p$. Expressing the two localized global
Selmer classes in the five local Kummer generators and then using
their logarithm residues gives a two-dimensional image in
$\Lambda/2\Lambda$. Adjoining the rational axis $(1,0,0,0)$
raises the rank to three. This simultaneously proves injectivity
of the Selmer-to-logarithm map and exclusion of that primitive
axis.

The chain equations force exactly the excluded direction. For
$z=-T/Y'$ the ratio satisfies
\[
 (z/s)^2=\frac{g(T)}{h(T)}\in1+2^{2k}\mathcal O_D.
\]
Both cubics are monic integral and $v_2(T)=-2k$, which proves this
congruence by dividing them by $T^3$. Of $z/s-1,z/s+1$, the one
with smaller valuation has valuation one, using the normalization
$v_2(2)=1$. Thus for one sign
\[
 z=\pm s(1+2^{2k-1}\eta),\qquad\eta\in\mathcal O_D.
\]
The same leading congruence holds for $\log_E R$. Consequently
$2^{-k}m\log_E R$ is primitive in $\Lambda$ and its residue is
$(1,0,0,0)$.

To apply the global Selmer obstruction, divide $R$ by the largest
possible global power of two, writing $R=2^jR_0$. The nonzero
logarithm and finite generation make $j$ finite. A globally
primitive point has nonzero Kummer class, hence a nonzero residue
in $\Lambda/2\Lambda$ by injectivity of the map just computed.
Therefore $m\log_E R_0$ is primitive, forcing $j=k$.
Its residue would be the excluded axis. This contradiction
exhausts the even branch and proves Type 17 at all rational
heights. The two-cover projections, local generator independence,
logarithm valuations and final three-dimensional rank assertion
are all checked by \cert{check_type17_even_branch.m}.
The odd cover is \cert{check_type17_odd_branch.m}, and the
unconditional relation range is
\cert{check_type17_factor_base.gp}, together with the
\texttt{--type17} mode of
\cert{type13_prym_class/check_bound.py}.

\subsection{Type 12. Two cubic-field curves of rank one}\label{detail:type12}

\subsubsection{The two parity classes}\label{arith:61}

Write \(b\) for the square of the length-four B center.
A valid type-12 lift gives a point \(mP_1+nP_9\) with \(m+n\) odd on the
rank-two elliptic curve
\[E_{12}:Y^2=X^3-240X^2+215aX+90a^2,\qquad a=160,\quad X=ab,\]
where \(P_1=(160,-2400)\) and \(P_9=(1440,50400)\).

\emph{Proof.} From \S\ref{app:type12}, \(b\) and \(10-b\) are squares. The three
singleton squares have product \(f(b)/64\), where
\(f(b)=160b^3-240b^2+215b+90\). Indeed their monic cubic, multiplied
by 64, is
\[64Z^3-(240b+320)Z^2+(320b^2+280b+436)Z-f(b).\]
Thus \(f(b)\) is a square, giving the point with \(X=ab\).
The full Mordell--Weil group is free on \(P_1,P_9\).

Let \(w^3+38w-504=0\) and \(r=(4-w)/8\), so \(f(r)=0\).
At the ramified prime of \(\mathbf Q(w)\) over two, take the integer
valuation with \(v(2)=2\). We have \(v(w)=1,\ v(r)=-5,\ v(a)=10\).
The equation \(b+(10-b)=10\) expresses 10 as a sum of two rational
squares. Their roots must both be dyadic units. Indeed unequal
valuations give an even valuation for the sum, while two roots of equal
valuation \(j\) give valuation \(2j+1\), since two odd squares sum to
2 modulo 8. Hence \(v(b)=0\). The Kummer character at the root \(ar\)
is represented by \(a(b-r)\), whose valuation is five. The same holds
at both generators (with \(b=1,9\)). Valuation modulo two is additive
on the Mordell--Weil group, so \(m+n\equiv1\pmod2\).

For a surviving class indexed by \(b_0=1\) or \(9\), put \(d=b_0-r\).
The same abscissa character makes \((b-r)/d\) square. Multiplying by the
two center squares gives a point with \(X=-db\) on
\[Y^2=X(X+10d)(X+rd).\]
The two-torsion has order four, and unconditional two-descent gives
rank upper bounds three and four for the two covers, respectively.
\cert{check_type12_center_parity.m}

\subsubsection{Two rank-one curves and their rational images}\label{arith:62}

Suppose that \(b,10-b\) are rational squares and the
singleton cubic in \S\ref{arith:61} splits into three rational squares,
counted with multiplicity. Then \(b\in\{1,9\}\). Both values give a
collision. Type 12 has no valid lift at any rational height.

\emph{Proof.} Retain \(K,r,f\) from \S\ref{arith:61} and write
\[q(B)=f(B)/(160(B-r))=B^2+q_1B+q_0,\quad
 q_1=r-3/2,\qquad q_0=r^2-3r/2+43/32.\]
The proof of \S\ref{arith:61} uses just the square conditions in the
statement. Thus for one of \(b_0=1,9\), the quantity \((b-r)/(b_0-r)\)
is a square in \(K\). Since \(f(b)\) is square, \(160(b_0-r)q(b)\)
is also square. The rational number \(b\) is nonzero since 10 is
not a rational square.

For \(b_0=1\), put \(d=-160(1-r)\). The product \(d(b-10)q(b)\) is a
square. With \(x=d(b-10)\) it gives a point on
\[E^{(1)}:\quad y^2=x(x^2+dq'(10)x+d^2q(10)).\]
For \(b_0=9\), put \(D=160(9-r)\). This time take the square
\(D b(10-b)q(b)\). Division by \(b^4\) gives the square
\(D(10t-1)(q_0t^2+q_1t+1)\) with \(t=1/b\). Put \(c=10Dq_0\) and
\[F(U)=(U-c/10)(U^2+10Dq_1U+10Dc).\]
Then \(F(ct)=c^2D(10t-1)(q_0t^2+q_1t+1)\). We obtain a point on
\(E^{(9)}:y^2=F(x+c/10)\) with \(x=c/b-c/10\). In each case let
\(P_0\in E^{(b_0)}(K)\) be a point over \(b=b_0\).

\textbf{Ranks.} For the first curve do two-descent directly. For the second
do it on the quotient by the order-two subgroup generated by \((0,0)\)
and use the image of \(P_0\). Each descent gives a 2-Selmer group of
order 16 and rational two-torsion of order 2. The quotient of the
Selmer group by torsion and the specified point has dimension two.
On this quotient the alternating Cassels--Tate pairing takes the
nonzero value on a pair of basis elements. The image of the global
points lies in the radical, so the rank is at most one. Each \(P_0\)
has infinite order, and isogeny preserves rank. Thus both original
curves have rank one over \(K\).

The second descent uses the Minkowski proof bound throughout. For the
first, the quadratic cofactor of the two-division algebra has absolute
field \(L\) defined by \(T^6+228T^4+12996T^2+7077920\). PARI verifies
its discriminant \(-5540347111683392000\), checks that every ideal of
its factor base has norm at most 100, and certifies that the factor base
generates the class group. Magma uses the resulting proof bound 100
after testing field isomorphism, and the Minkowski bound elsewhere.
\cert{check_type12_fields.gp}; \cert{check_type12_exclusion.m}; \href{https://pari.math.u-bordeaux.fr/dochtml/html-stable/General_number_fields.html\#bnfcertify}{PARI class-group certification}

\textbf{The residue disks.} Take \(p=907\). Since \(-1\) is a nonsquare,
a pole of \(b\) is impossible. Indeed \((10-b)/b\) is square and would
reduce to \(-1\). The field \(K\) splits completely at \(p\). For each
residue \(b\), test the two center squares, the three squareclasses
\((b-r_j)(b_0-r_j)\) at the three roots \(r_j\), and whether the
singleton cubic has three roots counted with multiplicity, each a square in
\(\mathbf F_p\). These reductions are necessary since the cubic
has unit leading coefficient. In each class the exact enumeration of
the 907 residues leaves only \(b\equiv b_0\pmod p\).

For the three embeddings \(j:K\hookrightarrow\mathbf Q_p\), let \(n_j\)
be the order of the reduced elliptic curve and write
\(\ell_j(R)=\log_E([n_j]R_j)\), using the differential \(dx/(2y)\).
The displayed models have good reduction, and \(y(P_{0,j})\) is a unit
at all three places. With \(x(b)=d(b-10)\) or \(c/b-c/10\), compute
\[L_j=\ell_j(P_0)/p\bmod p,\qquad
 V_j=n_j x'(b_0)/(2y(P_{0,j}))\bmod p .\]
For \(L_j\) the formal parameter \(-x/y\) of \([n_j]P_{0,j}\) suffices
modulo \(p^2\). The certificate gives the following rows in a common
ordering of the embeddings.

{
\begin{longtable}[]{@{}
  >{\raggedright\arraybackslash}p{(\linewidth - 6\tabcolsep) * \real{0.2500}}
  >{\raggedright\arraybackslash}p{(\linewidth - 6\tabcolsep) * \real{0.2500}}
  >{\raggedright\arraybackslash}p{(\linewidth - 6\tabcolsep) * \real{0.2500}}
  >{\raggedright\arraybackslash}p{(\linewidth - 6\tabcolsep) * \real{0.2500}}@{}}
\toprule\noalign{}
\begin{minipage}[b]{\linewidth}\raggedright
\(b_0\)
\end{minipage} & \begin{minipage}[b]{\linewidth}\raggedright
\((L_1,L_2,L_3)\)
\end{minipage} & \begin{minipage}[b]{\linewidth}\raggedright
\((V_1,V_2,V_3)\)
\end{minipage} & \begin{minipage}[b]{\linewidth}\raggedright
\(L_iV_j-L_jV_i\) for \(ij=12,13,23\)
\end{minipage} \\
\midrule\noalign{}
\endhead
\bottomrule\noalign{}
\endlastfoot
1 & 823, 146, 636 & 453, 716, 192 & 698, 516, 760 \\
9 & 236, 627, 838 & 800, 274, 656 & 238, 499, 300 \\
\end{longtable}
}

For any global point \(R\), rank one gives a rational \(k\) such that
\(\ell_j(R)=k\ell_j(P_0)\) at every place. This follows by clearing
denominators in the group modulo torsion. At each place a candidate
reduces to \(\epsilon_jP_{0,j}\) with \(\epsilon_j=\pm1\). Choose two
places \(i,j\) with the same sign. Set \(A_i=\ell_i(P_0)/p\) and likewise
for \(j\). The function \(A_i\ell_j(R(b))-A_j\ell_i(R(b))\) vanishes
at the anchor with those signs, and at any global candidate on these
two local branches. Its derivative at the anchor reduces to
\(\epsilon_i(L_iV_j-L_jV_i)\), a unit in every case.

To see uniqueness on the entire disk, put \(h=b-b_0\). The derivatives
of the logarithms are \(n_j x'(b)/(2y_j(b))\). As power series in \(h\)
they have integral coefficients, since the ordinates are units and
Hensel lifting takes their square roots with unit derivative. For the
reciprocal formula \(b_0=9\) is a unit as well. Integration gives terms
of degree \(k\) in \(h\) with coefficients in \(k^{-1}\mathbf Z_p\).
For \(h\in p\mathbf Z_p\) and \(k\ge2\),
\((k-1)v_p(h)>v_p(k)\) when \(h\ne0\). The linear term therefore has
strictly smaller valuation than every higher term. This proves
\(b=b_0\) in both classes.

The singleton squares at the two anchors are respectively
\(\{1,9,25\}/4\) and \(\{25,49,81\}/4\).
At \(b=1\) the B length-four chain centered at 1 contains \(1/2\); at
\(b=9\) the one centered at 3 contains \(3/2\). Both vertices belong to
the centered A chain. The exceptional branches \(b=0,1,4\) of the fourth
moment already give collisions by \S\ref{app:type12}. This proves the type exclusion.

\subsection{Type 1. Split descent and a further division}\label{detail:type1}

\subsubsection{The parity of the squared center}\label{arith:63}

Put \(b=v^2\) for the square of the length-two B
center in type 1, and set \(a=16384\) and
\[f_1(B)=aB^3+38672B^2+33144B+11025,\qquad
 E_1:Y^2=X^3+38672X^2+33144aX+11025a^2.\]
For any valid lift, the point with \(X=ab\) lies in one of two classes
of \(E_1(\mathbf Q)/2E_1(\mathbf Q)\), represented by the points above
\(b=0,1\). Fourteen of the sixteen global classes are excluded.

\emph{Proof.} By \S\ref{app:type1}, \(5b+4\) is square. The constant coefficient of the
singleton cubic of \S\ref{app:type1} is \(-f_1(b)\), after scaling its leading
coefficient to 64. Hence \(f_1(b)\) is square, and
\((ab,a\sqrt{f_1(b)})\) gives the claimed point. Changing its ordinate
sign has no effect modulo twice the group. The full group is torsion-free
of rank four. The certificate computes a basis and represents its sixteen
parity classes by sums with coefficients zero or one.

At each prime in the table below, reduce into
\(E_1(\mathbf F_p)/2E_1(\mathbf F_p)\). For integral \(b\), retain the
residues where \(b,5b+4\) are squares and all three singleton roots
are squares, allowing multiplicity. For a pole, \((5b+4)/b\) reduces
to 5 and is square. The three singleton squares divided by \(b\) are
squares and integral at \(p\), since their cubic after division by
\(b^3\) has integral coefficients and unit leading coefficient. They
must reduce to the roots of \(64z^3-1152z^2+7232z-16384\). Retain the
group identity for a pole whenever both necessary tests pass.

{
\begin{longtable}[]{@{}llllll@{}}
\toprule\noalign{}
Prime & 7 & 67 & 73 & 107 & 257 \\
\midrule\noalign{}
\endhead
\bottomrule\noalign{}
\endlastfoot
Global classes left after intersection & 8 & 6 & 5 & 4 & 2 \\
\end{longtable}
}

Exact rational division by two identifies these last classes with the
two stated representatives. For a root \(r\) of \(f_1\) the abscissa character
at \(ar\) now makes \((b-r)(b_0-r)\) square in \(\mathbf Q(r)\).
At 73 the polynomial \(f_1\) splits with distinct roots, and 5 is a
nonsquare, so there is no pole. Testing the character at all three roots
and imposing the center and singleton squares gives just \(b=b_0\)
modulo 73 in each class. Sections~\ref{arith:64} and~\ref{arith:65} determine the eligible
centers in these disks.
\cert{check_type1_parity.m}

\subsubsection{Split descent in the cubic field}\label{arith:64}

In the type-1 parity class represented by \(b_0=1\) in
\S\ref{arith:63}, the only squared center satisfying the necessary square
and splitting conditions is \(b=1\). It gives overlapping chains. The
cubic-field elliptic cover for \(b_0=1\) has rank exactly one. For
\(b_0=0\) the corresponding cover has unconditional rank at most two.

\textbf{The two covers.} Put \(K=\mathbf Q(w)\), where
\(w^3-w^2-330w-149578=0\), and set
\(r=(-w^2+20w-5426)/7168\). Then \(f_1(r)=0\). By \S\ref{arith:63},
\((b-r)(b_0-r)\) is square in \(K\). Since \(b\) and \(5b+4\) are
rational squares, \(d b(b+4/5)(b-r)\) is square for \(d=5(b_0-r)\).
Consequently a lift gives a point on the following elliptic curve:
\[E_d:Y^2=X(X-dr)(X+4d/5),\qquad X=db .\]

\textbf{The global 2-Selmer groups.} Use the roots
\((e_1,e_2,e_3)=(0,dr,-4d/5)\). The split Kummer map has coordinates
\(X-e_i\) modulo squares, with their product square. At a two-torsion
point replace its zero coordinate by the product of the other two.
Take \(S\) to consist of all prime ideals over
\(2,3,5,7,11,12101,1122137\). Outside \(S\) the roots are integral with
unit differences and the residue characteristic is odd. There the
local image consists of the unramified triples of square product.
Indeed the image is unramified by good reduction, and both spaces have
dimension two. Thus the global classes lie in two copies of
\(K(S,2)\), the squareclasses of even valuation outside \(S\).

PARI certifies the class group and units of \(K\). The \(S\)-class
group is trivial, so the fundamental \(S\)-units and the torsion unit
span \(K(S,2)\). They give dimension 17. At each finite place in \(S\),
the certificate constructs independent Kummer images of local points
until their number is \(2+[K_v:\mathbf Q_2]\) at a dyadic place or
two at an odd place. These are the exact dimensions of
\(E_d(K_v)/2E_d(K_v)\), by the local Lie group structure and the four
two-torsion points. Local unit characters at \(v\mid2\) are computed
modulo \(\mathfrak p^{2v(2)+1}\), which suffices by Hensel's lemma.
At the real embedding the two sign patterns are given by the unbounded
component and by an abscissa between the two smallest roots.
Intersecting these linear conditions gives Selmer dimensions 8 and 7
for \(b_0=0,1\), respectively. The point selection stops only after
the required local dimension is attained.

\textbf{The Cassels--Tate pairing.} For a Selmer triple \(s=(s_1,s_2,s_3)\)
use the two-cover \(s_i z_i^2=X-e_i\). For cyclic indices \(i,j,k\)
let \(L_i\) be the tangent at a specified \(K\)-point of the conic
\[s_j z_j^2-s_k z_k^2+(e_j-e_k)t^2=0 .\]
Cassels' formula for the pairing against another triple \(s'\) is the
product over the places of the three Hilbert symbols
\((L_i(P_v),s'_i)_v\), where \(P_v\) is a local point of the cover.
The data file supplies four triples on each curve and three conic
points per triple. The verifier checks the conic identities and the
square-product conditions. It constructs local points at the bad
places, at primes supporting the triples and at primes supporting the
fractional ideal generated by the coefficients of each tangent.
At the remaining odd primes the reductions of the conics are smooth,
their tangent forms are primitive, and all entries of the triples
are units. On the conic a tangent cuts out twice its point of contact.
Thus the evaluation of a tangent on primitive integral coordinates
has even valuation whenever it is nonzero. The three conic coordinates
are primitive for a primitive point of the intersection, by the other
quadric equation. Changing projective scale has trivial total symbol
because the target triple has square product. This accounts for every
place. The verifier uses Hensel error bounds for the finite symbols
and rational isolating intervals for the real signs.
\href{https://www.dpmms.cam.ac.uk/~taf1000/papers/casselspairing.pdf}{Cassels pairing over number fields}

Using 1 for a negative symbol product, the resulting pairing minors
in the order of the data file are
\[\begin{pmatrix}0&1&0&0\\1&0&0&1\\0&0&0&1\\0&1&1&0\end{pmatrix},
 \qquad
 \begin{pmatrix}0&1&1&1\\1&0&0&1\\1&0&0&0\\1&1&0&0\end{pmatrix}.\]
Both are invertible over \(\mathbf F_2\). The images of global points
are in the radical. Subtracting four and then the two torsion
dimensions gives rank bounds two and one. For \(b_0=1\),
\(P=(d,3d^2/5)\) has infinite order, proving equality.
\cert{type1_descent/verify.gp},
\cert{type1_descent/tangents.gp}

\textbf{The disk at 73 for \(b_0=1\).} \S\ref{arith:63} gives
\(b\equiv1\pmod{73}\). At the three embeddings of \(K\) into
\(\mathbf Q_{73}\) the model is good and \(Y(P)\) is a unit.
Use the notation of \S\ref{arith:62}, with \(n_j=\#E_d(\mathbf F_{73})\) at
place \(j\). For the chart \(X=db\), the rows
\(L_j=\log([n_j]P_j)/73\bmod73\) and
\(V_j=n_jd/(2Y(P_j))\bmod73\) are
\[L=(30,34,22),\qquad V=(55,61,36).\]
Their pairwise determinants are \(33,16,28\) modulo 73. A global point
has logarithm proportional to that of \(P\) at all three places, with
the same rational proportionality factor. At least two embeddings
reduce to the same sign of \(P\). The corresponding determinant
function on that disk vanishes at 1 and its derivative is a unit.
Its derivative has integral power-series coefficients, so the
linear term dominates on the entire disk by the valuation estimate
of \S\ref{arith:62}. Thus \(b=1\). The B length-two chain centered at 1 shares
\(1/2\) with the fixed A chain. This excludes the class at every
rational height.
\cert{type1_descent/disk73.m}

\subsubsection{The discriminant-field cover}\label{arith:65}

Type 1 has no valid distinct-entry lift at any rational height.
In the two disks of \S\ref{arith:63} the only squared centers satisfying the
necessary center and singleton conditions are respectively 0 and 1.

\textbf{The discriminant character.} \S\ref{arith:64} handles the class of 1.
In the class of 0 write \(b=v^2\). If \(v\ne0\), then \(k=v_2(v)\ge3\).
In fact \(k<0\) would make \(5+4/v^2\) a square congruent to 5 modulo 8.
At the ramified prime of the cubic field of \S\ref{arith:64}, the normalized
valuation of its root \(r\) is \(-4\), and \(1-1/r\) is nonsquare.
For odd \(v\), substituting \(v^2\) for 1 does not change this local
squareclass, as \(v^2-1\in8\mathbf Z_2\). This contradicts \S\ref{arith:63}.
For \(k=1,2\), respectively, \((5v^2+4)/4\) is 6 or 5 modulo 8.

Use the sextic
\[h(B)=-800B^6+2625B^5-12050B^4+529B^3+54448B^2+53744B+14400.\]
The discriminant of the singleton cubic with leading coefficient 64 is
\(2^{26}h(b)\), so \(h(b)\) is a square. The sextic is irreducible.
Let \(F=\mathbf Q(\theta)\) with \(h(\theta)=0\). We claim that
\(s=1-b/\theta\) is square in \(F\).

Its norm is \(h(b)/14400\). Away from the primes above
\(D=\{2,3,5,23,41,2111,\allowbreak 11801,144461\}\), the polynomial has unit leading
and constant coefficients and unit discriminant. For integral \(b\), at
most one prime over a given rational prime divides \(b-\theta\), and it
has residue degree one. The square norm forces that valuation to be even.
At a pole the valuation is even because \(b=v^2\). This puts \(s\) in
the group of squareclasses unramified outside \(D\).

The discriminant of \(F\) is 14755078257381100. PARI certifies that its
factor base generates the class group, and all its generators are
principal. For each set of places used in this proof, the verifier
checks the expected number of independent unit squareclasses by
valuations and local squareclass maps. Thus all the global upper
spaces used below are unconditional.

Start with the norm-square subspace. The class of \(s\) is trivial at
every place above 2 and 73. At 2 one has
\(v_{\mathfrak p}(b/\theta)>2v_{\mathfrak p}(2)\) because \(k\ge3\);
at 73 the numerator \(b\) reduces to zero and \(\theta\) is a unit.
These conditions leave dimension eight. For each odd \(l\in D\), put
\(t_i=v_i(\theta)\) at the primes above \(l\). For integral \(b\), the
valuation vector of \(s\) modulo two is \((\max(t_i,0))_i\), up to
corrections supported where \(b\) and the integral \(\theta\) have the
same residue. Such a residue must be rational and must pass the two
center-square tests, the product-square test and the split square-root
test for the singleton cubic. For a pole at \(l\ne5\) the vector is
\((t_i)_i\) modulo two. A pole at 5 would give \(5+4/v^2\) of valuation one.
Taking the linear span of all these possibilities suffices. The
successive dimensions at \(3,5,23,41,2111,11801,144461\) are
\(6,5,4,2,1,0,0\). This proves the claim.
\cert{type1_descent/sextic.gp}

\textbf{The sextic elliptic cover.} Put \(d=-5\theta\) and
\[E:\quad Y^2=X(X-d\theta)(X+4d/5),\qquad X=db .\]
The three square conditions on \(b,5b+4,s\) give a point \(R\in E(F)\).
The split two-descent of \S\ref{arith:64} gives Selmer dimension five and hence
rank at most three. The five explicit basis triples are indexed by bits
of weights \(1,2,4,8,16\). The three nonidentity two-torsion points have
masks 27, 7 and 28. Their classes are distinct, so they span the image
of all torsion modulo doubling.
\cert{type1_descent/sextic_tangents.gp}

Use the completions \(F_1,F_2\) of ramification and residue degrees
\((2,1),(1,3)\) above 2. At each place scale the abscissa by \(c^2\)
and ordinate by \(c^3\), choosing the positive rational power of 2 for
which \(c\theta\) is a unit. Both models are integral. Write
\(m=105\cdot2^7\) and \(\Lambda_i=m\log E(F_i)\) on the scaled models.
The local abscissas in the data file give Kummer bases of dimensions
4 and 5. Hence their scaled logarithms generate the lattices.
Locally number these points starting at one. The logarithms of points
2,3 in the first list and 3,4,5 in the second form lattice bases.
The matrices of all the local basis logarithms modulo \(2\Lambda_i\) are
\[U_1=\begin{pmatrix}0&1&0&0\\0&0&1&0\end{pmatrix},
 \qquad U_2=(0_{3\times2}\mid I_3).\]
For the computation use \(z=-X/Y\) of the \(m\)-multiple. The invariant
differential is an integral even power series in \(z\), with constant
term one, so \(\log(z)-z\in z^3\mathcal O_{F_i}\). All nonzero
parameters computed are in \(2^7\mathcal O_{F_i}\). Multiplying either
inverse basis matrix by 256 makes it integral. Thus the indicated
bits, including those after one more division by 2, follow from the
precision in the certificate.
\cert{type1_descent/sextic_local_points.txt},
\cert{type1_descent/sextic_logs.m}

Let \(T=(0,0)\). The parameter of \(R-T\) is \(X/Y\) on the scaled
model and equals
\[\frac{\pm v}{2c\theta}
 \left((1-b/\theta)(1+5b/4)\right)^{-1/2}.\]
At both places the square-root correction is a unit congruent to
\(\pm1\) modulo 4, since \(k\ge3\) and \(v_2(\theta)\) is respectively
0 and 2. In the product of the two lattices, dividing
\(m\log(R-T)\) by \(2^{k-2}\) gives the primitive vector
\(u=(1,1,1,1,1)\) modulo two. Indeed the scaled leading term is an
odd rational unit times \(256/(c\theta)\), with either sign at each
place. The certificate checks its five bits. The cubic error in the
logarithm cannot change them since \(v_2(z)=k-1\) and the inverse
basis matrices have denominators dividing 256.

In particular \(R-T\) has infinite order. In the free group modulo
torsion write it as \(2^j Q\) with \(Q\) primitive at 2. The leading
vector of the product logarithm of \(Q\) must also be \(u\). On the
five-dimensional Selmer group the joint logarithm map modulo two has
rank two and misses \(u\). Its kernel consists of the four torsion
masks and their translates by 13. Thus the logarithm of \(Q\) is
even, and after a torsion translation its Kummer mask is 13.

\textbf{One more division.} Write \(e_i\) for the three roots of the elliptic
cubic and \(d_i\) for the entries of the mask-13 triple. Its projective
two-cover \(C\) has equations \(d_i z_i^2=X-e_i\) on the affine chart.
For each \(i\), omit \(z_i\) and take the tangent \(L_i\) to the resulting
conic at the supplied field point. Use the two squareclass functions
\(F_C=(L_2L_3,L_1L_3)\). Their values are unchanged modulo squares
by projective rescaling. They extend as squareclasses at a zero by
removing the even order of vanishing.

Here is the local comparison used in the calculation. If two points
on \(C\) differ by \(V\) in its Jacobian \(E\), the quotient of their
\(F_C\)-values is the ordinary Kummer class of \(V\). Indeed the
tangent section on \(C\) cuts out \(2D_i\), where \(D_i\) is a fiber
of the degree-two projection to the conic. The class \(D_j-D_k\) is
the two-torsion point \((e_i,0)\) for the complementary index \(i\):
the two deck involutions flip signs \(z_j,z_k\), and their product is
the deck translation of the two-cover fixing the square root of
\(X-e_i\). Thus \(L_j/L_k\), equal to \(L_jL_k\) modulo squares,
has divisor twice a representative of that two-torsion class.
Evaluation on a degree-zero divisor gives its usual Kummer coordinate.
On the target elliptic curve the images of the two points differ by
\(2V\).
\href{https://doi.org/10.1112/jlms/jds063}{Descent functions on coverings}

Take a local reference lift on \(C\) at each place in the calculation.
Write \(F_p\) for its pair of squareclasses and \(M_p\) for the complete
local Kummer image. Any global point on \(C\) gives a single global
pair \(g\) with \(g_p-F_p\in M_p\) in additive squareclass notation.
The support consists of the model's bad places, the primes of the
entries \(d_i\) and the primes of the coefficient ideals of the
tangents. Away from these places the primitive tangent has even
valuation on the smooth conic, as in \S\ref{arith:64}, and both functions have
even valuation. The verifier ranges over the entire resulting pair of
global squareclass spaces. It imposes the affine linear constraints
at each finite place in the support and at both real places.

At our two dyadic places the final abscissa in each data list fixes
the reference elliptic image \(P_i\). In the corresponding bases
\(m\log P_i/2\) has bits \((1,0)\) and \((1,1,1)\).
For every compatible global pair add to these bits the local images
under \(U_i\) of \(g_p-F_p\). The result is exactly the residue of
\(m\log Q/2\) for any global lift. In every solution of the linear
system the first coordinate at the ramified place plus the second
coordinate at the unramified cubic place is \textbf{one}. This functional
vanishes on zero and on \(u\). The logarithm of \(Q\) can therefore
neither be divisible by four nor have the required leading vector.
This contradiction rules out every nonzero \(v\) in the zero disk.
At \(v=0\) the reflected B chains coincide. Together with \S\ref{arith:64}
this proves the exclusion.
\cert{type1_descent/sextic_four.gp}

\subsection{Type 6. Descent on two Prym twists}\label{detail:type6}

\subsubsection{The three classes on the base cubic}\label{arith:66}

Give the cubic \(H_6(a,c)=0\) of \S\ref{app:type6} the origin
\(O=(1,0)\). It has rank two and torsion of order two. In its group
modulo 216 every valid type-6 configuration belongs to one of the
three classes represented by \((1,1),(9,1),(16,4)\).

\emph{Proof.} On a valid configuration the coefficient \(E_6\) in \S\ref{app:type6}
is nonzero. Reduction of the third moment equation modulo \(H_6\)
simplifies the B singleton square and the A singleton sum and product to
\[b=(5a-c+5)/4,\quad S=(-3a+15c+22)/4,\quad
 T=(8a^2-40ac-67a+80c^2+95c+149)/16.\]
Put \(c=v^2\). For \(v\ne0\) the substitution \(w=(a-4)/v\) yields
the smooth plane quartic with affine equation
\[vw^3-5v^2w^2+10v^3w-10v^4+3w^2-25vw+50v^2-40=0.\]
Its inverse is \(a=4+vw,\ c=v^2\). This gives a genus-three model for
the required square cover. In particular, at an odd prime \(p\nmid30\)
an integral candidate reducing to \((4,0)\) requires 40/3 to be a
square modulo \(p\). Indeed \(H_6(a,0)=(a-4)^2(a-1)\) and the
coefficient of \(c\) is \(-40\) at \((4,0)\). It follows that
\(w\) is a unit there and reduces to a root of \(3w^2-40\).

The minimal elliptic model is
\[y^2+xy+y=x^3+x^2-2913x+80031.\]
The rank bounds both equal two and the torsion subgroup has order two.
Put \(P=(4,0),Q=(4,1)\) on the plane cubic and let \(U=(1,1)-P\).
The point \(U\) has order two. Exact rational division verifies that
no nonzero binary combination of \(P,Q,U\) is divisible by two and
no nonzero combination of \(P,Q\) with coefficients in
\(\{0,1,2\}\) is divisible by three. Hence \(P,Q\) give bases modulo
9 and 27. Modulo 216 the classes are parametrized uniquely by
\(iP+jQ+tU\) with \(i,j\bmod216\) and \(t\bmod2\).

At each good prime below, enumerate the projective plane cubic.
Keep all points at infinity. For an affine point, test the three
center squares, the two singleton square roots of \(X^2-SX+T\)
with multiplicity, and the quartic condition at \((4,0)\).
Reduce the surviving points in the finite elliptic group modulo the
given multiple. Intersect with the reductions of the coefficient
classes, lifting the surviving classes between stages.

{
\begin{longtable}[]{@{}lll@{}}
\toprule\noalign{}
Multiple & Primes & Global classes left \\
\midrule\noalign{}
\endhead
\bottomrule\noalign{}
\endlastfoot
9 & 7, 11, 13, 17, 41, 61, 73, 307 & 4 of 81 \\
27 & 197, 271, 397, 647 & 3 of 729 \\
216 & 7, 13, 23, 31, 41, 43, 67, 73 & 3 of 93,312 \\
\end{longtable}
}

The final classes are \(P+U,P-Q+U,-P+2Q+U\). The group law gives the
representatives asserted in the statement. At those three points
\(c=1\) or 4, so the length-four B chain meets the centered A chain.
The next two arguments impose the covering conditions on all points in these cosets.
\cert{check_type6_classes.m}

\subsubsection{Two disks on a common Prym twist}\label{arith:67}

In the classes of \((1,1)\) and \((16,4)\) of
\S\ref{arith:66}, the only points of the base cubic with \(c\) a rational
square are the two class representatives, respectively. Both are
collisions. A valid type-6 configuration must lie in the class of
\((9,1)\).

\textbf{A common twist.} With origin \((4,0)\), the cubic has Weierstrass
coordinates \(t=(a-1)/c,\ Y=2A(t)c+B(t)\), where
\[A=t^3-5t^2+10t-10,\qquad B=-6t^2+5t+20.\]
Substitution gives \(Ac^2+Bc+9t-10=0\) and
\(Y^2=5t(32t^2-155t+192)\). Write \(K=\mathbf Q(r)\),
\(r^2=-551\), \(e=(155+r)/64\), and \(h=-B(e)/(2A(e))\).
For a representative \(R_0\), the two-torsion character
\((t(R)-e)/(t(R_0)-e)\) is a square in \(K\) throughout
\(R_0+216E(\mathbf Q)\). Changing the origin translates both points
and preserves their difference. With \(c=v^2\), the function
\(q_1=a-1-ec=c(t-e)\) therefore has constant squareclass.
For representatives 1, 2, 3 in the order of \S\ref{arith:66} we may use
\(\delta=-e,8-e,-e\), respectively. At \((16,4)\),
\((15-4e)/(-e)=((5+r)/16)^2\).

On the quartic of \S\ref{arith:66} in coordinates \((v,w,z)\), put
\[q_1=3z^2+vw-ev^2,\quad q_2=v^2-hz^2,\qquad
 q_3=\frac{q_2^2-H_6(4z^2+vw,v^2,z^2)/(A(e)v^2)}{q_1}.\]
Both divisions are exact. The quartic is \(q_1q_3=q_2^2\), and a
candidate lifts to \(R^2=\delta q_1, RS=\delta q_2,S^2=\delta q_3\).
If \(M_i\) represents \(q_i\), let
\(M(X)=M_1+2XM_2+X^2M_3\), \(g=-\det M(X)\). The sextic is
squarefree. Bruin's construction gives the genus-two Jacobian
\(J_\delta=\operatorname{Jac}(y^2=\delta g(X))\).
\cite[Theorem~5.1 and \S6]{BruinPrym}

For clarity, the divisors used below can be constructed without
choosing a point at infinity on the genus-two curve. Write
\(Z=(v,w,1), W=(1,w',0), L=R+XS\), differentiating with respect
to \(v\). Take the monic associate \(U\) of
\(\delta W M(X) W^t-(L')^2\), and reduce the second coordinate of
\((L'Z-LW)M(X)\) modulo \(U\) to obtain \(V\).
Then \(V^2=\delta g\pmod U\), by the tangent identity in the Type~13 construction of \S\ref{app:type13-details}
applied to \(\delta M(X)\) and division of the ordinate by \(\delta\).
The two support points minus a hyperelliptic fiber give a degree-zero
class, including when the support is nonsplit.

\textbf{Unconditional descent.} The quadratic and quartic factors of \(g\)
over \(K\) give absolute number fields of degrees four and eight.
They are isomorphic to the fields defined by
\[x^4-x^3+63x^2+313x+1544,\]\[
 x^8-3x^7+4x^6-328x^5+3937x^4-26339x^3+91290x^2-140562x+86400.\]
PARI certifies generation of the ideal class groups by the factor
bases with norms at most 5 and 83. Magma checks the field
isomorphisms, uses these bounds for the two fields and Minkowski
bounds for other fields, and obtains the rank upper bound 2 for
both twists by 2-descent.

\textbf{The disks at 197.} This prime splits in \(K\), and the reduced base
elliptic group has invariants \((2,108)\). Thus every point in each
global class specializes to its designated representative. Choose
\(v\equiv v_0\) with \((v_0,w_0)=(1,-3),(1,5),(2,6)\), respectively.
There is a single global sign of \(R\) for which the lift specializes
at \emph{both} places to the chosen
\(R_0^*=e,8-e,(7+5r)/32\), respectively. Indeed write the base point as
\(R_0+2D\) with \(D\in108E(\mathbf Q)\), so \(D\) reduces to zero.
The function
\[F(D)=\frac{t(R_0+2D)-e}{t(R_0)-e}\]
is a square of a rational function \(H(D)\) normalized by \(H(0)=1\).
The divisor of \(t-e\) is twice the divisor of the corresponding
two-torsion point minus the origin. Translation preserves this class
in \(\operatorname{Pic}^0\) and pullback by \([2]\) kills it. The
remaining constant is fixed by evaluation at zero. At each prime
above 197, \(F\) is a unit near the specialization of zero, and on
the smooth integral model the local ring is normal. Thus \(H\) and
its inverse are regular there and \(H(D)\equiv1\). Taking
\(R=R_0^*(v/v_0)H(D)\) gives the stated sign.

For the two places in the order \(r\equiv146,51\), use
\(\log([n_j]D)/197\) on each Jacobian block, where \(n_j\) is the
order of the reduced Jacobian. Use \(dX/(2y),X\,dX/(2y)\) as basis.
On the common twist the orders are 38196 and 42540. The collision
logarithms and the derivatives with respect to \((v-v_0)/197\)
reduce as follows.

{
\begin{longtable}[]{@{}lll@{}}
\toprule\noalign{}
Center & Collision logarithm & Derivative \\
\midrule\noalign{}
\endhead
\bottomrule\noalign{}
\endlastfoot
\((1,1)\) & \((123,77,4,92)\) & \((66,116,10,167)\) \\
\((16,4)\) & \((75,54,2,128)\) & \((148,119,66,137)\) \\
\end{longtable}
}

The verifier computes the logarithms by translating the multiple
to a split degree-two divisor and displacing its two abscissas modulo
\(197^2\). For the derivative write \(\dot U=dU/dv\).
Implicit differentiation at each root of \(U\) and interpolation give
\(-n_j[X](X^k\dot U/(2V)\bmod U)\) for \(k=0,1\).
All denominators in this chart are units throughout the disk. In
particular the two roots of \(U\) are distinct in the residue
algebra and their ordinates are units.

The determinants of the first three columns of the two logarithms
followed by the respective derivative are 7 and 100 modulo 197.
Hence the two global classes span the rational Mordell--Weil group
up to torsion, and on each disk an integral linear functional
annihilates their exact logarithms and is a unit on the tangent.
Its integral vanishes at the collision and at any rational lift.
The pulled-back differential has integral power-series coefficients
in \(v-v_0\). For \(m\ge2\) and \(k=v_{197}(v-v_0)\ge1\),
\(mk-v_{197}(m)>k\), so the linear term of the integral dominates
every higher term. Consequently \(v=v_0\); the quartic has a unique
\(w\) in its disk since its \(w\)-derivative at these centers is a
unit. This proves the proposition.

For the twist \(\delta=8-e\) treated in \S\ref{arith:68}, the same calculation gives
orders 39780 and 42540, collision logarithm \((184,142,27,148)\),
and tangent \((69,56,174,18)\). The logarithm proves that the
collision divisor has infinite order. The certified rank interval
for this twist is \([1,2]\).

\cert{type6_descent/class_bounds.gp},
\cert{type6_descent/prym.m}

\subsubsection{The remaining dyadic disk}\label{arith:68}

Type 6 has no valid distinct-entry chain at any rational
height. In the class of \((9,1)\) the full necessary chain conditions force
\((a,c)=(9,1)\).

\textbf{The base disk.} Use the generators and origin of \S\ref{arith:66}.
None of the 48 nonzero combinations of the generators modulo 7 is
divisible by 7. Of the 49 lifts of the class of \((9,1)\) from 216 to
1512, the reduction tests of \S\ref{arith:66} leave seven at 67 and one after 79.
It is the class of \((9,1)\) itself. On the minimal elliptic model,
the multiples by 1512 of both generators lie in the fifth formal
subgroup at 2, as checked with the parameter \(-x/y\).
The finite quotient by this subgroup has order divisible only by 2
and 7: the reduction is split multiplicative of type \(I_{14}\).
The saturation checks at 2 and 7 therefore give the same inclusion
for the full rational Mordell--Weil group.
\cert{check_type6_classes.m}

In the Weierstrass coordinates of \S\ref{arith:67} the minimal coordinates satisfy
\(x=40t-65,\ 2y+x+1=40Y\). Hence the invariant differential is
\[\frac{dx}{2y+x+1}=\frac{dt}{Y}=-\frac{dc}{(H_6)_a}.\]
The last identity follows by differentiating \(A c^2+B c+9t-10=0\).
At \((9,1)\), the partial derivatives of \(H_6\) are 40 and \(-120\).
Putting \(c=1+64z,\ a=9+64u\) gives, after division by \(64\cdot40\),
an equation \(u-3z+8G(z,u)=0\) with integral higher terms.
It has a unique integral analytic branch through the origin and
\(40/(H_6)_a\) has integral series with constant 1 there. For
\(c=1+256s\), the elliptic logarithm divided by \(-32/5\) is therefore
\(s\) plus a strictly contracting integral series: each coefficient
of degree \(j+1\ge2\) has valuation at least \(2j-v_2(j+1)\).
This parametrizes the whole translate of the fifth formal subgroup,
by the elliptic formal exponential and the inverse function theorem.
Indeed the parametrization and the formal exponential of the integral
agree near zero and then on the disk by analytic continuation.
Consequently \(c\equiv1, a\equiv9\pmod{256}\). Choose the rational
sign of \(v=\sqrt c\) with \(v\equiv1\pmod{128}\); then
\(w=(a-4)/v\equiv5\pmod{128}\).

\textbf{The possible second global direction.} Take \(d=8-e\) and use the
reciprocal model \(y^2=f(X)=X^6 d g(1/X)\) of \S\ref{arith:67}. Write \(J\)
for its Jacobian, \(D\) for its collision divisor, \(T\) for its
nonzero rational two-torsion point, and \(F\) for a hyperelliptic
fiber. Write \(\theta\) for the class of \(X\) in the sextic algebra.
For a divisor of odd degree one can evaluate \(X-\theta\) modulo
squares and base-field scalars. The norm of this class equals the
leading coefficient of \(f\) modulo squares. Differences of two such
classes belong to the usual even descent space.
\href{https://math.mit.edu/~poonen/papers/descent.pdf}{Poonen and Schaefer, §§5, 11--13}

Here is the affine descent calculation. Put \(S\) equal to the ten
places over \(2,3,5,7,31\). The monic sextic is integral, its
discriminant a unit, and the leading coefficient of \(f\) a unit
outside \(S\). At such a place an integral abscissa specializing
to a root in the étale algebra has even order in that component
because the branch point is simple and \(y^2=f(X)\). At a pole all
components have the same valuation. Applying this to closed points
over their residue fields and taking norms gives diagonal valuation
parity for arbitrary divisors. The class group of \(K\) is cyclic
of order 26 and a dyadic prime generates it. Thus a base-field
scalar clears these diagonal parities outside \(S\). The two factor
fields of \S\ref{arith:67} give a space of dimension 40 of squareclasses with
even valuations outside \(S\), before taking the quotient by scalars.

At each listed place the verifier takes an odd local point as origin.
It adjoins the scalar classes and the differences of odd points, the
class of \(D\), and the class of \(T\). At one dyadic place an additional
degree-two divisor has support \(X^2+5/2\); extracting a square in
\(\mathbf Q_2(\sqrt{-10})\) verifies its ordinate. The local spaces
have dimensions 8 and 6 at 2, 4 and 4 at 3, 3 and 3 at 5, 2 and 4
at 7, and 4 and 2 at 31, in the order of the certificate.

These spaces are complete. With \(n\) irreducible factors at an odd
place the dimension including scalars is \(n\), and at the two dyadic
places it is \(n+3\). In fact, when there is an odd factor the scalar
map is injective, rational two-torsion has dimension \(n-2\), and
the local descent is injective modulo doubling. At the two remaining
odd places the degrees are \((2,4)\). There the two-torsion dimension
is one. If both fields contain the same quadratic extension then the
scalar map loses one dimension and the roots admit a conjugate
three-plus-three partition. Otherwise the scalar map is injective
and there is no such partition. The kernel criterion of
Poonen--Schaefer gives total dimension two in either case.
All ten places have local points; at all other finite places good
reduction gives an odd-degree divisor by passing to a sufficiently
large odd unramified extension. Thus any global rational divisor
class of degree one is represented by a divisor, by the local
Brauer obstructions and the Brauer--Hasse--Noether theorem.

The norm and local affine constraints, homogenized with one parity
coordinate, have a 14-dimensional nullspace. There is an odd vector,
so the even subspace has dimension 13. Scalar classes account for 11:
the quartic relative factor field has no quadratic intermediate field
(the absolute octic field has just one quadratic and no quartic
subfields). The images of \(D,T\) are independent, as the local
coordinates below show. They therefore span the even space modulo
scalars.

At the two dyadic places in the order \(r\equiv155,101\pmod{256}\),
choose a local Weierstrass point \(W_0\) as in the certificate. The
free local coordinates use the differences from \(W_0\) of the points
with abscissas 352 and 4704 at the first place. At the second use
the point at 9 minus \(W_0\), and the quadratic divisor at
\(X^2+5/2\) minus a fiber. Together with three and one independent
Weierstrass differences, respectively, these give bases modulo
doubling. In particular the two-primary local torsion has exponent
two. Sign choices of these free generators give the same coordinates
modulo two. For local torsion coordinates followed by free coordinates,
\(D\) has rows \((0,1,1;0,0)\), \((0;0,0)\).
The rows of \(T\) are \((0,1,1;0,0)\), \((1;0,0)\).
Every possible odd class \(W\) has
\[\log(W-W_0)\equiv(1,1,0,1)\pmod {2\Lambda},\]
where \(\Lambda\) is the product of the free local logarithm lattices.
Indeed the verifier takes any odd solution of the affine system,
subtracts the local Weierstrass images and reads the two free
coordinates at each place. Adding either even generator preserves them.
\cert{type6_descent/beta_selmer.m}

If the rank is two, \(J(K)/2J(K)\) has dimension three. Its even
descent image has dimension two, and Poonen--Schaefer's Theorem 11.3
writes the extra kernel generator as \(B=2W-F\) for a degree-one
class, modulo doubling. Choose that representative. Since \(F=2W_0\)
locally, \(\log(B)/2\) gives the displayed row. The classes
\(D,T,B\) span modulo doubling and \(D,B\) span up to finite index
and torsion. In rank one, \(D\) alone spans up to finite index and
torsion.

\textbf{The disk calculation.} In the same free local coordinates the
collision logarithm divided by two and the derivative at \(v=1\)
multiplied by eight have respective reductions
\[(1,1,1,0),\qquad (1,0,0,0).\]
The first two proposed global rows \((1,1,1,0),(1,1,0,1)\) are
independent modulo two. Thus even after saturation the reduction of
the global logarithm space is contained in their span, also in the
rank-one case.

For details of precision, the verifier multiplies each divisor by
\(2^{18}\cdot315\) and translates it to two split points with ordinates
\(y_j\). If their abscissa displacements are \(h_j\), integrate in
the variable \(X=x_j+h_j z\). The degree-\(k\) coefficients of
\(f(x_j+h_j z)/y_j^2-1\) have valuations at least \(6k\) for
\(1\le k\le6\). The coefficients of the inverse square root of \(f(x_j+h_j z)/y_j^2\)
with constant 1 have valuations at least \(4k\). After division by
the integration index and the multiplier, the discarded terms of
degrees eight and higher give error in \(2^8\mathbf Z_2\) in both
logarithms. The inverse basis matrices have entry valuations at
least \(-1\), so this precision determines the displayed reductions.

On the reciprocal curve the two pulled-back differentials are
\[\gamma(R,-S)\frac{dv}{F_w},\qquad
 \gamma=(58773-2161r)/128000,\]
where \(F(v,w)=0\) is the quartic and \(R^2=d q_1, RS=d q_2\).
Indeed the canonical complete intersection of three quadrics has
canonical bundle \(\mathcal O(1)\). Its two anti-invariant
differentials have basis \(R\,dv/F_w,S\,dv/F_w\), obtained by
multiplying the plane-quartic differential by the extra linear
coordinates. Changing the sheet negates the Mumford image, so the
two pullbacks lie in this space. Their constant and linear jets at
the collision determine the coefficients: the basis jet determinant
is nonzero and the verifier evaluates both using the monic tangent
support.

Here is a uniform bound for integration throughout the required disk.
Put \(v=1+32z, w=5+32u\). The scaled quartic has unit \(u\)-derivative
and gives an integral series \(u=z+4z^2\mathbf Z_2[[z]]\).
Also \(40/F_w\) is integral with constant one. Write \(L_j\) for
the local basis matrix. The vectors \(\gamma d(1,0)L_j^{-1}\),
\(\gamma(1-h)(0,1)L_j^{-1}\), \(\gamma(0,1)L_j^{-1}\) have entry
valuations respectively \((-5,-2),(-5,-2),(-5,-2)\) at the first
place and \((1,2),(4,3),(0,-1)\) at the second. On the positive
sheet \(R/d=1\) at the origin. At the first place
\(v_2(e)=v_2(d)=-5\) and \(q_1/d-1\in64z\mathbf Z_2[[z]]\).
Thus \(R/d-1\) and \(S-(1-h)\) belong to \(32z\mathbf Z_2[[z]]\).
At the second place \(v_2(e)=6,v_2(d)=3,v_2(1-h)=4\);
\(q_1/d-1\in8z\mathbf Z_2[[z]]\) and \(S=(v^2-h)/(R/d)\).
Together with integrality of the constant row, these bounds make
eight times the differential coefficient integral in the basis on
both disks. These parametrizations on the open disk \(|z|_2<1\)
have analytic coordinates and the parameter is recovered from \(v\).
On such a disk the abelian logarithm displacement equals the
termwise primitive, by functoriality from the Abel--Jacobi map.
\href{https://www.mathe2.uni-bayreuth.de/stoll/papers/uniform-bounds-2015-11-25.pdf}{Stoll, §7, integration on an open disk}

For \(k=v_2(v-1)\ge7\) we have \(v_2(z)\ge2\). After
integration and division by \((v-1)/8\), each nonconstant term vanishes
modulo two since \(2j-v_2(j+1)\ge1\) for \(j\ge1\). The displacement
therefore has primitive row \((1,0,0,0)\), outside the global span.

Finally the two local signs of \(R\) need not be prescribed. Negate
the global divisor if necessary to obtain the positive sign at the
first place. If the sign is negative at the second, subtraction of
the global collision logarithm adds the second block of \(-2\log D\).
After division by four this has row \((0,0,1,0)\), also outside the
global span, while the varying displacement is divisible by 16.
If the signs agree the preceding tangent test rules out every
\(v\ne1\). At \(v=1\) the unique local quartic branch has \(w=5\),
so \(c=1,a=9\) and the chains collide. \S\ref{arith:67} handled the
other two classes.
\cert{type6_descent/beta_logs.m}

\subsection{Type 7. Descent from a rational branch point}\label{detail:type7}

\subsubsection{Five classes on the cubic}\label{arith:69}

With origin \((a,c)=(1,1/4)\), the type-7 cubic
of \S\ref{app:type7} has rank two and trivial torsion. A valid chain must be in
one of five classes modulo 216, represented by
\[(1,1/4),\quad(9,9/4),\quad(4,49/4),\quad(4,1/4),\quad(1,25/4).\]

\emph{Proof.} On \(H_7=0\) the linear subresultant of the two remaining
moment equations gives
\[b=(10a-4c+15)/6,\qquad S=(16a+20c+81)/12,\]\[
 T=-(128a^2-640ac+48a+80c^2-1380c-1701)/144 .\]
The resultant in \(a\) of its leading coefficient and \(H_7\) is
proportional to
\[(4c-9)^4(4c-1)^6(4c+11)
 (144c^2+152c+529)(144c^2+440c+1225).\]
The first two factors give collisions, the linear factor at negative
\(c\) is impossible for a square, and the two quadratics have negative
discriminants. The formulas therefore apply to every valid lift.
\cert{check_all_shift_chain_types.py}

The minimal elliptic model is \(y^2=x^3+x^2-13533x-323937\).
The two rank bounds agree. With
\(P=(-7/2,1/4),Q=(9,9/4)\) on the plane cubic, exact division at
2 and 3 proves that these points generate the quotient modulo 216.
At good primes apply the projective reduction test of \S\ref{arith:66} using
the three center squares and the two square roots of \(X^2-SX+T\).
The successive moduli 9, 27 and 216 leave seven, five and five
classes. At 216 they are \(0,Q,-P+Q,2P-2Q,3P-2Q\), respectively.
All five representatives are collisions. For \((4,49/4)\) one has
\(b=1\), so the two length-two chains share \(3/2\). For
\((1,25/4)\) one has \(b=0\), making the reflected B chains
coincide. The other representatives have \(c=1/4\) or \(9/4\).
Sections~\ref{arith:70} and~\ref{arith:71} determine the points satisfying the necessary chain conditions in all five classes.

The square cover \(b=v^2\) has a rational branch point at
\((1,25/4)\). Restriction to the zero line
\(a=(4c-15)/10\) gives
\[H_7((4c-15)/10,c)=-27(4c-25)(144c^2+440c+1225)/125.\]
The other two finite branch points are defined over
\(\mathbf Q(\sqrt{-5})\), and the three simple poles lie at infinity.
Riemann--Hurwitz gives genus four. The descent in \S\ref{arith:70} starts from the rational branch point.
\cert{check_type7_classes.m}

\subsubsection{Descent at a rational branch point}\label{arith:70}

In four of the five classes in \S\ref{arith:69}, the
square and splitting conditions for a type-7 chain force the collision
representative. The remaining class is that of \((9,9/4)\).

Write \(u=a-1\). The line \(a=1\) cuts the cubic in \((1,25/4)\) and
twice the origin \((1,1/4)\). On the genus-four cover of \S\ref{arith:69}, put
\(t=b/u\). Then \(c=(25+(10-6t)u)/4\), and the cubic divided by
\(72u\) reads
\[A u^2+B u+C=0,\qquad
 A=45t^3-145t^2+135t-27,\quad
 B=-8(45t^2-80t+21),\quad C=144(5t-2).\]
Putting \(y=(2Au+B)/8\) gives
\[y^2=5F(t),\qquad F(t)=27t^3-79t^2+57t-9.\]
The inverse recovers \(u\) generically, so this is an isomorphism
of the smooth projective base curves. Its only pole of \(t\) is the
origin. At the other point on \(a=1\), \(t=2/5\). The Weierstrass
coordinates \((135t,135y)\) have coefficients
\([0,-395,0,38475,-820125]\).

For any other lift, \(t=d z^2\), where \(d\) is a positive squarefree
divisor of 30. Indeed \(5F(t)<0\) for \(t\le0\). At a prime \(p>5\)
a pole of \(t\) has even order by the cubic equation. At a zero of
its reduction the leading and constant coefficients of the quadratic
in \(u\) are units, so \(u\) is a unit and \(tu=b\) is square.
To select \(d\), reduce the coset \(R_0+216E(\mathbf Q)\) at
\(11,17,19,41\). At a pole the unit squareclass is 15 by the leading
coefficient 135. At a zero use \(27u^2+168u+288=0\) and the
squareclass of \(u\). The five cosets give the following results in the
order of \S\ref{arith:69}. For the additional disk tests in the last
column impose all five square conditions of that proposition on every
finite fiber. A vanishing leading coefficient is retained.

{
\begin{longtable}[]{@{}lrrr@{}}
\toprule\noalign{}
\((a,c)\) & \(t_0\) & \(d\) & Prime forcing \(t\equiv t_0\) \\
\midrule\noalign{}
\endhead
\bottomrule\noalign{}
\endlastfoot
\((1,1/4)\) & \(\infty\) & 15 & \\
\((9,9/4)\) & 2 & 2 & 17 \\
\((4,49/4)\) & \(1/3\) & 3 & 17, 83 \\
\((4,1/4)\) & 3 & 3 & 61 \\
\((1,25/4)\) & \(2/5\) & 10 & 389 \\
\end{longtable}
}

\textbf{Cubic-field elliptic covers.} Let \(e\) be a root of \(F\) and write
\(F(t)=27(t-e)(t^2+g_1t+g_0)\). In \(K=\mathbf Q(e)\) the
two-descent character on the monic base model is \(135(t-e)\).
Every point in a fixed coset has the character of its representative
because 216 is even. Thus \(t-e\) has class \(\delta=t_0-e\) for a
finite representative and class 135 at the origin. Put \(k=135d\delta\).
Multiplying the square conditions gives the elliptic lift
\[X=kt,\qquad Y^2=X(X^2+kg_1X+k^2g_0).\]
In the five rows its ranks over \(K\) are respectively \(0,3,3,2,2\).
The first curve has torsion of order two, with abscissas zero and
infinity. Since \(5F(0)=-45\), only the origin remains.

The primes 61 and 389 are inert in \(K\). For the corresponding
rank-two curves take the two global points in the logarithm certificate.
At the anchor the ordinate is a unit. If \(N\) is the local reduction
group order, the coordinates of \(\log([N]P_i)/p\) modulo \(p\)
are obtained from the formal parameter \(-X/Y\) modulo \(p^2\).
The third row is the tangent \(k/(2Y_0)\). In the power basis of
\(h^3-8h-2=0\), the determinants, up to sign, are 31 at 61 and
216 at 389. Both reduction orders are prime to \(p\).
The first two logarithms thus span the global logarithm space over
\(\mathbf Q_p\), and the tangent lies outside it. There is a single
place, so negating the lift chooses the branch through the global
anchor. A linear form annihilating the first two rows and taking a
unit value on the third gives a function with a unique zero on the
whole disk. Indeed the invariant differential has integral coefficients
in \(t-t_0\), since its denominator is a unit. After integration a term
of degree \(j\ge2\) has valuation at least
\(jv_p(t-t_0)-v_p(j)>v_p(t-t_0)\).

\textbf{A second descent for the rank-three rows.} On
\(C_d:y^2=5F(dz^2)\), the two-cover calculation gives four possible
classes for \(d=2\) and eight for \(d=3\). For the rows with
\(t_0=2,1/3\), write \(z_0=1,1/3\), respectively. At 17, the
only admissible classes in the sextic algebra are represented by
\(z_0-\theta\) and \(-z_0-\theta\), where \(F(d\theta^2)=0\).
For a component of even degree test that the norm of the ratio is
square. For the odd-degree components the square characters must agree,
since the class is defined modulo a common rational scalar. The
certificate tests every class returned by two-cover descent. Changing
the global sign of \(z\) exchanges the two classes, so take
\(z_0-\theta\). This is the two-cover map of Bruin and Stoll.
\href{https://arxiv.org/abs/0803.2052}{Two-cover descent on hyperelliptic curves}

Put \(n^2=d/3\), \(s=nz\), and let \(r\) satisfy
\(9r^3-5r^2-3r+1=0\). Use \(n=1\) for \(d=3\). The factorization
\[F(3s^2)=9(9s^3-5s^2-3s+1)P(s),\qquad P(s)=9s^3+5s^2-3s-1\]
yields an elliptic quartic over \(L=\mathbf Q(r,n)\). In the chosen
two-cover class the product of the four factors corresponding to
\((s-r)P(s)\) has the squareclass of its value at \(s_0=nz_0\).
Indeed the four factors form a Galois-stable subset over \(L\) and
the common rational scalar occurs to the fourth power. With
\(k=(s_0-r)P(s_0)\), \(m=kP(r)\) the lift and its Weierstrass
coordinates are
\[q^2=k(s-r)P(s),\qquad X=m/(s-r),\quad Y=mq/(s-r)^2,\]
\[Y^2=X^3+kP'(r)X^2+km(27r+5)X+9km^2 .\]
For \(d=3\) its rank is one. At 83 the selected class forces
\(z\equiv1/3\): the component characters for the other sign disagree.
All three places have degree one. Two lift signs agree, and the minors
of the anchor logarithm and its slope for the three pairs of places
are nonzero (in cyclic order, 67, 70, 22). The same integral estimate
as above bounds each disk by the single anchor. This proves \(t=1/3\).

Finally the pairs of quadratic roots \(u\) at \(t=1/3,3,2/5\)
are \(\{3,-9/2\},\{3,13/6\},\{0,-760/167\}\). In each the second
value makes \(1+u\) a nonsquare. This proves the proposition.
\cert{type7_descent/classes.m},
\cert{type7_descent/covers.m},
\cert{type7_descent/odd_logs.gp}

\subsubsection{The final sextic cover}\label{arith:71}

Type 7 admits no valid distinct-entry chain at any
rational height. In its remaining class the necessary chain conditions
force \((a,c)=(9,9/4)\).

We use the local divisibility criterion of Selmer-group Chabauty.
Here the two-isogeny and a lift of one abscissa class determine the
local information it requires.
\href{https://arxiv.org/abs/1506.04286}{Stoll, \emph{Chabauty without the Mordell--Weil group}}

\textbf{Global descent on the sextic cover.} Use the last quartic of \S\ref{arith:70}
with \(n^2=2/3\) and \(s_0=n\). Shift its abscissa by
\(j=-m/(2r)\), writing \(x=X-j\). The point over \(s=-r\) is now
\(T=(0,0)\), and the curve and its two-isogenous curve have models
\[E:y^2=x(x^2+A x+B),\qquad E':y'^2=x'(x'^2-2A x'+D),
 \qquad D=A^2-4B .\]
Here \(A=kP'(r)+3j\), \(B=3j^2+2jkP'(r)+km(27r+5)\).
Write \(s_i\) for the four elements listed in the certificate, and
\(s_* = s_2\). The global image under the abscissa character
\(\alpha\) for \(E\) lies in \(\langle s_1,s_2,s_3,s_4\rangle\).
The isogeny Selmer space for \(E'\) is exactly \(\langle D\rangle\).
In particular a point of \(E(L)\) of trivial abscissa class belongs to
\(2E(L)\): it lifts along the dual isogeny to \(E'(L)\), where the
two possible classes are represented by the identity and the kernel
point of the dual isogeny.

Here are details of the finite computations. The field has model
\(h^6-96h^4+144h^2-54=0\), class number one and one place above 2,
of ramification index six. PARI certifies its class group and units.
On each curve the abscissa class is unramified outside the primes
dividing \(2BD\) (including denominator ideals). Form the squareclasses
of units with these primes inverted and intersect with the local
images at those primes and all real places. For the two local images
at a finite prime \(v\), the dimensions sum to
\(2+[L_v:\mathbf Q_2]\) above 2 and to 2 elsewhere. This follows by
composing the two isogenies and using the order of the local group
modulo doubling. The certificate finds local points until the paired
images attain this dimension. At a real place the negative abscissas
occur precisely when the model has a negative root and a nonempty
interval of negative abscissas with positive cubic value. The two
resulting dimensions are 4 and 1, giving rank at most three for the sextic cover. The certificate identifies the basis
by testing squares in the number field. Locally at 2 the kernel on
the four-dimensional space is exactly \(\langle s_*\rangle\).

\textbf{Separating the two classes in this kernel.} Let
\(e=(-A+\sqrt D)/2\). A global point of nonzero abscissa in the
class \(s_*\) can be written \(x=s_*u^2,\ y=s_*uv\), with
\(v^2=s_*u^4+A u^2+B/s_*\). Choose the global point \((U,V,1)\)
on \(s_*U^2+AUV+(B/s_*)V^2=1\) specified in the certificate.
Put \(\gamma=s_*U-eV\). The norm of \((x-e)/\gamma\) is \(v^2\).
Hilbert 90, or the identity \((w+v)^2/w=\operatorname{Tr}(w)+2v\)
when \(w\bar w=v^2\), shows that
\[x-e=\lambda\gamma\cdot\text{square in }L(\sqrt D),
 \qquad \lambda\in L^\times .\]
Where the tangent value is nonzero, its squareclass in the quadratic
extension is that of \(\lambda\):
\[\mathcal L(u,v)=(2s_*U+AV)u^2+AU+2BV/s_*-2v.\]
In a fixed local abscissa class the possible \(\lambda\)'s form one
coset of the local abscissa image of \(E'\). Indeed the second
two-descent character is \(x-e\). On the image of \(E'\) under the
dual isogeny it has class \(x'\), by the identity
\(x(\widehat\phi(P'))-e=(x'-\sqrt D)^2/(4x')\).
The kernel of squareclass extension is generated by \(D\), itself an
abscissa class of \(E'\). These identities also describe the characters
at the exceptional points by continuation of the usual descent map.

The certificate exhibits one element \(\kappa\) in every required local
coset. It evaluates \(\mathcal L\) on a local quartic point and tests
membership of \(\kappa\mathcal L\) in the local image of \(E'\).
At untested odd places the curves have good reduction, the coefficients
of the model are integral, and \(s_*,D,\kappa\) and \(N(\gamma)=s_*\) are units.
The conic point is integral there: the minimum valuation of its tangent
coefficients is zero and the conic matrix is invertible. Thus the
second character and \(\gamma\) are unramified and the local image on
\(E'\) consists of both unramified classes. The finite test includes
the support of the fractional ideal generated by the tangent coefficients.
The real signs are tested by exact number-field signs. Consequently the
global ratio \(\lambda/\kappa\) is in the dual Selmer space
\(\langle D\rangle\). Every global point in this class satisfies
\(x-e=\kappa\gamma\cdot\mathrm{square}\) in the quadratic extension.

Write \(M=L_v\) at the unique dyadic place, with valuation \(v(2)=6\).
Here \(s_*\) is square but \(\kappa\gamma\) is nonsquare in
\(M(\sqrt D)\). Both \(\kappa(2s_*U+AV\pm2\sqrt{s_*})\) are
nonsquares in \(M\), the trace-and-norm criterion for this assertion.
Hence \(E(L)\cap2E(M)=2E(L)\). The local two-torsion consists of
\(0,T\) because \(D\) is nonsquare. Since \(B=\alpha(T)\) is
nonsquare, \(T\) has no local half. On quotienting by \(\langle
T\rangle\), the local group has no two-torsion and the global subgroup
still has the same divisibility property. It applies repeatedly.

\textbf{The whole dyadic disk.} On the minimal rational base model the
component group at 2 has order 3 and its nonsingular reduction
has order 2 (additive reduction). The third formal subgroup has index
24, so contains \(216E_{\rm base}(\mathbf Q_2)\).
At the representative \((138,\pm675)\) the first abscissa derivative
in the formal group parameter is \(\pm1350\), and higher coefficients
are integral. As the minimal abscissa is \(135t-132\), this yields
\(t\equiv2\pmod {16}\), hence \(z\equiv\pm1\pmod4\).
For \(z\equiv-1\) the unit \((nz-r)P(nz)/k\) has the nonsquare
class of its value at \(-1\): changing \(z\) by a multiple of four
changes this unit by an element of valuation at least 15. Thus
\(z=1+4\xi\) with \(\xi\in\mathbf Z_2\).

Choose the sign of the lift so that \(q\) is on the branch through
\(q(n)=k\). Write \(H(\xi)=P(1+4\xi)-P(1)\) for the difference of the
elliptic lifts, using \(P(z)\) here for the point. For \(\xi\ne0\)
its formal parameter has valuation \(9+v(\xi)\). Indeed the
abscissa difference before subtraction has valuation \(25+v(\xi)\)
and the ordinate sum has valuation 16 by the quartic formulas. The
chord slope therefore has valuation \(-9-v(\xi)\).
On the seventh formal subgroup the logarithm and exponential give
inverse isometries with \(\mathfrak p^7\) for the displayed integral
model, since \(7>v(2)\). In particular \(\mathfrak p^{19}\) maps into
four times this subgroup.

Set \(g(X)=(n+4nX-r)P(n+4nX)/k\). Its constant is 1 and the higher
coefficients have valuations \(17,32,47,60\). Pulling back
\(dx/(2y)\) gives
\[\log H(\xi)=-\frac{2n}{k}\int_0^\xi g(X)^{-1/2}\,dX .\]
Its linear and quadratic coefficients have valuations 9 and 14.
All coefficients of higher degree have valuation at least 19. For
degrees 3 through 12 the valuations are
\(25,24,35,40,51,44,55,60,71,70\). For degree \(i+1\), \(i\ge12\),
the binomial expansion gives the lower bound \(9+3i-6v_2(i+1)\):
a degree-\(i\) product contributes at least \(15i\), and its binomial
denominator costs at most \(12i\). This bound is at least 19 and tends
to infinity.

Write \(\xi=2^j u\) with \(u\) an odd dyadic unit. Modulo
\(\mathfrak p^{19}\), \(2^{-j}\log H(\xi)\) equals
\(u\log H(1)\) for \(j=0\) and \(u\log H(2)/2\) for \(j\ge1\).
For odd \(u\) the quadratic difference has the extra factor \(u^2-u\);
for \(j\ge1\) the normalized quadratic term already has valuation at
least 20. Modulo four times the local group, \(u\) can be replaced by
an odd integer in its residue class modulo four.

It remains to check two points. Modulo \(T\), \(H(1)\) admits one
successive local halving and \(H(2)\) admits two. In both cases choose
the final representative \(R\) modulo \(T\) having square abscissa.
Neither \(R\) nor \(R+T\) has a local half. Moreover both
\(\kappa\mathcal L(u,v)\) and \(D\kappa\mathcal L(u,v)\) for this
representative are nonsquare, using \(u^2=x(R)/s_*\). Thus its class
modulo doubling cannot come from a global point. Only the global
abscissa classes 1 and \(s_*\) could give its local abscissa;
the first would be globally divisible, and the second fails the
tangent test. Translation by \(T\) preserves the image of the global
subgroup. The certificate checks halving by extracting square roots
of \(x\) and \(2x+A-2y/\sqrt x\), trying both signs, and checks the
precision of each square test. By the logarithm congruence every
\(H(\xi)\) with \(\xi\ne0\) has this same obstruction after division
modulo \(T\): division is unique there and multiplication by an odd
integer acts trivially modulo two. Global divisibility proved above
would require the primitive class to be global. Hence \(\xi=0\).

This forces \(t=2\). The quadratic gives \(u=a-1=8\) or \(144/23\).
The latter makes \(a=167/23\) a nonsquare. The former gives
\((a,c)=(9,9/4)\), a collision by \S\ref{app:type7}. \S\ref{arith:70} treats the
other four classes.
\cert{type7_descent/sextic.gp},
\cert{type7_descent/dyadic.m}

\subsection{Type 8. Reciprocal elliptic covers}\label{detail:type8}

\subsubsection{Product curves and three possible squareclasses}\label{arith:72}

Let \(b=v^2\) be the squared B center in a valid
type-8 chain. Every odd prime in the reduced denominator of \(v\) is
1 modulo 20, and \(v_2(v)\ge-1\). Put
\[G(b)=5(256b^4-1152b^3+1824b^2-388b+405).\]
The curve \(Y^2=G(b)\) has rational rank five, and the Jacobian of
\(C:Y^2=bG(b)\) has rational rank four. In the two-cover descent on
\(C\), each valid chain belongs to the class represented by
\(b_0-\theta\) for one of \(b_0=1/4,9/4,25/4\), where
\(\theta G(\theta)=0\) in the associated étale algebra.

\emph{Proof.} For the B polynomial put
\(B_v(X)=\prod_{j=-2}^2(X-(v+j)^2)\). The singleton polynomial of
\S\ref{app:type8} satisfies
\[(X-1/4)K_b(X)=B_v(X)-B_v(1/4),\qquad 256K_b(0)=G(b).\]
All four roots of \(K_b\) are rational squares. Their product gives
the quartic model; multiplication by the square \(b\) gives the
genus-two model. A vanishing B center already repeats vertices.

Suppose \(p\) is an odd denominator prime, or \(p=2\) and
\(v_2(v)<-1\). After substituting \(X=bU\), the monic quartic
\(K_b(bU)/b^4\) is integral locally. Its four roots are integral
squares. The five normalized B vertices have square tending to 1
modulo \(p\), and the normalized central A square is zero modulo
\(p\). The reduction is therefore \(((U-1)^5+1)/U\). At 2 this
polynomial has no linear root. At 5 it is \(U^4\); all four roots
would be divisible by 5, while the constant term \(G(b)/(256b^4)\)
has valuation exactly one. At every other odd prime the four roots
are \(1-\zeta^i\), for \(1\le i\le4\) and a primitive fifth root of
unity \(\zeta\). Splitting requires \(5\mid p-1\). The ratio
\((1-\zeta^{-1})/(1-\zeta)=-\zeta^{-1}\) is square, and \(\zeta\)
is square because its order is odd. Thus \(-1\) is square as well,
giving \(p\equiv1\pmod {20}\).

The quartic with the point over \(b=1/4\) has minimal model
\(y^2=x^3-105150x+7613125\), with both rank bounds equal to five.
For the genus-two Jacobian the rank upper bound is four and the
rational torsion has order two. The points over
\(1/4,3/4,9/4,4,25/4\), based at the rational point at infinity,
give five independent images under the two-descent map. This proves
the matching rank lower bound.

Two-cover descent gives seven possible classes on \(C\). At
\(13,17,19,43\) the denominator result ensures an affine reduction.
Enumerate the square residues \(b\) for which \(K_b\) splits into
squares. In each unramified factor of the descent algebra compare
the component norm of \(b-\theta\) with that of the representative,
allowing a common scalar as in \S\ref{arith:70}. A factor where \(b-\theta\)
vanishes is omitted from the test. The successive numbers of classes
are 6, 5, 4 and 3, and the last three are precisely the stated ones.
All three representatives are collisions by their explicit chain
vertices.
\cert{check_type8_product.m}

These classes give three explicit elliptic covers over the quartic
field \(K=\mathbf Q(\rho)\), \(G(\rho)=0\). Write
\(q(T)=G(T)/(1280(T-\rho))\). In the two-cover test the component at
zero is \(b\); since both \(b\) and \(b_0\) are nonzero squares, the
common rational scalar is square. Hence \((b-\rho)/(b_0-\rho)\)
is square over \(K\). As \(G(b)\) and \(G(b_0)\) are also squares,
the corresponding elliptic equation is \(W^2=q(b_0)q(b)\).
The next step uses the full singleton polynomial
and a second elliptic cover.

\subsubsection{A singular fiber and the reciprocal covers}\label{arith:73}

The class represented by \(b_0=1/4\) contains no
valid chain at any rational height. For the three classes of \S\ref{arith:72},
in the order \(b_0=1/4,9/4,25/4\), the quartic-field elliptic curves
there have 2-Selmer dimensions \(5,4,5\). Each class also gives an
elliptic curve with 2-Selmer dimension three. These computations are
unconditional.

\emph{Proof.} Scale the center and singleton squares by four. Write
\(B=4b\), \(B_0=4b_0\) and
\[P_B(X)=(X-B)((X-B-4)^2-16B)((X-B-16)^2-64B),\qquad
 F(B,X)=\frac{P_B(X)-P_B(1)}{X-1}.\]
All four roots of this monic quartic must be rational squares. Its
constant term is \(5f(B)\), where
\(f(T)=T^4-18T^3+114T^2-97T+405\). For a root \(r\) of \(f\),
\S\ref{arith:72} gives \((B-r)/(B_0-r)\in K^{\times2}\) in
\(K=\mathbf Q(r)\).

\textbf{The singular fiber.} \S\ref{arith:72} gives \(B\in\mathbf Z_{23}\),
so all four roots are integral there. Enumerating the twelve square
residues modulo 23, including zero, gives precisely \(B=1,2,9\)
for which all four roots are squares, counted with multiplicity.
Since \(-4\) is a simple root of \(f\bmod23\), the degree-one
completion there tests the squareclass by \((B+4)(B_0+4)\).
For \(B_0=1\) only \(B\equiv1\) survives. The special fiber factors as
\[F(1,X)=(X-1)(X-9)^2(X-25).\]
In particular some root satisfies \(X\equiv9\pmod{23}\).
Put \(U=B-1, V=X-9\). The terms of degree at most two of
\(F(1+U,9+V)\) are \(-64(2V^2-15U^2)\).
Since 30 is nonsquare modulo 23, this quadratic form has no
nontrivial zero over the residue field. If \((U,V)\ne(0,0)\),
put \(m=\min(v_{23}(U),v_{23}(V))\ge1\). Its quadratic term has
valuation exactly \(2m\), whereas all higher terms have valuation
at least \(3m\). Thus \(B=1\) and \(X=9\), and the repeated
singleton violates distinctness.
\cert{type8_descent/node.m}

\textbf{Two elliptic curves per class.} Write \(q(T)=f(T)/(T-r)\) and
\[g(T)=T^3q(1/T)/q(0).\]
Both cubics are monic and irreducible over \(K\). The direct
cover has equation \(W^2=q(B_0)q(B)\). Since \(B\) and \(B_0\)
are nonzero squares, the factor \((B B_0)^{-3}/q(0)^2\) is a
square as well. Thus the reciprocal cover is
\(W^2=g(1/B_0)g(1/B)\). For either monic cubic
\(h(T)=T^3+c_2T^2+c_1T+c_0\), writing \(k=h(t_0)\) gives
the Weierstrass model \(y^2=x^3+kc_2x^2+k^2c_1x+k^3c_0\),
with \(x=kt\). Irreducibility gives trivial rational two-torsion.

The cubic root field over \(K\) is the same for all six curves:
scaling a root by \(k\) or taking its reciprocal preserves the
field. It is the degree-twelve field \(M=\mathbf Q(r,r')\)
of two distinct ordered roots of \(f\). The certificate verifies
an absolute model, signature \((0,6)\), and discriminant
\(5528109^5\). Apply the explicit formula of \S\ref{sec:marked} with support
\(2^{22}\), generating bound \(2^{17}\) and no real places.
The rational interval bound gives
\(A_{\rm arch}+S-\log|D_M|>862/1000\).
Consequently the primes of norm at most \(2^{17}\) generate the
class group. PARI verifies their reduction to a factor base of
norms at most 1000. With this certified bound for \(M\) and the
Minkowski bound for other fields, the class group is cyclic of
order three and two-descent gives the dimensions asserted above.

\cert{type13_prym_class/prime_counts.cpp},
\cert{type13_prym_class/check_bound.py},
\cert{type8_descent/factor_base.gp},
\cert{type8_descent/covers.m}

\subsubsection{Global sign constraints and transverse logarithms}\label{arith:74}

No valid distinct-entry type-8 chain exists at any
rational height. In the classes represented by \(b_0=9/4,25/4\), the
reciprocal elliptic curves of \S\ref{arith:73} both have rank three.

\emph{Proof.} Continue with \(B=4b\), \(t=1/B\), \(t_0=1/B_0\),
\(k=g(t_0)\) and the monic reciprocal cubic \(g\) of \S\ref{arith:73}.
Write \(E\) for the model with abscissa \(x=kt\), and put
\(P_0=(kt_0,k^2)\). Use
\[K=\mathbf Q(h),\quad h^4-h^3+28h^2-54h+231=0,
 \qquad r=(-2h^3+h^2-33h+129)/15 .\]
For each branch the certificate supplies three points \(P_i\) of
\(E(K)\) by giving their parameters \(t_i\) and extracting a square
root of \(k^3g(t_i)\) in \(K\). One is \(P_0\); for \(B_0=9\) another
has rational parameter \(-59/64\). The logarithm calculation below
proves independence, so the upper bound three is attained in both cases.
\cert{type8_descent/logarithms.m}

\textbf{Signs for \(B_0=9\).} At each of \(p=17,131,197,269\), enumerate
square residues \(B\) for which \(F(B,X)\) splits into four squares.
For every monic irreducible factor \(d\) of \(f\bmod p\), require
\(d(B)d(9)\) to be square, allowing zero. This follows by taking
the residue-field norm of \((B-r)/(9-r)\) at the corresponding
unramified place when both entries are units. \S\ref{arith:72} excludes
poles at these primes. The sole surviving residue is \(B=9\) in each
case. All four primes give good reduction of \(E\). At each place
above them a putative global point therefore reduces to \(P_0\) or
\(-P_0\).

At 131 the residue degrees are \(1,1,2\). To determine the relative
signs, take the subgroup generated by the three independent points.
For \(\ell=2,3\) separately the certificate saturates this subgroup at
\(\ell\): for each line of coefficient vectors modulo \(\ell\), it
tests exact division of the indicated sum in \(E(K)\). When divisible,
replace the generator whose coefficient is one by the divided point
and repeat. This enlarges the subgroup by index \(\ell\). The final
tests have no further division and \(E(K)[\ell]=0\). Since the rank
is three, the resulting points represent all global points modulo
\(mE(K)\), with \(m=16\) or \(9\), respectively.

At each finite place with residue group of order \(N\), multiply by
\(N/\gcd(N,m)\) to annihilate \(m\) times the residue group. Enumerate
coefficient triples modulo \(m\) and require agreement with a signed
reduction of \(P_0\) at every place. For \(m=16\), the primes
\(17,197,131\) leave four triples, and the signs at the two linear
factors of 131 agree. For \(m=9\), the primes \(269,131\) leave
36 triples and tie the quadratic factor to one linear factor.
The only simultaneous patterns are \((+,+,+)\) and \((-,-,-)\).
Negating the global point if necessary makes all signs positive.
\cert{type8_descent/signs.m}

\textbf{Transversality.} Use \(p=131\) for \(B_0=9\) and \(p=2\) for
\(B_0=25\). In the second case there is just one place. \S\ref{arith:72}
makes \(B\) integral at 2. An even \(B\) would give
\(F(B,X)\equiv X^4+X^3+X^2+X+1\pmod2\), which has no root.
Thus \(B\) is an odd square and \(t\equiv t_0\pmod8\).

Let \(V=\prod_{v\mid p}K_v\), of dimension four over \(\mathbf Q_p\),
and take local elliptic logarithms for \(dx/(2y)\). The certificate
expresses these in integral power bases of the unramified completions.
For each \(P_i\) it takes a common integer multiple across the places,
computes the formal parameter \(z=-x/y\), and divides the concatenated
row by the same power of \(p\) at all places to make it primitive.
Every formal parameter used has valuation at least eight. For the
integral Weierstrass model the coefficient of \(z^j\) in the formal
logarithm lies in \(j^{-1}\mathcal O_{K_v}\), by integration of the
invariant differential. Hence replacing the logarithm by
\(z\) preserves each normalized coordinate modulo \(p^5\).

Let \(D\) be the resulting three-row matrix of true normalized logarithms
and let \(a\) be the coordinate vector of \(1/k\) in \(V\). For 131 the
minimum valuation of its maximal minors is zero, and for 2 it is one.
Adjoining \(a\) gives a determinant of exactly the same valuation in
both cases. These statements are determined by the calculation modulo
\(p^5\). In particular the rows are independent. Their signed maximal
minors, divided by the common power of \(p\), define an integral linear
functional \(\lambda\) annihilating all three rows, with
\(\lambda(1/k)\) a unit. As the three points span \(E(K)\otimes\mathbf Q\),
the functional kills the logarithm of every global point.

\textbf{The full parameter disks.} Write \(t=t_0+d\xi\), with
\(d=131\) or \(8\) and \(\xi\in\mathbf Z_p\). At 131 the sign result
puts the lift on the branch through \(P_0\) at every place. At 2 choose
the sign giving \(y/k^2\in1+4\mathcal O_{K_v}\).
On these branches \(g(t_0+d\xi)/k\) has constant term one and its
square root takes value one at zero. For \(t\ne t_0\) the chord used
in subtracting \(P_0\) has slope \((y+k^2)/(k(t-t_0))\). Its valuation
is \(-v_p(t-t_0)\) at 131 and \(1-v_2(t-t_0)\) at two. The difference
therefore belongs to the formal group throughout the disk. Pulling
back the differential, and dividing out the scalar \(d/2\), its
logarithm has the same zeros under \(\lambda\) as
\[\lambda\!\left(\frac1k\int_0^\xi
     \left(\frac{g(t_0+dX)}{k}\right)^{-1/2}dX\right).\]
The linear coefficient is a unit. At 131 the coefficient of \(X^i\)
in the integrand has valuation at least \(i\), since the binomial
coefficients of exponent \(-1/2\) are 131-integral. At two, for
\(i\ge1\), the valuation is at least \(i+1\): a term with \(j\)
nonconstant factors gains \(3i\) from \(d^i\), and
\(v_2\binom{-1/2}{j}=s_2(j)-2j\ge1-2i\), where \(s_2\) counts
binary ones. After division by \(i+1\) the bounds are still positive
and tend to infinity. Thus every nonlinear coefficient belongs to
\(p\mathbf Z_p\); for \(\xi\ne0\) the linear term has strictly smallest
valuation. Global points require \(\xi=0\).

The only parameters compatible with valid splitting in the two classes
are therefore \(B=9\) and \(25\), respectively. At both values
\(P_B(1)=0\), so the A and B vertex squares coincide. The class
\(B_0=1\) was excluded in \S\ref{arith:73}.


\section{Reproducing the arithmetic}\label{app:verification}
The companion repository contains the \dossier{}, supporting proofs in \texttt{proofs/}, and executable certificates in \texttt{scripts/}. Each linked filename in this paper is relative to \texttt{scripts/} and opens the version used in the proof. From its repository root, Python files run with Python~3, algebraic identity checkers require SymPy, \texttt{.gp} files run with PARI/GP, and \texttt{.m} files in the supplement are Magma programs. The class-group and descent calculations for the marked vertex use PARI/GP 2.19.1 and Magma V2.30-1. A C++20 compiler linked against PARI is required for the parallel prime counts. All computations refer to exact polynomial or number-field models; local computations are carried out to the precision stated in the corresponding proofs.

\paragraph{Classification.} Run \cert{check_all_shift_chain_types.py} to enumerate the integer partitions in Table~\ref{tab:parts} and check the elimination identities in the normalized models. Its genus computations establish necessary conditions. The certificates for excluding the resulting rational parameters are linked in the individual cases.

\paragraph{Type 18 and the marked vertex.}\label{app:type18}
For the partition \((5,\varnothing,(1,1,1,1,1))\) the centered side is \(O_{10}\) and the other side consists of five pairs \(\pm v_i\). Doubling the entries gives \(\{\pm1,\pm3,\pm5,\pm7,\pm9\}\) on the centered side. With \(y_i=2v_i\), the four equations \eqref{eq:chains} are exactly \eqref{eq:four-even-moments}. The marked-vertex calculation in Sections~\ref{sec:vertices} and~\ref{sec:marked} therefore forces the two sides to coincide. This proves the Type~18 exclusion for every rational choice of centers.

The five-class calculation, prime-ideal bound, factor-base verification, elliptic descent, pairing and logarithm must be read in that order. From the root of a fresh checkout the following commands reproduce these steps, assuming the indicated programs and development library are available on the usual paths.
\begin{samepage}
\begin{Verbatim}[fontsize=\footnotesize,breaklines=true]
mkdir -p build/type18-class
magma scripts/check_type18_product_classes.m
c++ -O3 -std=c++20 scripts/type13_prym_class/prime_counts.cpp \
    -lpari -o build/type18-class/prime_counts
python3 scripts/type13_prym_class/check_bound.py \
    build/type18-class/prime_counts --type18 --jobs 12
gp -fq scripts/type18_descent/check_factor_base.gp
magma scripts/type18_descent/rank.m
magma scripts/type18_descent/pairing.m
gp -fq scripts/type18_descent/logarithm.gp
\end{Verbatim}
\end{samepage}
The explicit-formula margin exceeds \(0.284\). The proof bound supplied to the elliptic descent is 1000 for the identified degree-twelve field; applying it presupposes both the explicit-formula calculation and the prime-ideal reduction up to \(85{,}000{,}000\). The pair of Selmer representatives in the pairing check is given by fixed quartics. The last command returns the primitive residue \(\bar h^3+\bar h\). Parallel interval verification of the factor base is described in \cert{type18_descent/README.md}.

\paragraph{Type 10.} Appendix~\ref{app:type10} gives, in order, the elliptic quotient, the two height calculations, their local targets, the root counts, and the elliptic-index test. Given those upstream inputs, the final two commands are
\begin{samepage}
\begin{Verbatim}[fontsize=\footnotesize,breaklines=true]
python3 scripts/check_sixth_mode_qc_roots13.py
python3 scripts/check_sixth_mode_candidate_sieve.py
\end{Verbatim}
\end{samepage}
The first isolates roots of the necessary height equations. The second relates each surviving local candidate to a residue of the global elliptic multiplier. Root counting on its own does not certify that every isolated local root is rational.

\paragraph{Other types.} The links beside the equations in Appendices~\ref{app:cases} and~\ref{app:details} identify the local inputs, ranks and global descent calculations. Where a Magma command returns a finite-index subgroup, the proof either verifies the index condition at the relevant primes or uses linear dependence of logarithms over the ground field. A congruence defining a residue disk leaves infinitely many local parameters. The final implications in the proofs use an exact rank-zero computation, a bounded height proved from the equations, a certified rational-image computation, or an obstruction on a whole disk.

\phantomsection\addcontentsline{toc}{section}{References}
\begin{thebibliography}{99}
\bibitem{cmsv}
D. Coppersmith, M. J. Mossinghoff, D. Scheinerman and J. M. VanderKam,
\emph{Ideal solutions in the Prouhet--Tarry--Escott problem}, Math. Comp. \textbf{93} (2024), 2473--2501.
\href{https://doi.org/10.1090/mcom/3917}{Journal article}.
\href{https://arxiv.org/html/2304.11254v1}{Preprint, version 1}.
\bibitem{banks}
W. D. Banks and A. J. van der Poorten, \emph{Squares from products of integers}, Gaz. Aust. Math. Soc. \textbf{31} (2004), 40--42.
\href{https://austms.org.au/wp-content/uploads/Gazette/2004/Mar04/Mar2004.pdf}{Journal issue}.
\bibitem{wright}
E. M. Wright, \emph{Equal sums of like powers}, Canad. Math. Bull. \textbf{8} (1965), 193--202.
\href{https://doi.org/10.4153/CMB-1965-015-7}{Journal article}.
\bibitem{choudhry}
A. Choudhry, \emph{A new approach to the Tarry--Escott problem}.
\href{https://arxiv.org/abs/1603.00206v2}{arXiv:1603.00206v2} (2017).
\bibitem{BruinStoll}
N. Bruin and M. Stoll, \emph{Two-cover descent on hyperelliptic curves}, Math. Comp. \textbf{78} (2009), 2347--2370.
\href{https://arxiv.org/abs/0803.2052}{Author preprint}.
\bibitem{magma}
W. Bosma, J. Cannon and C. Playoust, \emph{The Magma algebra system. I. The user language}, J. Symbolic Comput. \textbf{24} (1997), 235--265.
\href{https://doi.org/10.1006/jsco.1996.0125}{Journal article}.
\href{https://magma.maths.usyd.edu.au/magma/handbook/}{Magma Handbook}.
\bibitem{miller}
J. C. Miller, \emph{Class numbers of totally real fields and applications to the Weber class number problem}, Acta Arith. \textbf{164} (2014), 381--397.
\href{https://doi.org/10.4064/aa164-4-4}{Journal article}.
\bibitem{pari}
The PARI Group, \emph{PARI/GP}, University of Bordeaux.
\href{https://pari.math.u-bordeaux.fr/}{Software and source}.
\href{https://pari.math.u-bordeaux.fr/dochtml/html-stable/General_number_fields.html}{Number-field documentation}.
\bibitem{blp}
P. Borwein, P. Lisoněk and C. Percival, \emph{Computational investigations of the Prouhet--Tarry--Escott problem}, Math. Comp. \textbf{72} (2003), 2063--2070.
\href{https://www.daemonology.net/papers/pte.pdf}{Author copy}.
\bibitem{caley}
T. Caley, \emph{The Prouhet--Tarry--Escott problem}, Ph.D. thesis, University of Waterloo, 2012.
\href{https://uwspace.uwaterloo.ca/items/b24da5f2-8532-426b-ad80-3009df50bdba}{University repository}.

\bibitem{BruinPrym}
N. Bruin, \emph{The arithmetic of Prym varieties in genus 3}, Compos. Math. \textbf{144} (2008), 317--338.
\href{https://doi.org/10.1112/S0010437X07003314}{Journal article}.
\bibitem{Cantor}
D. G. Cantor, \emph{Computing in the Jacobian of a hyperelliptic curve}, Math. Comp. \textbf{48} (1987), 95--101.
\href{https://doi.org/10.1090/S0025-5718-1987-0866101-0}{Journal article}.
\bibitem{Silverman}
J. H. Silverman, \emph{The Arithmetic of Elliptic Curves}, 2nd ed., Graduate Texts in Mathematics \textbf{106}, Springer, 2009.
\href{https://doi.org/10.1007/978-0-387-09494-6}{Publisher record}.
\bibitem{BalakrishnanDogra2018}
J. S. Balakrishnan and N. Dogra,
\emph{Quadratic Chabauty and rational points I: $p$-adic heights},
Duke Math. J. \textbf{167} (2018), no.~11, 1981--2038.
\href{https://doi.org/10.1215/00127094-2018-0013}{doi:10.1215/00127094-2018-0013}.
The cited text is also available as
\href{https://arxiv.org/abs/1601.00388v2}{arXiv:1601.00388v2}.

\bibitem{BDMTV}
J. S. Balakrishnan, N. Dogra, J. S. M\"uller, J. Tuitman and J. Vonk,
\emph{Quadratic Chabauty for modular curves: algorithms and examples},
Compos. Math. \textbf{159} (2023), no.~6, 1111--1152.
\href{https://doi.org/10.1112/S0010437X23007170}{doi:10.1112/S0010437X23007170}.

\bibitem{GajovicMuller}
S. Gajovi\'c and J. S. M\"uller,
\emph{Computing $p$-adic heights on hyperelliptic curves},
\href{https://arxiv.org/abs/2307.15787v4}{arXiv:2307.15787v4} (2024).

\bibitem{DDMM}
T. Dokchitser, V. Dokchitser, C. Maistret and A. Morgan,
\emph{Arithmetic of hyperelliptic curves over local fields},
Math. Ann. \textbf{385} (2023), 1213--1322.
\href{https://doi.org/10.1007/s00208-021-02319-y}{doi:10.1007/s00208-021-02319-y}.

\bibitem{BPR}
M. Baker, S. Payne and J. Rabinoff,
\emph{On the structure of nonarchimedean analytic curves},
in \emph{Tropical and Non-Archimedean Geometry},
Contemp. Math. \textbf{605}, Amer. Math. Soc., Providence, RI, 2013,
93--121. \href{https://doi.org/10.1090/conm/605/12113}{doi:10.1090/conm/605/12113}.
The cited expanded text is
\href{https://arxiv.org/abs/1404.0279v1}{arXiv:1404.0279v1}.

\bibitem{Schaefer1995}
E. F. Schaefer,
\emph{2-Descent on the Jacobians of hyperelliptic curves},
J. Number Theory \textbf{51} (1995), no.~2, 219--232.
\href{https://doi.org/10.1006/jnth.1995.1044}{doi:10.1006/jnth.1995.1044}.

\bibitem{PoonenSchaefer}
B. Poonen and E. F. Schaefer,
\emph{Explicit descent for Jacobians of cyclic covers of the projective line},
J. Reine Angew. Math. \textbf{488} (1997), 141--188.
\href{https://math.mit.edu/~poonen/papers/descent.pdf}{Author manuscript}.

\end{thebibliography}

\end{document}
